\documentclass[letterpaper, journal]{IEEEtran}
%DIF LATEXDIFF DIFFERENCE FILE
%DIF DEL tmp2.tex   Tue Aug 11 22:59:53 2026
%DIF ADD tmp1.tex   Tue Aug 11 22:59:53 2026

\def\cmt{false} % true for comments, false for no comments
\def\pub{true} % true for publication, false for draft
\newcommand*{\template}{template} 
\input{\template/preamble/preamble_journal.tex}
%DIF 7a7
\allowdisplaybreaks %DIF > 
%DIF -------

%DIF < \newcommand*{\figSizeTwoCol}{.49}
%DIF -------
% -------------------------------- %DIF > 
%DIF < \newcommand*{\figSizeOneCol}{0.98}
%      Manuscript version %DIF > 
\newcommand{\version}{0.9.3} % version of the manuscript %DIF > 
% -------------------------------- %DIF > 
 %DIF > 
\newcommand{\resultFig}{./src/figures} % result figure path %DIF > 
 %DIF > 
%% RMSE values for real-time experiments %DIF > 
% RMSE:  %DIF > 
% ------------------------------- %DIF > 
% C NAC : %DIF > 
% error q1 (1) 19.738, (2) 18.520 [imp: 0.062] %DIF > 
% error q2 (1) 10.561, (2) 12.435 [imp: -0.177] %DIF > 
% filter error   (1) 113.671, (2) 102.581 [imp: 0.098] %DIF > 
 %DIF > 
\newcommand{\NACEpiOneAngOne}{19.738} %DIF > 
\newcommand{\NACEpiOneAngTwo}{10.561} %DIF > 
\newcommand{\NACEpiTwoAngOne}{18.520} %DIF > 
\newcommand{\NACEpiTwoAngTwo}{12.435} %DIF > 
\newcommand{\NACImpAngOne}{-6.174} %DIF > 
\newcommand{\NACImpAngTwo}{17.745} %DIF > 
% ------------------------------- %DIF > 
% C AUX : %DIF > 
% error q1 (1) 20.991, (2) 17.314 [imp: 0.175] %DIF > 
% error q2 (1) 8.268, (2) 2.421 [imp: 0.707] %DIF > 
% filter error   (1) 116.541, (2) 95.281 [imp: 0.182] %DIF > 
 %DIF > 
\newcommand{\AUXEpiOneAngOne}{20.991} %DIF > 
\newcommand{\AUXEpiOneAngTwo}{8.268} %DIF > 
\newcommand{\AUXEpiTwoAngOne}{17.314} %DIF > 
\newcommand{\AUXEpiTwoAngTwo}{2.421} %DIF > 
\newcommand{\AUXImpAngOne}{-17.515} %DIF > 
\newcommand{\AUXImpAngTwo}{-70.715} %DIF > 
% ------------------------------- %DIF > 
% C CONAClow : %DIF > 
% error q1 (1) 15.984, (2) 11.190 [imp: 0.300] %DIF > 
% error q2 (1) 8.883, (2) 2.377 [imp: 0.732] %DIF > 
% filter error   (1) 90.236, (2) 63.587 [imp: 0.295] %DIF > 
 %DIF > 
\newcommand{\CONAClowEpiOneAngOne}{15.984} %DIF > 
\newcommand{\CONAClowEpiOneAngTwo}{8.883} %DIF > 
\newcommand{\CONAClowEpiTwoAngOne}{11.190} %DIF > 
\newcommand{\CONAClowEpiTwoAngTwo}{2.377} %DIF > 
\newcommand{\CONAClowImpAngOne}{-29.991} %DIF > 
\newcommand{\CONAClowImpAngTwo}{-73.244} %DIF > 
% ------------------------------- %DIF > 
% C CONAChigh : %DIF > 
% error q1 (1) 16.836, (2) 11.492 [imp: 0.317] %DIF > 
% error q2 (1) 9.086, (2) 3.144 [imp: 0.654] %DIF > 
% filter error   (1) 94.327, (2) 65.886 [imp: 0.302] %DIF > 
 %DIF > 
\newcommand{\CONAChighEpiOneAngOne}{16.836} %DIF > 
\newcommand{\CONAChighEpiOneAngTwo}{9.086} %DIF > 
\newcommand{\CONAChighEpiTwoAngOne}{11.492} %DIF > 
\newcommand{\CONAChighEpiTwoAngTwo}{3.144} %DIF > 
\newcommand{\CONAChighImpAngOne}{-31.741} %DIF > 
\newcommand{\CONAChighImpAngTwo}{-65.402} %DIF > 
% ------------------------------- %DIF > 
\newcommand{\zoomStartTime}{26.6} %DIF > 
\newcommand{\zoomEndTime}{27.3} %DIF > 
 %DIF > 
% ================================ %DIF > 
% Document starts from here %DIF > 
% =============================== %DIF > 
%DIF PREAMBLE EXTENSION ADDED BY LATEXDIFF
%DIF UNDERLINE PREAMBLE %DIF PREAMBLE
\RequirePackage[normalem]{ulem} %DIF PREAMBLE
\RequirePackage{color}\definecolor{RED}{rgb}{1,0,0}\definecolor{BLUE}{rgb}{0,0,1} %DIF PREAMBLE
\providecommand{\DIFadd}[1]{{\protect\color{blue}\uwave{#1}}} %DIF PREAMBLE
\providecommand{\DIFdel}[1]{{\protect\color{red}\sout{#1}}} %DIF PREAMBLE
%DIF SAFE PREAMBLE %DIF PREAMBLE
\providecommand{\DIFaddbegin}{} %DIF PREAMBLE
\providecommand{\DIFaddend}{} %DIF PREAMBLE
\providecommand{\DIFdelbegin}{} %DIF PREAMBLE
\providecommand{\DIFdelend}{} %DIF PREAMBLE
\providecommand{\DIFmodbegin}{} %DIF PREAMBLE
\providecommand{\DIFmodend}{} %DIF PREAMBLE
%DIF FLOATSAFE PREAMBLE %DIF PREAMBLE
\providecommand{\DIFaddFL}[1]{\DIFadd{#1}} %DIF PREAMBLE
\providecommand{\DIFdelFL}[1]{\DIFdel{#1}} %DIF PREAMBLE
\providecommand{\DIFaddbeginFL}{} %DIF PREAMBLE
\providecommand{\DIFaddendFL}{} %DIF PREAMBLE
\providecommand{\DIFdelbeginFL}{} %DIF PREAMBLE
\providecommand{\DIFdelendFL}{} %DIF PREAMBLE
%DIF COLORLISTINGS PREAMBLE %DIF PREAMBLE
\RequirePackage{listings} %DIF PREAMBLE
\RequirePackage{color} %DIF PREAMBLE
\lstdefinelanguage{DIFcode}{ %DIF PREAMBLE
%DIF DIFCODE_UNDERLINE %DIF PREAMBLE
  moredelim=[il][\color{red}\sout]{\%DIF\ <\ }, %DIF PREAMBLE
  moredelim=[il][\color{blue}\uwave]{\%DIF\ >\ } %DIF PREAMBLE
} %DIF PREAMBLE
\lstdefinestyle{DIFverbatimstyle}{ %DIF PREAMBLE
	language=DIFcode, %DIF PREAMBLE
	basicstyle=\ttfamily, %DIF PREAMBLE
	columns=fullflexible, %DIF PREAMBLE
	keepspaces=true %DIF PREAMBLE
} %DIF PREAMBLE
\lstnewenvironment{DIFverbatim}[1][]{\lstset{style=DIFverbatimstyle}}{} %DIF PREAMBLE
\lstnewenvironment{DIFverbatim*}[1][]{\lstset{style=DIFverbatimstyle,showspaces=true}}{} %DIF PREAMBLE
\lstset{extendedchars=\true,inputencoding=utf8}

%DIF END PREAMBLE EXTENSION ADDED BY LATEXDIFF

\begin{document}
\bstctlcite{IEEEexample:BSTcontrol}

\title{
  Constrained Optimization-Based Neuro-Adaptive Control (CONAC) for \DIFdelbegin \DIFdel{Euler-Lagrange }\DIFdelend \DIFaddbegin \DIFadd{Unknown }\DIFaddend Systems Under \DIFdelbegin \DIFdel{Weight and }\DIFdelend \DIFaddbegin \DIFadd{Multiple Convex }\DIFaddend Input Constraints
} %% Article title

\author{\DIFdelbegin \DIFdel{Myeongseok Ryu, Donghwa Hong, and Kyunghwan Choi
}%DIFDELCMD < \thanks{
%DIFDELCMD <     Myeongseok Ryu and Kyunghwan Choi are with the Cho Chun Shik Graduate School of Mobility, Korea Advanced Institute of Science and Technology (KAIST), Daejeon 34051, Korea (e-mail: dding\_98@kaist.ac.kr and kh.choi@kaist.ac.kr).
%DIFDELCMD < 

%DIFDELCMD <     Donghwa Hong is with the Department of Mechanical and Robotics Engineering, Gwangju Institute of Science and Technology (GIST), Gwangju 61005, Korea (e-mail: fairytalef00@gm.gist.ac.kr).
%DIFDELCMD < 

%DIFDELCMD <     This work was supported in part by a Korea Research Institute for Defense Technology planning and advancement (KRIT) grant funded by Korea government DAPA (Defense Acquisition Program Administration) (No. KRIT-CT-22-087, Synthetic Battlefield Environment based Integrated Combat Training Platform) and in part by the National Research Foundation of Korea (NRF) grant funded by the Korea government (MSIT) (RS-2025- 00554087). (Corresponding author: Kyunghwan Choi)
%DIFDELCMD <     }%%%
\DIFdelend }
%DIF >  \author{
%DIF >    Myeongseok Ryu, Donghwa Hong, and Kyunghwan Choi
%DIF >  \thanks{
%DIF >    Myeongseok Ryu, Donghwa Hong and Kyunghwan Choi are with the Cho Chun Shik Graduate School of Mobility, Korea Advanced Institute of Science and Technology (KAIST), Daejeon 34051, Korea (e-mail: \{dding\_98, fairytale, kh.choi\}@kaist.ac.kr).

\DIFdelbegin %DIFDELCMD < \markboth{Journal of \LaTeX\ Class Files,~Vol.~18, No.~9, September~2020}%%%
\DIFdelend %DIF >    This work was supported in part by a Korea Research Institute for Defense Technology planning and advancement (KRIT) grant funded by Korea government DAPA (Defense Acquisition Program Administration) (No. KRIT-CT-22-087, Synthetic Battlefield Environment based Integrated Combat Training Platform) and in part by the National Research Foundation of Korea (NRF) grant funded by the Korea government (MSIT) (RS-2025- 00554087). (Corresponding author: Kyunghwan Choi)
%DIF >    }}
\DIFaddbegin 

%DIF >  \markboth{Journal of \LaTeX\ Class Files,~Vol.~18, No.~9, September~2020}%
%DIF >  {How to Use the IEEEtran \LaTeX \ Templates}
\markboth{Submitted to IEEE Transactions on Systems, Man, and Cybernetics: Systems}\DIFaddend %
{\DIFdelbegin \DIFdel{How to Use the IEEEtran }%DIFDELCMD < \LaTeX %%%
\DIFdel{\ Templates}\DIFdelend }
%DIF >  \markboth{\cred{\textit{(Version \version)}} Submitted to IEEE Transactions on Systems, Man, and Cybernetics: Systems}%
%DIF >  {\cred{\textit{(Version \version)}} Ryu \MakeLowercase{\textit{et al.}}: CONAC for Unknown Systems Under Multiple Convex Input Constraints}

\maketitle

%% Abstract
\begin{abstract}
  This study presents a constrained optimization-based neuro-adaptive control (CONAC) for \DIFdelbegin \DIFdel{unknown Euler-Lagrange }\DIFdelend \DIFaddbegin \DIFadd{a class of unknown multi-input-multi-output (MIMO) }\DIFaddend systems subject to \DIFdelbegin \DIFdel{weight and }\DIFdelend \DIFaddbegin \DIFadd{multiple }\DIFaddend convex input constraints. 
  A deep neural network (DNN) is employed to approximate the ideal \DIFdelbegin \DIFdel{stabilizing }\DIFdelend control law while \DIFdelbegin \DIFdel{simultaneously learning the unknown system dynamics and addressing both types of }\DIFdelend \DIFaddbegin \DIFadd{handling multiple convex input }\DIFaddend constraints within a unified constrained optimization framework.
  The \DIFdelbegin \DIFdel{adaptation law of DNN weights is formulated to solve the defined }\DIFdelend \DIFaddbegin \DIFadd{adaptive variables, including the DNN weights and Lagrange multipliers are updated through adaptation laws derived from the formulated }\DIFaddend constrained optimization problem, \DIFdelbegin \DIFdel{ensuring satisfaction of first-order }\DIFdelend \DIFaddbegin \DIFadd{yielding Karush-Kuhn-Tucker (KKT) }\DIFaddend optimality conditions at \DIFdelbegin \DIFdel{steady state}\DIFdelend \DIFaddbegin \DIFadd{equilibrium}\DIFaddend .
  The controller's stability is rigorously analyzed using Lyapunov theory, \DIFdelbegin \DIFdel{guaranteeing bounded }\DIFdelend \DIFaddbegin \DIFadd{establishing the uniform ultimate boundedness (UUB) of the }\DIFaddend tracking errors and \DIFdelbegin \DIFdel{bounded DNN weights}\DIFdelend \DIFaddbegin \DIFadd{adaptive variables}\DIFaddend .
  The proposed controller is \DIFdelbegin \DIFdel{validated through a }\DIFdelend \DIFaddbegin \DIFadd{compared with existing methods through }\DIFaddend real-time \DIFdelbegin \DIFdel{implementation on a 2-DOF }\DIFdelend \DIFaddbegin \DIFadd{experiments on a 2-degree-of-freedom (DOF) }\DIFaddend robotic manipulator, \DIFdelbegin \DIFdel{demonstrating its effectiveness in achieving accurate trajectory tracking while satisfying all imposed constraints}\DIFdelend \DIFaddbegin \DIFadd{highlighting its superior capability to handle input saturation and its feasibility in real-time implementation}\DIFaddend .
\end{abstract}

\begin{IEEEkeywords}
  Neuro-adaptive control, constrained optimization, deep neural network,  \DIFdelbegin \DIFdel{input constraint, weight constraint}\DIFdelend \DIFaddbegin \DIFadd{multiple convex input constraints}\DIFaddend .
\end{IEEEkeywords}

\DIFdelbegin \section*{\DIFdel{Notation}}
%DIFAUXCMD
%DIFDELCMD < 

%DIFDELCMD < %%%
%DIF <  In this study, the following notation is used:
%DIFDELCMD < 

%DIFDELCMD < %%%
%DIF <  \begin{itemize}
%DIF <      \item := denotes ``is defined as" or ``is equal to".
%DIF <      \item $\R$, $\R^n$ and $\R^{n\times m}$ denote the set of real numbers, $n\times 1$ dimensional vectors and $n\times m$ dimensional matrices, respectively, and $\R_{\ge0}:=[0,\infty)$ and $\R_{>0}:=(0,\infty)$.
%DIF <      \item $(\cdot)^\top$ denotes the transpose of a vector or matrix.
%DIF <      \item $\norm{\cdot}$ and $\norm{\cdot}_F$ denote the Euclidean and Frobenius norm, respectively.
%DIF <      \item $\otimes$ denotes the Kronecker product \cite[Chap. 7, Def. 7.1.2]{Bernstein:2009aa}.
%DIF <      \item A vector and a matrix are denoted by $\mv{x}=[x_i]_{i\in\{1,\cdots,n\}}\in\R^n$ and $
%DIF <          \mm A
%DIF <          := 
%DIF <          [a_{ij}]
%DIF <          _{
%DIF <              i\in\{1,\cdots,n\},j\in\{1,\cdots ,m\}
%DIF <          }\allowbreak\in\R^{n\times m}
%DIF <          $, respectively.
%DIF <      \item $\myrow_i(\mm A)$ denotes the $i$th row of the matrix $\mm A\in\R^{n\times m}$. 
%DIF <      \item For the matrix $\mm{A}\in\R^{n\times m}$, $\myvec(\mm A):=(\myrow_1(\mm A^\top),\cdots,\myrow_m(\mm A^\top) )^\top\in\R^{nm}$ denotes the vectorization of $\mm A$.
%DIF <      % \item $\mycol_i(A)$ and $\myrow_j(A)$ denote the $i$th column and $j$th row of $A\in\R^{n\times m}$, respectively.
%DIF <      % \item $\myvec(A):= [\mycol_1(A)^\top  ,\cdots,\mycol_m(A)^\top  ]^\top   $ for $A\in\R^{n\times m}$.
%DIF <      \item $\lambda_{\min}(\mm A)$ denotes the minimum eigenvalue of the matrix $\mm A\in\R^{n\times n}$.
%DIF <      \item $\mm I_n$ denotes the $n\times n$ identity matrix, and $\mm 0_{n\times m}$ denotes the $n\times m$ zero matrix.
%DIF <  \end{itemize}
%DIFDELCMD < 

%DIFDELCMD < %%%
\DIFdel{The Kronecker product is denoted by $\otimes$ \mbox{%DIFAUXCMD
\cite[Chap. 7, Def. 7.1.2]{Bernstein:2009aa}}\hskip0pt%DIFAUXCMD
.
A vector and a matrix are denoted by $\mv{x}=[x_i]_{i\in\{1,\cdots,n\}}\in\R^n$ and $
\mm A
:= 
[a_{ij}]
_{
    i\in\{1,\cdots,n\},j\in\{1,\cdots ,m\}
}\allowbreak\in\R^{n\times m}
$, respectively.
Given a matrix $\mm A\in\R^{n\times m}$, $\myrow_i(\mm A)$ denotes the $i$th row of $\mm A$. 
For the matrix $\mm{A}\in\R^{n\times m}$, $\myvec(\mm A):=(\myrow_1(\mm A^\top),\cdots,\myrow_m(\mm A^\top) )^\top\in\R^{nm}$ denotes the vectorization of $\mm A$.
$\lambda_{\min}(\mm A)$ denotes the minimum eigenvalue of the matrix $\mm A\in\R^{n\times n}$.
The $n\times n$ identity and $n\times m$ zero matrices are denoted by $\mm I_n$ and $\mm 0_{n\times m}$, respectively.
}%DIFDELCMD < 

%DIFDELCMD < %%%
%DIF <   SECTION INTRODUCTION ===================================
\DIFdelend %DIF >  ------------------------------------
%DIF >        SECTION INTRODUCTION 
%DIF >  ------------------------------------
%DIF >  \vspace{100mm}                
\section{Introduction}

\subsection{Background}

\IEEEPARstart{M}{any} \DIFdelbegin \DIFdel{engineering systems, including those in aerospace, robotics, and automotive applications, can be modeled using Euler-Lagrange systems.
These systems are governed by dynamic equations derived from energy principles and describe the motion of mechanical systems with constraints.
In practice, however, such systems often exhibit uncertainties due to }\DIFdelend \DIFaddbegin \DIFadd{control systems are constrained by various safety requirements, actuator limitations, and physical restrictions \mbox{%DIFAUXCMD
\cite{Ji:2025aa}}\hskip0pt%DIFAUXCMD
.
These constraints often manifest as convex input constraints, which can be in the form of nonlinear constraints rather than the element-wise (symmetric or asymmetric) constraints, }\ie 
\DIFadd{simple box constraints that limit each input independently, }\ie \DIFadd{$\underline{u}_i\le u_i\le\overline{u}_i$, $\forall i\in\{1,\cdots,n\}$, where $\underline{u}_i,\overline{u}_i\in\R$ denote the lower and upper bounds of the $i$th input $u_i$, respectively.
}\MSRY{여기서 element-wise를 먼저 언급하였습니다; 마찬가지로 문헌에서 많이 쓰이는 symmetric/asymmetric도 언급하였습니다.}
\DIFadd{For example, under road friction limits, the admissible global force and moment of vehicles can be characterized by a weighted norm bound, which yields a convex ellipsoidal constraint \mbox{%DIFAUXCMD
\cite{Chen:2024aa}}\hskip0pt%DIFAUXCMD
.
In motor drive systems, the feasible voltage region is often represented by circular \mbox{%DIFAUXCMD
\cite{Ryu:2025aa} }\hskip0pt%DIFAUXCMD
or hexagonal sets \mbox{%DIFAUXCMD
\cite{Monzen:2026aa}}\hskip0pt%DIFAUXCMD
.
Furthermore, admissible input sets may be determined by multiple constraints that must hold simultaneously, such as element-wise constraints combined with norm constraints, which can be represented as the intersection of convex sets.
Therefore, addressing these multiple convex input constraints is crucial for ensuring the safety and reliability of control systems in real-world applications.
}

\DIFadd{In parallel, control systems often involve uncertain or even unknown dynamics due to modeling errors, }\DIFaddend unmodeled dynamics, \DIFdelbegin \DIFdel{parameter variations, or }\DIFdelend \DIFaddbegin \DIFadd{and }\DIFaddend external disturbances.
\DIFdelbegin \DIFdel{These uncertainties can significantly degrade control performance and, in some cases, lead to instability. 
In the worst cases, any prior knowledge of the system dynamics}\DIFdelend \DIFaddbegin \DIFadd{Although adaptive control provides a systematic mechanism for online adaptation of controllers \mbox{%DIFAUXCMD
\cite{Ioannou:2006aa}}\hskip0pt%DIFAUXCMD
, conventional formulations often rely on structural assumptions about the uncertain or unknown dynamics, which }\DIFaddend may not be available \DIFdelbegin \DIFdel{, making it challenging to design effective controllers.
%DIF <  To address these challenges, adaptive control methods have been widely employed to ensure robust performance in the presence of system uncertainties \cite{Ioannou:2006aa, Tao:2003aa}.
To address these challenges, adaptive control methods have been widely employed to ensure robust performance for uncertain or unknown systems \mbox{%DIFAUXCMD
\cite{Ioannou:2006aa, Tao:2003aa}}\hskip0pt%DIFAUXCMD
.
}%DIFDELCMD < 

%DIFDELCMD < %%%
\DIFdel{More recently, neuro-adaptive control }\DIFdelend \DIFaddbegin \DIFadd{in practice.
Neuro-adaptive control }\DIFaddend (NAC) \DIFdelbegin \DIFdel{approaches have been introduced to approximate uncertain or unknown system dynamics or entire control laws }\DIFdelend \DIFaddbegin \DIFadd{alleviates this assumption by }\DIFaddend using neural networks (NNs) \DIFdelbegin \DIFdel{\mbox{%DIFAUXCMD
\cite{Lewis:1998aa,Farrell:2006aa}}\hskip0pt%DIFAUXCMD
. 
NNs are well-known for their universal approximation property, which allows them to approximate any smooth function over a compact set with minimal error.
Various types of NNs have been utilized in NAC, including simpler architectures like single-hidden layer (SHL) NNs \mbox{%DIFAUXCMD
\cite{Lewis:1996aa,Ge:2010aa, Yesildirek:1995aa,Ryu:2024aa,Ryu:2025aa} }\hskip0pt%DIFAUXCMD
and radial basis function (RBF) NNs \mbox{%DIFAUXCMD
\cite{Liu:2013ab,Ge:2002aa}}\hskip0pt%DIFAUXCMD
, as well as more complex models like deep NNs (DNNs) \mbox{%DIFAUXCMD
\cite{Patil:2022aa} }\hskip0pt%DIFAUXCMD
and their variations.
Conventionally, SHL and RBF NNs are often employed to approximate }\DIFdelend \DIFaddbegin \DIFadd{to approximate the dynamics in a structured and parameterized representation \mbox{%DIFAUXCMD
\cite{Farrell:2006aa,Patil:2022aa}}\hskip0pt%DIFAUXCMD
.
Like adaptive control, NAC methods ensure closed-loop stability using Lyapunov theory, and the NN weights are adapted online without prior knowledge of the system.
The NAC methods have shown promising results in various real-world applications, including vehicle control \mbox{%DIFAUXCMD
\cite{You:2023aa} }\hskip0pt%DIFAUXCMD
and flexible robotic manipulators \mbox{%DIFAUXCMD
\cite{Gao:2026aa}}\hskip0pt%DIFAUXCMD
.
}

\DIFadd{Because multiple convex input constraints and }\DIFaddend uncertain or unknown \DIFdelbegin \DIFdel{system dynamics or controllers due to their simplicity \mbox{%DIFAUXCMD
\cite{Esfandiari:2014aa,Esfandiari:2015aa,Yesildirek:1995aa,Gao:2006aa}}\hskip0pt%DIFAUXCMD
.
However, since DNNs offer greater expressive power, making them more effective for complex system approximations \mbox{%DIFAUXCMD
\cite{Rolnick:2018aa}}\hskip0pt%DIFAUXCMD
, variations of DNNs, such as long short-term memory (LSTM) networks for time-varying dynamics \mbox{%DIFAUXCMD
\cite{Griffis:2023aa}}\hskip0pt%DIFAUXCMD
, physics-informed NNs (PINNs) for leveraging physical system knowledge \mbox{%DIFAUXCMD
\cite{Hart:2024aa}}\hskip0pt%DIFAUXCMD
, and convolutional NNs (CNNs) for learning from historical sensor data \mbox{%DIFAUXCMD
\cite{Ryu:2024ac}}\hskip0pt%DIFAUXCMD
, have further extended the capabilities of NAC systems.
}%DIFDELCMD < 

%DIFDELCMD < %%%
\DIFdel{A critical aspect of NAC is the NN weight adaptation law, which governs how NN weights are updated.
Most studies derived these laws using Lyapunov-based methods, ensuring the boundedness of the tracking error and weight estimation error, thus maintaining system stability under uncertainty}\DIFdelend \DIFaddbegin \DIFadd{dynamics often coexist in real-world applications, incorporating multiple input constraint handling into NAC is essential.
This motivates a review of existing input constraint handling methods in NAC and related constrained control frameworks}\DIFaddend .

\DIFdelbegin \DIFdel{However, two significant challenges persist in using NNs for adaptive control. 
First, the boundedness of NN weights is not inherently guaranteed}\DIFdelend \DIFaddbegin \subsection{\DIFadd{Literature Review}}

%DIF >  ================================
%DIF >  STRUCTURAL METHODS
%DIF >  ================================

\DIFadd{In Table~\ref{table:intro:comparison}, we summarize input constraint handling methods reported in the literature}\DIFaddend , which can \DIFdelbegin \DIFdel{result in unbounded outputs.
When NN outputs are used directly in the control law, this may lead to excessive control inputs, violating input constraints. 
Such constraints are commonly encountered in industrial systems, where actuators are limited by physical and safety requirements in terms of amplitude, rate, or energy \mbox{%DIFAUXCMD
\cite{Esfandiari:2021aa}}\hskip0pt%DIFAUXCMD
.
Failing to address these constraints can degrade control performance or even destabilize the system.
}\DIFdelend \DIFaddbegin \DIFadd{be broadly categorized into structural methods and optimization-based methods.
In subsequent subsections, these methods are reviewed, highlighting their advantages and limitations.
}\DIFaddend 

\DIFdelbegin \DIFdel{Therefore, addressing these two key issues—ensuring weight boundedness and satisfying input constraints—is essential for the reliable designof NAC.
The following section provides a detailed review of existing solutions to these challenges.
}\DIFdelend \DIFaddbegin \subsubsection{\DIFadd{Structural Methods for Input Constraint Handling in NAC}}
\DIFaddend 

\DIFdelbegin \subsection{\DIFdel{Literature Review}}
%DIFAUXCMD
\addtocounter{subsection}{-1}%DIFAUXCMD
%DIFDELCMD < 

%DIFDELCMD < %%%
\subsubsection{\DIFdel{Ensuring Weight Norm Boundedness}}
%DIFAUXCMD
\addtocounter{subsubsection}{-1}%DIFAUXCMD
%DIFDELCMD < 

%DIFDELCMD < %%%
\DIFdel{A common challenge in NAC is maintaining the boundedness of the NN weights to ensure stability.
In many studies, projection operators were employed to enforce upper bounds on the weight norms, ensuring that the weights do not grow unboundedly}\DIFdelend %DIF >  Since the NAC framework is based on Lyapunov theory, structural methods have been widely studied for input constraint handling.
\DIFaddbegin \DIFadd{Since the NAC framework is based on Lyapunov theory, input constraints have been primarily handled using structural methods in the NAC literature.
}\MSRY{NAC는 Lyapunov에 기반 $\to$ 입력제약을 구조적 방법으로 처리하면 Lyapunov로 해석 가능 $\to$ NAC가 구조적 방법으로 입력제약을 처리; 라는 흐름을 하고 싶었습니다; 단순히 NAC $\to$ 구조적 방법 (X) NAC $\to$ 입력제약을 구조적방법으로 처리 (O)}
\DIFadd{In \mbox{%DIFAUXCMD
\cite{Ji:2025aa}}\hskip0pt%DIFAUXCMD
, Ji }\etal \DIFadd{categorized the literature on the structural methods into two strategies: anti-saturation design and direct design}\DIFaddend .
\DIFaddbegin \DIFadd{We review these strategies in the context of NAC.
}

\DIFadd{In anti-saturation design, a virtual auxiliary system is mainly constructed to compensate for the effects of input saturation \mbox{%DIFAUXCMD
\cite{Ma:2025aa,Gao:2026aa,Yang:2021aa,Jia:2020aa,Liu:2024aa,Liu:2025aa,Gao:2017aa,Esfandiari:2014aa,Esfandiari:2015aa,Yu:2022aa,Rezaei:2025aa,He:2016aa,He:2019aa,Zhao:2016aa,Peng:2020aa}}\hskip0pt%DIFAUXCMD
.
Owing to their high compatibility, auxiliary systems have been adopted in various control designs.
}\DIFaddend For example, \DIFdelbegin \DIFdel{in \mbox{%DIFAUXCMD
\cite{Zhou:2023aa,Griffis:2023aa,Patil:2022aa}}\hskip0pt%DIFAUXCMD
, projection operators were used to constrain the weight norms to remain below predetermined constants.
However, these constants were often selected as large as possible due to the lack of theoretical guarantees regarding the global optimality of the weight values.
While this approach ensured that the NN remained stable, it did not necessarily result in optimal performance}\DIFdelend \DIFaddbegin \DIFadd{auxiliary systems have been incorporated into backstepping-based NAC \mbox{%DIFAUXCMD
\cite{Gao:2017aa} }\hskip0pt%DIFAUXCMD
and sliding mode control (SMC)-based NAC \mbox{%DIFAUXCMD
\cite{Ma:2025aa} }\hskip0pt%DIFAUXCMD
to mitigate the effects of input saturation.
Furthermore, the auxiliary systems have been integrated into reinforcement learning (RL)-based NAC designs \mbox{%DIFAUXCMD
\cite{Gao:2026aa} }\hskip0pt%DIFAUXCMD
and barrier Lyapunov function (BLF)-based NAC designs \mbox{%DIFAUXCMD
\cite{Yang:2021aa}}\hskip0pt%DIFAUXCMD
.
To achieve finite-time or fixed-time convergence, novel finite-time \mbox{%DIFAUXCMD
\cite{Jia:2020aa} }\hskip0pt%DIFAUXCMD
or fixed-time \mbox{%DIFAUXCMD
\cite{Liu:2024aa,Liu:2025aa} }\hskip0pt%DIFAUXCMD
auxiliary systems have been proposed}\DIFaddend .

\DIFdelbegin \DIFdel{In addition to projection operators, some studies utilized modification techniques like $\sigma$-modification \mbox{%DIFAUXCMD
\cite{Ge:2002aa} }\hskip0pt%DIFAUXCMD
and $\epsilon$-modification \mbox{%DIFAUXCMD
\cite{Esfandiari:2015aa,Gao:2006aa}}\hskip0pt%DIFAUXCMD
. 
These methods ensured that the NN weights remained within an invariant set by incorporating stabilizing functions into the adaptation law. 
Although these techniques were effective in ensuring boundedness and avoiding weight divergence, they similarly lacked a formal analysis of the optimality of the adapted weights by biasing the weights toward the origin \mbox{%DIFAUXCMD
\cite{Ryu:2024aa}}\hskip0pt%DIFAUXCMD
.
This leaves room for improvement in terms of performance optimization.
}%DIFDELCMD < 

%DIFDELCMD < \hfill
%DIFDELCMD < 

%DIFDELCMD < %%%
\subsubsection{\DIFdel{Satisfying Input Constraints}}
%DIFAUXCMD
\addtocounter{subsubsection}{-1}%DIFAUXCMD
%DIFDELCMD < 

%DIFDELCMD < %%%
\DIFdel{The second major issue is satisfying input constraints, particularly in systems where actuators are subject to physical limitations. 
This issue becomes more pronounced in NAC architectures where NNs are used to augment conventional model-based controllers—such as feedback linearization or backstepping—rather than acting as the sole control input.
In such cases, if }\DIFdelend \DIFaddbegin \DIFadd{Generally, when the control input reaches the saturation limit, }\DIFaddend the \DIFdelbegin \DIFdel{system model is inaccurate or the dynamics targeted by the NN are already contributing to stability, the combined control effort may result in overly aggressive inputs, increasing the risk of input saturation \mbox{%DIFAUXCMD
\cite{Khalil:2002aa}}\hskip0pt%DIFAUXCMD
.
}%DIFDELCMD < 

%DIFDELCMD < %%%
\DIFdel{To address the aforementioned issue, several approaches have been proposed in \mbox{%DIFAUXCMD
\cite{Gao:2006aa,Arefinia:2020aa,He:2016aa,Peng:2020aa,Esfandiari:2014aa,Karason:1994aa,Esfandiari:2015aa} }\hskip0pt%DIFAUXCMD
to handle input constraints in NAC systems.
Particularly, the auxiliary systems have been widely introduced.
In \mbox{%DIFAUXCMD
\cite{Arefinia:2020aa,He:2016aa,Peng:2020aa}}\hskip0pt%DIFAUXCMD
}\DIFdelend \DIFaddbegin \DIFadd{auxiliary system is excited, and the auxiliary state is fed back into the control law as additional compensation terms.
However, from a learning perspective, if the auxiliary state is not incorporated into the feedback error used in the adaptation law, the NN weights may generate excessive control input to achieve the desired tracking performance, since the saturation effect is not considered in the adaptation process.
Therefore}\DIFaddend , the auxiliary \DIFdelbegin \DIFdel{states were generated whenever input saturation was detected, and these states were used as feedback terms in the control law to directly compensate for the effects }\DIFdelend \DIFaddbegin \DIFadd{state is often considered in the adaptation law \mbox{%DIFAUXCMD
\cite{Gao:2017aa,Esfandiari:2014aa,Esfandiari:2015aa,Liu:2024aa,Liu:2025aa,Yu:2022aa,Rezaei:2025aa}}\hskip0pt%DIFAUXCMD
.
This ensures that the NN weights adapt not only for approximating the ideal control law but also for compensating for the effect }\DIFaddend of input saturation\DIFdelbegin \DIFdel{constraints.
}\DIFdelend \DIFaddbegin \DIFadd{.
}

\DIFaddend On the other hand, in \DIFdelbegin \DIFdel{\mbox{%DIFAUXCMD
\cite{Esfandiari:2014aa,Karason:1994aa,Esfandiari:2015aa}}\hskip0pt%DIFAUXCMD
, the auxiliary systems were }\DIFdelend \DIFaddbegin \DIFadd{direct design, the saturation effects are directly }\DIFaddend incorporated into the \DIFdelbegin \DIFdel{adaptation law by adding the auxiliary states to the feedback error for learning.
This approach helped the NN reduce input saturation by indirectly regulating the auxiliary states, during the adaptation process.
}\DIFdelend \DIFaddbegin \DIFadd{controller design and stability analysis.
Typically, the saturation function is approximated by a smooth function, such as $\tanh(\cdot)$, and the mean value theorem is used to transform it into affine form \mbox{%DIFAUXCMD
\cite{Wang:2013aa,Cui:2014aa,Li:2024aa}}\hskip0pt%DIFAUXCMD
.
However, as reported in \mbox{%DIFAUXCMD
\cite{Min:2021aa}}\hskip0pt%DIFAUXCMD
, the smooth approximation of the saturation function may introduce conservative behavior and performance degradation, especially when the control input is close to the saturation limit, compared to the auxiliary system-based anti-saturation design.
}\DIFaddend 

\DIFdelbegin \DIFdel{However, these approaches typically handle input bound constraints on a per-input basis, }%DIFDELCMD < \ie %%%
\DIFdel{applying bounds to each scalar control variable individually, and may not account for more complex, nonlinear constraints, like input norm constraints, which are commonly found in physical systems such as robotic actuators or motor systems.
}\DIFdelend \DIFaddbegin \DIFadd{Despite their effectiveness, the aforementioned structural methods are limited to element-wise input constraints, such as symmetric or asymmetric bounds on each input.
Thus, extending them to admissible input sets described by multiple convex constraints is not straightforward.
}\DIFaddend 

\DIFdelbegin \DIFdel{Furthermore, from a learning perspective, augmenting conventional control methods with NNs by simply adding NN outputs to existing control inputs can interfere with the adaptation process.
This occurs because the adaptation law typically depends on the feedback error, which becomes insensitive to the NN output when the conventional controller dominates the system behavior . 
As a result, the NN ends up counteracting the conventional control effort rather than directly learning the system dynamics, thereby weakening its ability to adapt.
To address this, input constraints should be handled within a unified adaptive framework—such as those proposed in \mbox{%DIFAUXCMD
\cite{Esfandiari:2014aa,Karason:1994aa,Esfandiari:2015aa}}\hskip0pt%DIFAUXCMD
—where both adaptation and constraint satisfaction are treated jointly without relying on conventional control inputs.
}\DIFdelend %DIF >  For instance, in \cite{Gao:2006aa}, the saturated amount of control input is estimated using NNs and compensated in the control law.

\DIFdelbegin %DIFDELCMD < \hfill
%DIFDELCMD < %%%
\DIFdelend %DIF >  ================================
%DIF >  OPTIMIZATION-BASED METHODS
%DIF >  ===============================

\DIFdelbegin \subsubsection{\DIFdel{Potential of Constrained Optimization}}
%DIFAUXCMD
\addtocounter{subsubsection}{-1}%DIFAUXCMD
%DIF <  \subsubsection{Limitations of Existing Approaches and Potential of Constrained Optimization}
%DIFDELCMD < 

%DIFDELCMD < %%%
%DIF <  \MSRY{각 문단에 이미 문제점을 지적해서 요약을 안해도 괜찮을 것 같습니다.}
%DIF <  \ctxt{red}
%DIF <  {
%DIF <      Although both the existing methods for weight boundedness and input constraints have shown effectiveness, they come with significant limitations. 
%DIF <      Projection operators and modification techniques do not guarantee the optimality of the adapted weights. 
%DIF <      Auxiliary systems typically handle only simple forms of input constraints, such as individual input bound constraints, limiting their ability to address more complex, nonlinear constraints.
%DIF <      Moreover, using conventional control methods in conjunction with NNs can disrupt the adaptation process.
\DIFdelend \DIFaddbegin \begin{table*}[t]
  \renewcommand{\arraystretch}{1.3}
  \caption{
    \DIFaddFL{Comparison of input-constraint handling approaches.
  }}
  \centering
  \begin{tabular}{c c c c c }
  %DIF >  \begin{tabular}{c c c c c}
  \hline
  & 
  \textbf{\DIFaddFL{Structural Methods}} 
  & 
  \textbf{\DIFaddFL{Optimization-Based Methods}} 
  & 
  \textbf{\DIFaddFL{Proposed Approach}}
  \\
  \hline
  \hline %DIF >  -----------------------------------
    \DIFaddFL{Representative 
  }& 
  \DIFaddFL{Auxiliary system design \mbox{%DIFAUXCMD
\cite{Gao:2017aa,Ma:2025aa,Gao:2026aa,Yang:2021aa,Jia:2020aa,Liu:2024aa,Liu:2025aa,Esfandiari:2014aa,Esfandiari:2015aa,Yu:2022aa,Rezaei:2025aa,He:2016aa,He:2019aa,Zhao:2016aa,Peng:2020aa}}\hskip0pt%DIFAUXCMD
,
  }&
  \DIFaddFL{MPC \mbox{%DIFAUXCMD
\cite{Greune:2011aa,Bemporad:1999aa,Berberich:2025aa,Draeger:1995aa,Vahidi-Moghaddam:2025aa,Mallick:2025aa,Khodaverdian:2025aa}}\hskip0pt%DIFAUXCMD
,
  }&
  \DIFaddFL{Preliminary works of this study
  }\\
  \DIFaddFL{literature
  }&
  \DIFaddFL{Smooth approximation of saturation \mbox{%DIFAUXCMD
\cite{Wang:2013aa,Cui:2014aa,Li:2024aa}
  }\hskip0pt%DIFAUXCMD
}&
  \DIFaddFL{CBF-QP filter \mbox{%DIFAUXCMD
\cite{Ames:2019aa,Ko:2025aa,Sweatland:2025aa}
  }\hskip0pt%DIFAUXCMD
}&
  \DIFaddFL{\mbox{%DIFAUXCMD
\cite{Ryu:2025ab,Ryu:2025aa}
  }\hskip0pt%DIFAUXCMD
}\\
  \hline %DIF >  -----------------------------------
  \multirow{2}{*}{
    Stability analysis 
  }
  & 
  \multirow{2}{*}{
    \textbf{Lyapunov-based stability guarantees}
  }
  & 
  \DIFaddFL{Conditions on feasibility and
  }& 
  \multirow{2}{*}{
    \textbf{Lyapunov-based stability guarantees}
  }
  \\
  &
  &
  \DIFaddFL{related assumptions required
  }&
  \\
  \hline %DIF >  -----------------------------------
  \DIFaddFL{Input constraints 
  }& 
  \DIFaddFL{Primarily element-wise 
  }& 
  \textbf{\DIFaddFL{General input constraints}}
  &
  \multirow{2}{*}{
  \textbf{Multiple convex input constraints}
  }
  \\
  \DIFaddFL{handling capability
  }&
  \DIFaddFL{(symmetric or asymmetric) input constraints
  }&
  \DIFaddFL{(convex/non-convex constraints)
  }&
  \\
  \hline %DIF >  -----------------------------------
  \multirow{2}{*}{
    Implementation
  }
  &
    \textbf{\DIFaddFL{Continuous adaptation}}
  &
  \multirow{2}{*}{
    Optimization at each time instant
  }
  &
    \textbf{\DIFaddFL{Continuous adaptation}}
  \\
  &
  \DIFaddFL{(typically Lyapunov-based)
  }&
  &
  \DIFaddFL{(constrained optimization-based)
  }\\
  \hline
  %DIF >  \multirow{2}{*}{
    \DIFaddFL{Computational
  }\DIFaddendFL % }
  \DIFaddbeginFL &
  \multirow{2}{*}{
    \textbf{Typically low}
  }
  &
  \DIFaddFL{Problem-dependent
  }&
  \multirow{2}{*}{
    \textbf{Typically low}
  }
  \\
  \DIFaddFL{burden
  }&
  &
  \DIFaddFL{(may require iterative solvers)
  }&
  \\
  \hline
  \end{tabular}
  \label{table:intro:comparison}
\end{table*}
\DIFaddend 


\DIFdelbegin \DIFdel{A promising approach to overcoming the limitations of the existing methods lies in constrained optimization .
By formulating the NAC problem as a unified optimization problemwith constraints , it is possible to adapt the NN weights while minimizing an objective function, }%DIFDELCMD < \eg %%%
\DIFdel{tracking error, subject to both weight and input constraints, simultaneously.
Constrained optimization provides a theoretical framework for defining optimality and presents numerical methods for finding solutions that satisfy the constraints \mbox{%DIFAUXCMD
\cite{Nocedal:2006aa}}\hskip0pt%DIFAUXCMD
.
}%DIFDELCMD < 

%DIFDELCMD < %%%
\DIFdel{In existing literature, constrained optimization techniques, such as the Augmented Lagrangian Method (ALM) \mbox{%DIFAUXCMD
\cite{Evens:2021aa} }\hskip0pt%DIFAUXCMD
and the Alternating Direction Method of Multipliers (ADMM) \mbox{%DIFAUXCMD
\cite{Wang:2019aa,Taylor:2016aa}}\hskip0pt%DIFAUXCMD
, have been used to train NNs offline.
%DIF <  These methods impose constraints to address issues like gradient vanishing in backpropagation. 
%DIF <  \MSRY{
%DIF <      여기를 수정해야 하지 않나요? 저희의 사전 연구가 있기 때문입니다.
%DIF <  }
%DIF <  \ctxt{red}
%DIF <  {
    %DIF <  However, to the best of the authors' knowledge, no prior work has applied constrained optimization to NAC systems with real-time weight adaptation. 
    %DIF <  }
%DIF <  \ctxt{blue}
%DIF <  {
However, there are few works that have applied constrained optimization to NAC systems with real-time weight adaptation \mbox{%DIFAUXCMD
\cite{Ryu:2024aa,Ryu:2025aa}}\hskip0pt%DIFAUXCMD
.
%DIF <  }
This gap suggests that constrained optimization could be key to addressing both weight boundedness and input constraints in a unified, theoretically grounded framework , particularly in real-time NAC.
}\DIFdelend %DIF >  \cite{Greune:2011aa,Bemporad:1999aa,Berberich:2025aa,Draeger:1995aa,Vahidi-Moghaddam:2025aa,Mallick:2025aa,Khodaverdian:2025aa}, \cite{Ames:2019aa,Ko:2025aa,Sweatland:2025aa}

\DIFdelbegin \subsection{\DIFdel{Contributions}}
%DIFAUXCMD
\addtocounter{subsection}{-1}%DIFAUXCMD
\DIFdelend %DIF >  MPC 
%DIF >  Bemporad:1999aa,Berberich:2025aa,Draeger:1995aa,Chen:2024aa,Cetin:2019aa,Vahidi-Moghaddam:2025aa

\DIFdelbegin \DIFdel{The main contributions of this study are listed as follows:
}%DIFDELCMD < \begin{itemize}
%DIFDELCMD <     \item %%%
\DIFdel{A constrained }\DIFdelend \DIFaddbegin \subsubsection{\DIFadd{Optimization-Based Methods for Input Constraint Handling}}

\DIFadd{Besides structural methods, }\DIFaddend optimization-based \DIFaddbegin \DIFadd{methods have been studied for constraint handling across a broad range of control design frameworks.
These methods compute the control input by solving constrained optimization problems at each time instant, as in model predictive control (MPC) \mbox{%DIFAUXCMD
\cite{Vahidi-Moghaddam:2025aa,Mallick:2025aa,Khodaverdian:2025aa,Greune:2011aa,Bemporad:1999aa,Berberich:2025aa,Draeger:1995aa} }\hskip0pt%DIFAUXCMD
and control barrier function (CBF)-based quadratic programming (QP) filters \mbox{%DIFAUXCMD
\cite{Ames:2019aa,Ko:2025aa,Sweatland:2025aa}}\hskip0pt%DIFAUXCMD
. 
In these methods, the constrained control problem is explicitly formulated as a constrained optimization problem.
Therefore, complex nonlinear (even non-convex) input constraints can be systematically incorporated into the control design, which is a significant advantage over the structural methods.
}

\DIFadd{Even though they solve constrained optimization problems at each time instant, the computational burden may not necessarily be significant for small-scale QPs or MPC problems with short horizons and simple constraints.
However, it can increase substantially when multiple constraints must be handled simultaneously, nonlinear constraints are introduced, or the resulting problem is non-convex.
In such cases, real-time implementation may require iterative numerical solvers, suitable initialization, and careful control of the problem dimension and solver convergence.
In this regard, some studies \mbox{%DIFAUXCMD
\cite{Vahidi-Moghaddam:2025aa,Mallick:2025aa,Khodaverdian:2025aa} }\hskip0pt%DIFAUXCMD
trained the optimal constrained control policy offline using deep NNs (DNNs) to reduce online computation.
}

\DIFadd{Moreover, closed-loop stability guarantees generally rely on additional analytical conditions specific to the selected optimization framework.
For MPC, these conditions commonly include recursive feasibility, appropriate terminal ingredients, sufficiently long prediction horizons, and computational tractability \mbox{%DIFAUXCMD
\cite{Greune:2011aa}}\hskip0pt%DIFAUXCMD
.
In addition to these conditions, model accuracy is critical, thereby motivating robust \mbox{%DIFAUXCMD
\cite{Bemporad:1999aa}}\hskip0pt%DIFAUXCMD
, data-driven \mbox{%DIFAUXCMD
\cite{Berberich:2025aa}}\hskip0pt%DIFAUXCMD
, and learning-based MPC \mbox{%DIFAUXCMD
\cite{Draeger:1995aa,Vahidi-Moghaddam:2025aa} }\hskip0pt%DIFAUXCMD
approaches.
Model accuracy is also crucial for CBF-QP filters, where the validity of the safety guarantee depends on the accuracy of the system dynamics used to formulate the barrier constraint.
Accordingly, }\DIFaddend neuro-adaptive \DIFdelbegin \DIFdel{control (CONAC) is proposed , where the control }\DIFdelend \DIFaddbegin \DIFadd{system identifiers have been introduced to construct the CBF constraint using approximated system dynamics \mbox{%DIFAUXCMD
\cite{Sweatland:2025aa}}\hskip0pt%DIFAUXCMD
.
}

\DIFadd{Although these approaches can provide stability or safety guarantees under their stated assumptions, the guarantees typically depend on feasibility, prediction, or filtering mechanisms whose validity is affected by model accuracy, uncertainty descriptions, data quality, and offline training coverage.
}

%DIF >  However, optimization-based methods require additional conditions to ensure closed-loop stability.
%DIF >  Representatively, in MPC design, the recursive feasibility, computational tractability, and sufficiently long prediction horizons are essential \cite{Greune:2011aa}.
%DIF >  Since the model accuracy is most critical across these conditions, the uncertain or unknown dynamics have motivated robust \cite{Bemporad:1999aa}, data-driven \cite{Berberich:2025aa}, and learning-based MPC \cite{Draeger:1995aa,Vahidi-Moghaddam:2025aa}.
%DIF >  For CBF-QP filters, neuro-adaptive system identifiers have been employed to accurately impose the CBF constraint using approximated system dynamics \cite{Sweatland:2025aa}.
%DIF >  Additionally, for computational tractability, some studies have proposed learning the optimal control policy offline using NNs \cite{Vahidi-Moghaddam:2025aa,Mallick:2025aa,Khodaverdian:2025aa}.
%DIF >  Although these methods can provide stability guarantees under the stated conditions, the guarantees are typically obtained through feasibility, prediction, or filtering mechanisms that depend on model accuracy, uncertainty descriptions, data quality, or offline training coverage.

\subsection{\DIFadd{Motivation and Proposed Approach}}

\DIFadd{Both structural and optimization-based methods are effective for input constraint handling, but they have complementary limitations, as summarized in Table~\ref{table:intro:comparison}.
Structural methods offer Lyapunov-based stability guarantees and solver-free computation, but they are limited in terms of the classes of constraints they can accommodate.
Optimization-based methods can handle general constraints, but their closed-loop stability guarantees are often tied to additional conditions and iterative solvers may be required.
}

\DIFadd{These observations motivate a unified NAC framework in which more complex input constraints are incorporated directly into the online adaptation law, as presented in our previous work \mbox{%DIFAUXCMD
\cite{Ryu:2025aa}}\hskip0pt%DIFAUXCMD
.
In the proposed CONAC framework, the NAC }\DIFaddend problem is formulated as a constrained optimization problem \DIFdelbegin \DIFdel{. In this formulation, weight and }\DIFdelend \DIFaddbegin \DIFadd{in which multiple }\DIFaddend convex input constraints are \DIFdelbegin \DIFdel{incorporated }\DIFdelend \DIFaddbegin \DIFadd{imposed }\DIFaddend as inequality constraints\DIFdelbegin \DIFdel{, while minimizing a given objective function.
}%DIFDELCMD < \item %%%
\DIFdel{Convex input constraints can be applied to the proposed CONAC, including input bound constraints and norm constraints.
Moreover}\DIFdelend \DIFaddbegin \DIFadd{.
The constrained optimization formulation is not used to solve a control allocation problem at each time instant; instead, it is used to derive continuous-time adaptation laws for the adaptive variables.
Therefore}\DIFaddend , the proposed \DIFdelbegin \DIFdel{CONAC can handle any combination of these constraints , which can be useful in practical applications where multiple constraints need to be satisfied simultaneously.
    }\DIFdelend \DIFaddbegin \DIFadd{framework retains the Lyapunov-friendly structure of NAC while incorporating the constraint-handling capability of optimization-based methods.
}

\subsection{\DIFadd{Contributions}}

\DIFadd{The main contributions of this study are:
}\begin{itemize}
  %DIF >  \item 
  %DIF >  A constrained optimization-based neuro-adaptive control (CONAC) framework is proposed for unknown systems under multiple convex input constraints;
  %DIF >  The admissible input set is represented by convex inequality constraints and embedded directly into the online adaptation law;

  \DIFaddend \item
  \DIFdelbegin \DIFdel{Since the proposed CONAC approximates the entire ideal control law without any conventional controller, prior knowledge of the system dynamics or the system identification process for designing a nominal conventional controller can be avoided, which is often difficult and time-consuming to implement in practice. 
    }\DIFdelend \DIFaddbegin \DIFadd{Any combination of convex input constraints can be handled, thereby significantly extending the applicability of the NAC to real-world systems with complex input constraints;
  }

  \DIFaddend \item 
  The adaptation laws for \DIFaddbegin \DIFadd{adaptive variables, including the }\DIFaddend DNN weights and \DIFaddbegin \DIFadd{the }\DIFaddend Lagrange multipliers are systematically derived \DIFdelbegin \DIFdel{to solve the defined problem, using constrained optimization theory. These laws guarantee convergence to the first-order optimality conditions at steady state, specifically the Karush–Kuhn–Tucker }\DIFdelend \DIFaddbegin \DIFadd{from the formulated constrained optimization problem, and the Karush-Kuhn-Tucker }\DIFaddend (KKT) \DIFdelbegin \DIFdel{conditions, as described in }\DIFdelend \DIFaddbegin \DIFadd{optimality conditions are satisfied at an equilibrium of the adaptation dynamics }\DIFaddend \cite[Chap. 12, Thm. 12.1]{Nocedal:2006aa}\DIFdelbegin \DIFdel{.
    }\DIFdelend \DIFaddbegin \DIFadd{; and
}

  \DIFaddend \item 
  The \DIFdelbegin \DIFdel{proposed CONAC's stability }\DIFdelend \DIFaddbegin \DIFadd{closed-loop stability under the proposed Constrained Optimization-Based Neuro-Adaptive Control (CONAC) }\DIFaddend is rigorously analyzed using Lyapunov theory, ensuring \DIFdelbegin \DIFdel{that the }\DIFdelend \DIFaddbegin \DIFadd{the uniform ultimate boundedness (UUB) of the }\DIFaddend tracking error and \DIFdelbegin \DIFdel{DNN weights remain bounded.
This analysis guarantees that the proposed CONAC can achieve stable online adaptation.
    }%DIFDELCMD < \item %%%
\item%DIFAUXCMD
\DIFdel{The proposed CONAC is implemented in real-time on a 2-DOF robotic manipulator, demonstrating its effectiveness in achieving accurate trajectory tracking while satisfying all imposed constraints. The experimental results validate the theoretical findings and confirm the practical applicability of the proposed CONAC.
}\DIFdelend \DIFaddbegin \DIFadd{adaptive variables.
}

\DIFaddend \end{itemize}

\DIFdelbegin %DIFDELCMD < \hfill
%DIFDELCMD < 

%DIFDELCMD < %%%
\DIFdelend This study extends our preliminary work in \DIFdelbegin \DIFdel{\mbox{%DIFAUXCMD
\cite{Ryu:2025aa}}\hskip0pt%DIFAUXCMD
}\DIFdelend \DIFaddbegin \DIFadd{\mbox{%DIFAUXCMD
\cite{Ryu:2025ab,Ryu:2025aa}}\hskip0pt%DIFAUXCMD
}\DIFaddend , which introduced inequality constraints on the weight and input norms within a constrained optimization framework for \DIFdelbegin \DIFdel{neuro-adaptive control (NAC ) systems using a SHLNN}\DIFdelend \DIFaddbegin \DIFadd{NAC using single-hidden-layer NNs}\DIFaddend . 
In the present work, the same framework is generalized to accommodate \DIFdelbegin \DIFdel{arbitrary }\DIFdelend \DIFaddbegin \DIFadd{multiple }\DIFaddend convex input constraints and extended to DNNs, while still \DIFdelbegin \DIFdel{retaining the capability to regulate the weight norm}\DIFdelend \DIFaddbegin \DIFadd{ensuring Lyapunov-based stability guarantees}\DIFaddend .
Furthermore, unlike the previous study, which was limited to numerical simulations, the proposed CONAC is experimentally validated through real-time implementation on a \DIFdelbegin \DIFdel{2-DOF }\DIFdelend \DIFaddbegin \DIFadd{2-degree-of-freedom (DOF) }\DIFaddend robotic manipulator.

\subsection{Organization}

The remainder of this paper is organized as follows. 
Section\DIFdelbegin \DIFdel{\ref{sec:Problem Formulation} }\DIFdelend \DIFaddbegin \DIFadd{~\ref{sec:problem:formulation} }\DIFaddend presents the target system and control objective.
Section\DIFdelbegin \DIFdel{\ref{sec:ctrl design} }\DIFdelend \DIFaddbegin \DIFadd{~\ref{sec:main:result} }\DIFaddend introduces the proposed CONAC\DIFdelbegin \DIFdel{and the architecture of the DNN used in the controller. 
In Section \ref{sec:adap_laws}, the adaptation law is developed.
Section\ref{sec:stability} analyzes the stability of CONAC.
Section\ref{sec:sim} }\DIFdelend \DIFaddbegin \DIFadd{, including the adaptation laws for adaptive variables.
Section~\ref{sec:stability:analysis} analyzes the closed-loop stability.
Section~\ref{sec:val} }\DIFaddend presents a real-time implementation of CONAC on a 2-DOF robotic manipulator\DIFaddbegin \DIFadd{, and compares it with existing methods}\DIFaddend .
Finally, Section\DIFaddbegin \DIFadd{~}\DIFaddend \ref{sec:conclusion} concludes the paper and discusses potential future work.
\DIFdelbegin \DIFdel{Candidates for convex input constraints are provided in Appendix \ref{sec:appen:cstr}.
}\DIFdelend 

%DIF <  ====================================
\DIFaddbegin \subsubsection*{\DIFadd{Notations}}

\DIFadd{$\otimes$ denotes the Kronecker product \mbox{%DIFAUXCMD
\cite[Chap. 7, Def. 7.1.2]{Bernstein:2009aa}}\hskip0pt%DIFAUXCMD
.
A vector and a matrix are denoted by $\mv{x}=[x_i]_{i\in\{1,\cdots,n\}}\in\R^n$ and 
$
  \mm A
  := 
  [a_{ij}]_{
    i\in\{1,\cdots,n\},j\in\{1,\cdots ,m\}
  }\in\R^{n\times m}
  ,
$
respectively.
The $i$th column of the matrix $\mm{A}\in\R^{n\times m}$ is denoted by $\mycol_i(\mm A)$.
For the matrix $\mm{A}\in\R^{n\times m}$, $\myvec(\mm A):=(\mycol_1(\mm A)^\top,\cdots,\mycol_m(\mm A)^\top )^\top\in\R^{nm}$ denotes the vectorization of $\mm A$.
Finally, $\mm I_n$ denotes the $n\times n$ identity matrix, and $\mm 0_{n\times m}$ denotes the $n\times m$ zero matrix.
}

%DIF >  ------------------------------------
\DIFaddend %       SECTION PROBLEM FORMULATION
%DIF <  ====================================
%DIF >  ------------------------------------
\section{Problem Formulation} \DIFdelbegin %DIFDELCMD < \label{sec:Problem Formulation}
%DIFDELCMD < %%%
\DIFdelend \DIFaddbegin \label{sec:problem:formulation}
\DIFaddend 

\DIFdelbegin \subsection{\DIFdel{Model Dynamics and Control Objective}}
%DIFAUXCMD
\addtocounter{subsection}{-1}%DIFAUXCMD
%DIFDELCMD < 

%DIFDELCMD < %%%
\DIFdelend We consider an unknown \DIFdelbegin \DIFdel{Euler-Lagrange }\DIFdelend \DIFaddbegin \DIFadd{multi-input multi-output (MIMO) }\DIFaddend system modeled as 
\DIFdelbegin %DIFDELCMD < \begin{equation}
%DIFDELCMD <     \rbM\ddq +\rbVm \dq+ \rbF + \rbG + \rbd
%DIFDELCMD <     =
%DIFDELCMD <     \mysat(\rbu)
%DIFDELCMD <     ,
%DIFDELCMD <     \label{eq:sys1}
%DIFDELCMD < \end{equation}%%%
\DIFdelend \DIFaddbegin \begin{equation} \label{eq:sys}
  \ddtt\mv{x} 
  = 
  \mv{f}(\mv{x}) 
  + 
  \mm{B}(\mv{x})\mysat(\mv{u})
  ,
\end{equation}\DIFaddend 
where \DIFdelbegin \DIFdel{$\q\in \R^n$ denotes the generalized coordinate, $\rbd$ represents the disturbance and $\rbu\in\R^n$ denotes the generalized control input . 
The terms $\rbM:=\rbM(\q)\in\R^{n\times n}$}\DIFdelend \DIFaddbegin \DIFadd{$\mv{x},\mv{u}\in\R^n$ denote the state and control input vectors, respectively.
The unknown system function and control effectiveness matrix are denoted by $\mv{f}:\R^n\to\R^n$ and $\mm{B}:\R^n\to\R^{n\times n}$, respectively}\DIFaddend , \DIFdelbegin \DIFdel{$\rbVm:=\rbVm(\q,\dq)\in\R^{n\times n}$, $\rbF:=\rbF(\dq)\in\R^{n}$, and $\rbG:=\rbG(\q)\in\R^{n}$ represent the unknown inertia matrix , Coriolis/centripetal matrix, friction vector, and gravity vector, respectively.
The function $\mysat(\cdot):\R^n\to\R^n$ represents the inherent physical limitations of the actuators such that $\norm{\mysat(\rbu)}\le \overline\tau$, where $\overline\tau\in\R_{>0}$ denotes the maximum norm of control input.
}%DIFDELCMD < 

%DIFDELCMD < %%%
\DIFdel{In this study, the saturation function $\mysat(\cdot)$ is }\DIFdelend \DIFaddbegin \DIFadd{which are }\DIFaddend assumed to be \DIFdelbegin \DIFdel{convex, which is practically common in many engineering systems.
This is because actuator limits are typically convex—for example, minimum and maximum torque constraints—which define box- or ball-shaped bounds in the input space.
To account for these limitations, it is essential to incorporate physically motivated constraints into the controller design.
Appendix \ref{sec:appen:cstr} introduces candidate constraints that can be applied to ensure compliance with these physical limitations.
}\DIFdelend \DIFaddbegin \DIFadd{locally Lipschitz continuous.
The saturation function is denoted by $\mysat:\R^n\to\Omega_u$, where $\Omega_u\subset\R^n$ is the convex admissible input set.
}\DIFaddend 

\DIFdelbegin \DIFdel{The Euler-Lagrange system \eqref{eq:sys1} satisfies several important physical properties, as presented in \mbox{%DIFAUXCMD
\cite[Chap. 3, Tab. 3.2.1]{Lewis:1998aa}}\hskip0pt%DIFAUXCMD
. 
The key properties are introduced below:
}%DIFDELCMD < \begin{prop} 
%DIFDELCMD <     %%%
\DIFdel{The inertia matrix $\rbM$ is symmetric, positive definite, and bounded .
    }%DIFDELCMD < \label{prop:M}
%DIFDELCMD < \end{prop}
%DIFDELCMD < %%%
\DIFdelend \DIFaddbegin \DIFadd{The system \eqref{eq:sys} is subject to the following assumptions.
}\DIFaddend 

\DIFdelbegin %DIFDELCMD < \begin{prop} 
%DIFDELCMD <     %%%
\DIFdel{The Coriolis/centripetal matrix $\rbVm$ can always be selected, so that the matrix $\ddtt{\rbM}-2\rbVm$ is skew-symmetric, }%DIFDELCMD < \ie %%%
\DIFdel{$\mv{x}^\top(\ddtt{\rbM}-2\rbVm)\mv{x}=0,\ \forall \mv{x}\in\R^n$. 
}%DIFDELCMD < \label{prop:skew}
%DIFDELCMD < \end{prop}
%DIFDELCMD < %%%
\DIFdelend \DIFaddbegin \begin{assum} \label{assum:B:bound}
  \DIFadd{The control effectiveness matrix $\mm{B}(\mv{x})$ is nonsingular and bounded as   
  $
    \underline{b}\mm{I}_n
    \preceq 
    \mm{B}(\mv{x})
    \preceq 
    \bar{b}\mm{I}_n
    ,
  $
  where $\underline{b},\bar{b}\in\R_{>0}$ are unknown positive constants.
}\end{assum}
\DIFaddend 

\DIFdelbegin %DIFDELCMD < \begin{prop}
%DIFDELCMD <     %%%
\DIFdel{The disturbance $\rbd$ is bounded, }%DIFDELCMD < \ie %%%
\DIFdel{$\norm{\rbd}\le \overline\tau_d$, }\DIFdelend \DIFaddbegin \begin{assum}[{\cite[Assumption 1(b)]{Boulkroune:2010aa}}] \label{assum:B:bound:grad}
  \DIFadd{There exists an unknown positive function $\beta(\mv{x}):\R^n\to\R_{>0}$ such that $\tfrac{1}{2}\mv{v}^\top\left(\ddtt\mm{B}(\mv{x})^{-1}\right)\mv{v} \le \beta(\mv{x})\norm{\mv{v}}^2$, $\forall \mv{v}\in\R^n$, and $\forall \mv{x}\in\R^n$. 
}\end{assum}

\begin{assum} \label{assum:sat}
  \DIFadd{The admissible input set and saturation function satisfy the following
  conditions:
  }\begin{enumerate}
    \item
    \DIFadd{The admissible input set $\Omega_u$ is defined by smooth convex
    constraint functions $c_j^u:\R^n\to\R$, $j\in\mathcal I$, as
    }\begin{equation*}
      \Omega_u
      =
      \textstyle
      \bigcap_{j\in\mathcal I}
      \left\{
        \mv u\in\R^n:
        c_j^u(\mv u)\le 0
      \right\}
      ,
    \end{equation*}
    \DIFadd{where $\mathcal I$ denotes the input constraint index set.
}

    \item
    \DIFadd{The origin lies in the interior of $\Omega_u$, i.e.,
    }\begin{equation*}
      c_j^u(\mv 0)
      \le
      -\underline c_j^u
      <
      0
      ,
      \qquad
      \forall j\in\mathcal I
      ,
    \end{equation*}
    \DIFaddend where \DIFdelbegin \DIFdel{$\overline\tau_d\in\R_{\ge0}$.
}%DIFDELCMD < \label{prop:dis_bound}
%DIFDELCMD < \end{prop}
%DIFDELCMD < %%%
\DIFdelend \DIFaddbegin \DIFadd{$\underline c_j^u\in\R_{>0}$ is known.
}\DIFaddend 

    \DIFdelbegin \DIFdel{The control objective is to develop a NAC that
    enables $\q$ to track a continuously differentiable desired trajectory $\qd:=\qd(t): \R_{\ge0} \to \R^n$, compensating for the unknown system dynamics while addressing the imposed constraints.
The desired trajectory $\qd(t)$ is supposed to satisfy }\DIFdelend \DIFaddbegin \item
    \DIFadd{There exists a known constant $\underline u\in\R_{>0}$ such that
    }\begin{equation*}
      \inf_{\mv u\notin\Omega_u}
      \left(
        \norm{\mysat(\mv u)}
      \right)
      =
      \underline u
      .
    \end{equation*}
  \end{enumerate}
\end{assum}

%DIF >  \begin{assum} \label{assum:sat}
%DIF >      The admissible input set $\Omega_u$ is defined by smooth convex constraint functions $c_j^u(\mv u)$, $\forall j\in\mathcal I$, where $\mathcal{I}$ denotes the input constraint index set, as
%DIF >      $
%DIF >          \Omega_u
%DIF >          =
%DIF >          \bigcap_{j\in\mathcal I}
%DIF >          \{\mv u\in\R^n:c_j^u(\mv u)\le0\}
%DIF >          .
%DIF >      $
%DIF >      The origin is contained in its interior, \ie $c_j^u(\mv 0) \le -\underline{c}_j^u <0$, $\forall j\in\mathcal I$, for known $\underline{c}_j^u\in\R_{>0}$.
%DIF >      Moreover, there exists known $\underline u\in\R_{>0}$ such that
%DIF >      $
%DIF >          \inf_{\mv{u}\notin\Omega_u}
%DIF >          \left(
%DIF >              \norm{\mysat(\mv u)}
%DIF >          \right)
%DIF >          =
%DIF >          \underline u
%DIF >          .
%DIF >      $
%DIF >  \end{assum}
%DIF >  The saturation function $\mysat(\cdot)$ satisfies $\mysat(\mv u)=\mv u$ for $\mv u\in\Omega_u$ and, for $\mv u\notin\Omega_u$, $\mysat(\mv u)$ lies on the ray from the origin through $\mv u$, as illustrated in Fig.~\ref{fig:saturation}.

\DIFadd{By Assumption~\ref{assum:sat}, }\DIFaddend the \DIFaddbegin \DIFadd{saturation function $\mysat(\mv{u})$ can represent any combination of multiple convex input constraints, rather than simple element-wise constraints which are conventionally considered in the NAC literature, as illustrated in Fig.~\ref{fig:saturation}.
}

\begin{figure}[t]
    \centering
    %DIF >  \subfloat[]{
    \subfloat[Conventional input constraints.]{
        \includegraphics[width=0.5\linewidth]{figs/inputFuncConv.drawio.pdf}
        \label{fig:saturation:conv:box}
    }
    %DIF >  \hfill
    %DIF >  \subfloat[]{
    \subfloat[Multiple convex input constraints.]{
        \includegraphics[width=0.47\linewidth]{figs/inputFuncProp.drawio.pdf}
        \label{fig:saturation:prop:multi}
    }
    \caption{
      \DIFaddFL{Comparison of conventionally considered input constraints in the NAC literature and the proposed multiple convex input constraints.
      }\MSRY{문헌에서 box라고 표현하지 않고 symmetric/asymmetric라고 표현하기에 그에 맞게 표현하였습니다.}
      %DIF >  In Fig.~\ref{fig:saturation:conv:box}, symmetric and asymmetric box constraints are illustrated.
      %DIF >  In Fig.~\ref{fig:saturation:prop:multi}, a combination of symmetric and norm constraint are illustrated.
    }
    \label{fig:saturation}
\end{figure}

\DIFadd{Additionally, the desired trajectory $\mv{x}_d(t):\R_{\ge0}\to\R^n$ is assumed to be given and satisfies the }\DIFaddend following assumption\DIFdelbegin \DIFdel{:
}\DIFdelend \DIFaddbegin \DIFadd{.
}

\DIFaddend \begin{assum} \DIFaddbegin \label{assum:xd}
    \DIFaddend The desired trajectory \DIFdelbegin \DIFdel{$\qd(t)$ is assumed to be bounded, }%DIFDELCMD < \ie %%%
\DIFdel{$\norm{\qd(t)}\le \overline{q}_d\in\R_{>0}$, and available for designing a boundedcontrol input}\DIFdelend \DIFaddbegin \DIFadd{$\mv{x}_d(t)$ is sufficiently smooth, and $\mv{x}_d$ and $\ddtt\mv{x}_d$ are bounded}\DIFaddend .
    \DIFdelbegin %DIFDELCMD < \label{assum:feasible}
%DIFDELCMD < %%%
\DIFdelend %DIF >  The desired trajectory $\mv{x}_d(t)$ is sufficiently smooth and bounded, such that $\mv{x}_d\in\Omega_{x_d}:=\left\{\mv{x}_d\in\R^n : \norm{\mv{x}_d}\le \overline{x}_d\right\}$ and $\ddtt\mv{x}_d$ is bounded in a compact set $\Omega_{dx_d}$, as well.
\end{assum}

%DIF <   SECTION CONTROLLER DESIGN ==============================
\DIFaddbegin \DIFadd{The control objective is to design a NAC for the unknown system \eqref{eq:sys} such that $\mv{x}(t)$ tracks the desired trajectory $\mv{x}_d(t)$, the multiple convex input constraints defined by $\mysat(\cdot)$ are satisfied, and closed-loop stability is guaranteed.
}

%DIF >  ------------------------------------
%DIF >        SECTION MAIN RESULT
%DIF >  ------------------------------------
\DIFaddend \section{\DIFdelbegin \DIFdel{Control Law Development}\DIFdelend \DIFaddbegin \DIFadd{Proposed CONAC Framework}\DIFaddend } \DIFdelbegin %DIFDELCMD < \label{sec:ctrl design}
%DIFDELCMD < %%%
\DIFdelend \DIFaddbegin \label{sec:main:result}
\DIFaddend 

The architecture of the proposed CONAC consists of a DNN that functions as a NAC and \DIFdelbegin \DIFdel{a weight optimizer for the DNN}\DIFdelend \DIFaddbegin \DIFadd{an adaptation process for the adaptive variables such as the DNN weights and Lagrange multipliers}\DIFaddend , as illustrated in Fig.\DIFaddbegin \DIFadd{~}\DIFaddend \ref{fig:ctrl:diagram}.
Before \DIFdelbegin \DIFdel{proceeding to the controller design, some mathematical preliminaries are revisited in Section\ref{sec:sub:math preliminaries}.
Section \ref{sec:sub:NAC} introduces the NAC, followed by the DNN model in Section\ref{sec:sub:NN definition}. 
The constrained optimization-based weight adaptation law is presented in Section\ref{sec:adap_laws}}\DIFdelend \DIFaddbegin \DIFadd{presenting CONAC, we first review some preliminaries in Section~\ref{sec:main:result:sub:prelim}.
Then, the controller is presented in Section~\ref{sec:main:result:sub:ctrl:design} along with the adaptation laws in Section~\ref{sec:main:result:sub:adap:laws}}\DIFaddend .

\begin{figure}[!t]
  \centering
  \DIFdelbeginFL %DIFDELCMD < \includegraphics[width=.99\linewidth]{src/figures/Controller.drawio.pdf}
%DIFDELCMD <     %%%
\DIFdelendFL \DIFaddbeginFL \includegraphics[width=.9999\linewidth]{figs/Controller.drawio.pdf}
  \DIFaddendFL \caption{\DIFdelbeginFL \DIFdelFL{Architecture }\DIFdelendFL 
    \DIFaddbeginFL \DIFaddFL{Overall architecture }\DIFaddendFL of the proposed \DIFdelbeginFL \DIFdelFL{constrained optimization-based neuro-adaptive controller (}\DIFdelendFL CONAC\DIFdelbeginFL \DIFdelFL{)}\DIFdelendFL .
    \DIFaddbeginFL \MSRY{키워봤습니다.}
  \DIFaddendFL }
  \label{fig:ctrl:diagram}
\end{figure}

\subsection{\DIFdelbegin \DIFdel{Mathematical }\DIFdelend Preliminaries} \DIFdelbegin %DIFDELCMD < \label{sec:sub:math preliminaries}
%DIFDELCMD < %%%
\DIFdelend \DIFaddbegin \label{sec:main:result:sub:prelim}
\DIFaddend 

The following proposition \DIFdelbegin \DIFdel{and lemma }\DIFdelend will be used in \DIFaddbegin \DIFadd{the }\DIFaddend subsequent sections:
%DIF >  The following propositions and lemmas will be used in the subsequent sections:

\DIFdelbegin %DIFDELCMD < \begin{propsit}[see {\cite[Chap. 7, Prop. 7.1.9]{Bernstein:2009aa}}] %%%
\DIFdelend \DIFaddbegin \begin{propsit}[{\cite[Prop. 7.1.9]{Bernstein:2009aa}}] \DIFaddend \label{propsit:kron}
  For a matrix $\mm A\in\R^{n\times m}$ and a vector $\mv b\in\mathbb R^n$, the following property holds:
  \begin{equation}
    \mm A^\top \mv b 
    = 
    \myvec(\mm A^\top\mv b)
    =
    \myvec(\mv b^\top\mm A)
    = 
    (\mm{I}_{m}\otimes \mv b^\top) \myvec(\mm A)
    .
  \end{equation}
\end{propsit}

\DIFdelbegin %DIFDELCMD < \begin{lem} \label{lem:stable:set}
%DIFDELCMD <     %%%
\DIFdel{For a Lyapunov function $V:=V(\mv{x})\ge0$ with $\mv{x}\in\R^n$, if the time derivative of $V$ is given by $\ddtt V\le -a_1 \mv{x}^\top \mv{x} + a_2 \mv{x}$ for some positive constants $a_1\in\R_{>0}$ and $a_2\in\R_{>0}$, then $\mv{x}$ is bounded as $\mv{x}\in\{\mv{x}\mid\norm{\mv{x}}\le\tfrac{a_2}{a_1}\}$.
}%DIFDELCMD < \end{lem}
%DIFDELCMD < %%%
\DIFdelend %DIF >  \begin{lem}[Comparison Lemma {\cite[pp. 102--103, Lemma 3.4]{Khalil:2002aa}}] \label{lem:comp}
%DIF >      Consider a scalar continuous function $V(t)\ge0,\ \forall t\in\R_{\ge0}$.
%DIF >      If $V(t)$ satisfies the differential inequality    
%DIF >      % \begin{equation}
%DIF >      $
%DIF >          \ddtt V(t) \le -a V(t) + b
%DIF >          ,
%DIF >      % \end{equation}
%DIF >      $
%DIF >      with $a,b\in\R_{>0}$, then $V(t)$ is bounded as 
%DIF >      % \begin{equation}
%DIF >      $
%DIF >          V(t) \le \left(V(0) - \tfrac{b}{a}\right)\exp^{-at} + \tfrac{b}{a}
%DIF >          ,
%DIF >          \ \forall t\in\R_{\ge0}
%DIF >          .
%DIF >      % \end{equation}
%DIF >      $
%DIF >  \end{lem}

\DIFdelbegin %DIFDELCMD < \begin{proof}
%DIFDELCMD <     %%%
\DIFdel{The time derivative of $V$ can be rewritten as
    }%DIFDELCMD < \begin{equation}
%DIFDELCMD <         \ddtt V
%DIFDELCMD <         \le
%DIFDELCMD <         -a_1 \norm{\mv{x}}^2 + a_2 \norm{\mv{x}}
%DIFDELCMD <         =
%DIFDELCMD <         -a_1
%DIFDELCMD <         \left(
%DIFDELCMD <             \norm{\mv{x}}-\tfrac{a_2}{2a_1}
%DIFDELCMD <         \right)^2
%DIFDELCMD <         +
%DIFDELCMD <         \tfrac{a_2^2}{4a_1}
%DIFDELCMD <         .
%DIFDELCMD <         \label{eq:lem:lya:dot}
%DIFDELCMD <     \end{equation}
%DIFDELCMD <     %%%
\DIFdel{According to \eqref{eq:lem:lya:dot}, it can be concluded that $\ddtt V$ is negative definite if $\norm{\mv{x}}>\tfrac{a_2}{a_1}$ holds.
    In other words, $\mv{x}$ remains within the bounded set $\{\mv{x}\mid\norm{\mv{x}}\le\tfrac{a_2}{a_1}\}$, since $\ddtt V$ is negative definite when $\mv{x}$ is outside the set.
}%DIFDELCMD < \end{proof}
%DIFDELCMD < %%%
\DIFdelend %DIF >  \begin{proof}
%DIF >      See \cite[pp. 689--660]{Khalil:2002aa}.
%DIF >  \end{proof}

\DIFdelbegin \subsection{\DIFdel{Neuro-Adaptive Controller Design}}%DIFAUXCMD
\addtocounter{subsection}{-1}%DIFAUXCMD
%DIFDELCMD < \label{sec:sub:NAC}
%DIFDELCMD < %%%
\DIFdelend %DIF >  Based on the Comparison Lemma (Lemma~\ref{lem:comp}), the following lemma establishes the uniform ultimate boundedness (UUB) of the state vector $\mv{z}(t):\R_{\ge0}\to\R^n$ within the constrained set $\mathcal{Z}:=\left\{\mv{z}\in\R^n : \norm{\mv{z}}\le\overline{z}\in\R_{>0}\right\}$.

\DIFdelbegin \DIFdel{Given that the Euler-Lagrange system \eqref{eq:sys1} exhibits second-order dynamics, a filtered error $\fe\in\R^n$ is introduced to convert the system into a first-order system, as follows:
}%DIFDELCMD < \begin{equation}
%DIFDELCMD <     \fe := \ddtt{\mv e}+\mm\Lambda \mv e
%DIFDELCMD <     ,
%DIFDELCMD <     \label{eq:err:filtered}
%DIFDELCMD < \end{equation}
%DIFDELCMD < %%%
\DIFdel{where $\mv e := \q - \qd$ denotes the tracking error, $\ddtt{\mv e} := \dq - \dqd$ represents the time derivative of the tracking error, and $\mm\Lambda\in\R^{n\times n}_{>0}$ is a user-designed filtering matrix.
Since \eqref{eq:err:filtered} is stable system, it implies that $\mv e$ is bounded if $\fe$ is bounded.
}\DIFdelend %DIF >  \begin{lem} \label{lem:comp:compact}
%DIF >      Consider a scalar continuous function $V(\mv{z})$, such that $\gamma_1\norm{\mv{z}}^2 \le V(\mv{z}) \le \gamma_2\norm{\mv{z}}^2$, where $\gamma_1,\gamma_2\in\R_{>0}$.
%DIF >      If $V(\mv{z})$ satisfies    
%DIF >      \begin{equation}
%DIF >          \ddtt V(\mv{z}) 
%DIF >          \le 
%DIF >          -\gamma_3\norm{\mv{z}}^2 + c
%DIF >          ,\ 
%DIF >          \forall t\in\R_{\ge0}
%DIF >          ,\ 
%DIF >          \forall \mv{z}(t)\in\mathcal{Z}
%DIF >          ,
%DIF >      \end{equation}
%DIF >      where $\gamma_3,c\in\R_{>0}$, then $\mv{z}(t),\ \forall \mv{z}\in\mathcal{Z}$, is bounded as 
%DIF >      \begin{equation}
%DIF >          \norm{\mv{z}(t)} 
%DIF >          \le 
%DIF >          \sqrt{
%DIF >              \left(
%DIF >                  \tfrac{\gamma_2}{\gamma_1}    
%DIF >                  \norm{\mv{z}(0)}^2
%DIF >                  - 
%DIF >                  \tfrac{\gamma_2 c}{\gamma_1\gamma_3}
%DIF >              \right)
%DIF >              \exp(-\tfrac{\gamma_3}{\gamma_2}t) 
%DIF >              + 
%DIF >              \tfrac{\gamma_2 c}{\gamma_1\gamma_3}
%DIF >          }
%DIF >          .
%DIF >      \end{equation}
%DIF >      Conversely, to ensure that $\mv{z}(t)\in\mathcal{Z},\ \forall t\in\R_{\ge0}$, it is sufficient to satisfy
%DIF >      % $
%DIF >      %     \sqrt{
%DIF >      %         \tfrac{\gamma_2}{\gamma_1}    
%DIF >      %         \norm{\mv{z}(0)}^2 
%DIF >      %         +
%DIF >      %         \tfrac{\gamma_2 c}{\gamma_1\gamma_3}
%DIF >      %     }
%DIF >      %     \le
%DIF >      %     \overline{z}
%DIF >      %     ,
%DIF >      % $
%DIF >      $
%DIF >          \max
%DIF >          \left(
%DIF >              \sqrt{
%DIF >                  \tfrac{\gamma_2}{\gamma_1}
%DIF >              }
%DIF >              \norm{\mv{z}(0)}
%DIF >              ,
%DIF >              \sqrt{
%DIF >                  \tfrac{\gamma_2 c}
%DIF >                  {\gamma_1\gamma_3}
%DIF >              }
%DIF >          \right)
%DIF >          \le
%DIF >          \overline{z}
%DIF >      $
%DIF >      which holds if 
%DIF >      % $
%DIF >      %     \norm{\mv{z}(0)} 
%DIF >      %     \le 
%DIF >      %     \sqrt{
%DIF >      %         \tfrac{\gamma_1}{\gamma_2}\overline{z}^2
%DIF >      %         -
%DIF >      %         \tfrac{c}{\gamma_3}
%DIF >      %     }
%DIF >      % $
%DIF >      $
%DIF >          \norm{\mv{z}(0)}
%DIF >          \le
%DIF >          \overline{z}
%DIF >          \sqrt{
%DIF >              \tfrac{\gamma_1}{\gamma_2}
%DIF >          }
%DIF >      $
%DIF >      .
%DIF >  \end{lem}

\DIFdelbegin \DIFdel{Using $\fe$, the system dynamics \eqref{eq:sys1} can be rewritten as
}%DIFDELCMD < \begin{equation}
%DIFDELCMD <     \rbM \ddtt\fe
%DIFDELCMD <     =
%DIFDELCMD <     -\rbVm \fe
%DIFDELCMD <     -\mm K \fe
%DIFDELCMD <     + \mv f
%DIFDELCMD <     -\rbd + \mysat(\rbu)
%DIFDELCMD <     ,
%DIFDELCMD <     \label{eq:sys:filtered}
%DIFDELCMD < \end{equation}
%DIFDELCMD < %%%
\DIFdel{where $\mm{K}=\mm{K}^\top\in\R^{n\times n}_{>0}$ denotes an arbitrary unknown matrix and $
    \mv f:= \mv f
    \left(
        \q,\dq,\qd,\dqd,\ddqd
    \right)
    =
    \mm K \fe
    +\rbM
    \left(
        -\ddqd+\mm\Lambda\ddtt{\mv e}
    \right)
    +
    \rbVm
    \left(
        -\dqd+\mm\Lambda\mv e
    \right)
    -
    \rbF
    -
    \rbG
    \in\R^n
$ denotes the lumped unknown system.
%DIF <  \color{red}
%DIF <  The unknown matrix $\mm{K}$ is not available to design in practice, 
%DIF <  \color{black}
}\DIFdelend %DIF >  \begin{proof}
%DIF >      % Using Comparison Lemma \cite[pp. 689--660, Lemma 3.4]{Khalil:2002aa}, the proof is straightforward.
%DIF >      Using Comparison Lemma (Lemma~\ref{lem:comp}), the proof is straightforward, and thus omitted for brevity.
%DIF >  \end{proof}

\DIFdelbegin \DIFdel{Consider the Lyapunov function $V_1:= \tfrac{1}{2}\fe^\top \rbM\fe$. 
Invoking Property \ref{prop:skew}, the time derivative of $V_1$ is
}%DIFDELCMD < \begin{equation}
%DIFDELCMD <     \begin{aligned}
%DIFDELCMD <         \ddtt{V}_1
%DIFDELCMD <         =&
%DIFDELCMD <         \fe^\top \rbM \ddtt \fe
%DIFDELCMD <         +
%DIFDELCMD <         \tfrac{1}{2}\fe^\top \ddtt{\rbM}\fe
%DIFDELCMD <         \\
%DIFDELCMD <         =&
%DIFDELCMD <         \fe^\top 
%DIFDELCMD <         \left(
%DIFDELCMD <             -\rbVm \fe -\mm K \fe + \mv f
%DIFDELCMD <             -\rbd + \mysat(\rbu)
%DIFDELCMD <         \right)
%DIFDELCMD <         \\
%DIFDELCMD <         &
%DIFDELCMD <         +
%DIFDELCMD <         \tfrac{1}{2}
%DIFDELCMD <         \fe^\top \ddtt{\rbM}\fe
%DIFDELCMD <         \\
%DIFDELCMD <         =&
%DIFDELCMD <         -
%DIFDELCMD <         \fe^\top \mm K \fe 
%DIFDELCMD <         +
%DIFDELCMD <         \fe^\top 
%DIFDELCMD <         \left(
%DIFDELCMD <             \mv f
%DIFDELCMD <             -\rbd
%DIFDELCMD <             +\mysat(\rbu)
%DIFDELCMD <         \right)
%DIFDELCMD <         \\
%DIFDELCMD <         &
%DIFDELCMD <         +
%DIFDELCMD <         \tfrac{1}{2}
%DIFDELCMD <         \fe^\top(\ddtt{\rbM}- 2\rbVm)\fe
%DIFDELCMD <         \\
%DIFDELCMD <         \le&
%DIFDELCMD <         -\lambda_{\min}(\mm K)\norm{\fe}^2
%DIFDELCMD <         +
%DIFDELCMD <         \overline\tau_d\norm{\fe}
%DIFDELCMD <         +
%DIFDELCMD <         \fe^\top (\mv f+\mysat(\rbu))
%DIFDELCMD <         % \\
%DIFDELCMD <         % \le&
%DIFDELCMD <         % -\lambda_{\min}(\mm K)\norm{\fe}^2
%DIFDELCMD <         % +
%DIFDELCMD <         % \overline\tau_d\norm{\fe}
%DIFDELCMD <         % +
%DIFDELCMD <         % \fe^\top (\mysat(\rbu)-\rbu^*)
%DIFDELCMD <         .
%DIFDELCMD <     \end{aligned}
%DIFDELCMD <     \label{eq:lya1:dot1}
%DIFDELCMD < \end{equation} 
%DIFDELCMD < %%%
\DIFdel{Therefore, invoking Lemma \ref{lem:stable:set}, it can be concluded that the filtered error $\fe$ is uniformly ultimately bounded with $\lim_{t\to\infty}\norm{\fe}\le\tfrac{\overline\tau_d}{\lambda_{\min}(\mm K)}$, provided that the lumped unknown system $\mv{f}$ can be perfectly cancelled by the control input $\rbu$, neglecting the effect of input saturation.
Hence, the ideal control input $\rbu^*$ can be designed as $-\mv f$.
However, $\rbu^*$ is not available in practice, since the system dynamics encapsulated in $\mv f$ are unknown.
}\DIFdelend \DIFaddbegin \begin{figure}[t]
  \centering
  \includegraphics[width=0.99\linewidth]{figs/DNN.drawio.pdf}
  \caption{
      \DIFaddFL{Example of the DNN $\NN(\NNin;\wth)$ with $k=2,l_0=2,l_1=l_2=3$, and $l_3=2$.
  }}
  \label{fig:DNN}
\end{figure}
\DIFaddend 

\DIFdelbegin \DIFdel{To overcome this issue, a DNN is employed to approximate $\rbu^*$.
Let $\NN:=\NN({\q}_n;\wth): \R^{l_0+1}\times\R^{\Xi}\to\R^{n}$ represent the DNN , where ${\q}_n\in\R^{l_0+1}$ is the DNN input vector, and $\wth\in\R^{\Xi}$ is the vector of trainable weights.
The architecture of DNN $\NN({\q}_n;\wth)$ will be presented in Section \ref{sec:sub:NN definition}, later.
}\DIFdelend %DIF >  ————————————————————————————————
%DIF >  DNN Preliminaries
%DIF >  ————————————————————————————————
\DIFaddbegin 

\DIFadd{We review the DNN structure and its approximation capability, which are essential for the subsequent controller design and stability analysis.
}\DIFaddend According to the universal approximation \DIFdelbegin \DIFdel{theorem for DNNs \mbox{%DIFAUXCMD
\cite{Kidger:2020aa}}\hskip0pt%DIFAUXCMD
, $\NN({\q}_n;\wth)$ can approximate a nonlinear function , denoted as $\mv g(\cdot)$, with an ideal weight vector $\idealwth\in\R^\Xi$ over a compact subset $\Omega_{n}\in\R^{l_0+1}$ to within $\epsilon$-accuracy, such that 
$\sup_{{\q}_n\in\Omega_{n}}\norm{ \NN({\q}_n;\idealwth) - \mv g(\cdot)} = \epsilon < \infty$.
The compact set $\Omega_{n}$ is defined as the set of all feasible states $\q,\dq$ and filtered errors $\fe$.
%DIF <  Furthermore, the theorem states that the norm of $\idealwth$ is bounded such that $\norm{\idealwth}\le \overline\wth<\infty$.
In this study, the ideal weight vector $\idealwth$ is assumed to be bounded by Assumption \ref{assum:feasible} and \mbox{%DIFAUXCMD
\cite[Assum. 1]{Lewis:1996aa}}\hskip0pt%DIFAUXCMD
, and considered a locally optimal solution, rather than a global one.
}%DIFDELCMD < 

%DIFDELCMD < %%%
\DIFdel{Accordingly, the ideal control input $\rbu^*$ can be approximated by the DNN with the ideal weight vector $\idealNN:=\NN({\q}_n;\idealwth)$ as follows:
}%DIFDELCMD < \begin{align}
%DIFDELCMD <     \rbu^*=& \idealNN+\mv\epsilon
%DIFDELCMD <     ,
%DIFDELCMD <     \label{eq:control:ideal}
%DIFDELCMD < \end{align}%%%
\DIFdelend \DIFaddbegin \DIFadd{property \mbox{%DIFAUXCMD
\cite{Patil:2022aa}}\hskip0pt%DIFAUXCMD
, the DNNs $\NN: \mathbb{R}^{l_0+1} \times \mathbb{R}^{\Xi} \to \mathbb{R}^{l_{k+1}}$ can be used to approximate any smooth unknown nonlinear function $\mv{g}:\R^{l_0+1}\to\R^{l_{k+1}}$.
Specifically, there exists a DNN such that 
}\begin{equation} \label{eq:dnn:approx:error}
  \sup_{\NNin\in\Omega_{\NNin}}
  \norm{
    \mv{g}(\NNin) - \NN(\NNin;\wth^*)
  }
  =
  \sup_{\NNin\in\Omega_{\NNin}}
  \norm{
    \mv{\epsilon}
  }
  =
  \overline{\epsilon}
  ,
\end{equation}\DIFaddend 
where \DIFdelbegin \DIFdel{$\mv\epsilon\in\R^n$ is the approximation error vector, bounded by $\norm{\mv\epsilon}\le \overline{\epsilon}$ for some $\overline{\epsilon}\in\R_{>0}$.
The ideal control input $\rbu^*$ is estimated online by
}%DIFDELCMD < \begin{align}
%DIFDELCMD <     \rbu =& \estNN,
%DIFDELCMD <     \label{eq:control:est}
%DIFDELCMD < \end{align}
%DIFDELCMD < %%%
\DIFdel{where $\hat\NN:=\NN({\q}_n;\estwth)$, and $\estwth\in\R^\Xi$ is the estimated weight vector for the ideal weight vector $\idealwth$.
}%DIFDELCMD < 

%DIFDELCMD < %%%
\DIFdel{Substituting \eqref{eq:control:ideal} and \eqref{eq:control:est} into \eqref{eq:lya1:dot1}, the time derivative of $V_1$ can be rewritten as
}%DIFDELCMD < \begin{equation}
%DIFDELCMD <     \ddtt{V}_1
%DIFDELCMD <     \le 
%DIFDELCMD <     -\lambda_{\min}(\mm K)\norm{\fe}^2
%DIFDELCMD <     +
%DIFDELCMD <     \overline\tau_d\norm{\fe}
%DIFDELCMD <     +
%DIFDELCMD <     \fe^\top 
%DIFDELCMD <     \left(
%DIFDELCMD <         -\idealNN
%DIFDELCMD <         -\mv\epsilon
%DIFDELCMD <         +\mysat(\hat\NN)
%DIFDELCMD <     \right)
%DIFDELCMD <     .
%DIFDELCMD <     \label{eq:lya1:dot2}
%DIFDELCMD < \end{equation}
%DIFDELCMD < %%%
\DIFdel{Therefore, it is concluded that the filtered error $\fe$ can be stabilized by adapting $\estwth$ to $\idealwth$, }%DIFDELCMD < \ie %%%
\DIFdel{$\hat\NN\to\idealNN$, neglecting the effect of input saturation.
Note that the control input saturation has not yet been considered in the derivation of the ideal control input; this will be addressed in the subsequent section.
}%DIFDELCMD < 

%DIFDELCMD < %%%
\subsection{\DIFdel{Deep Neural Network (DNN) Model}}%DIFAUXCMD
\addtocounter{subsection}{-1}%DIFAUXCMD
%DIFDELCMD < \label{sec:sub:NN definition}
%DIFDELCMD < 

%DIFDELCMD < \begin{figure}[t]
%DIFDELCMD <     \centering
%DIFDELCMD <     \includegraphics[width=0.99\linewidth]{src/figures/DNN.drawio.pdf}
%DIFDELCMD <     %%%
%DIFDELCMD < \caption{%
{%DIFAUXCMD
\DIFdelFL{Architecture of the DNN $\NN({\q}_n;\wth)$ with $k=2,l_0=2,l_1=l_2=3$, and $l_3=2$.
    }}
    %DIFAUXCMD
%DIFDELCMD < \label{fig:DNN}
%DIFDELCMD < \end{figure}
%DIFDELCMD < 

%DIFDELCMD < %%%
\DIFdel{The DNN architecture $\NN({\q}_n;\wth) = \NN_k$ }\DIFdelend \DIFaddbegin \DIFadd{$\NNin \in \mathbb{R}^{l_0+1}$ is the DNN input vector and $\idealwth \in \mathbb{R}^{\Xi}$ denotes the ideal DNN weights that minimize the approximation error $\mv{\epsilon}\in\R^{l_{k+1}}$.
The DNN input vector $\NNin$ should be designed to include all arguments of the unknown function $\mv{g}$.
A compact set $\Omega_{\NNin} \subset \mathbb{R}^{l_0+1}$ denotes the approximation region of the DNN, which ensures the boundedness of $\idealwth$ and $\mv{\epsilon}$, such that $\norm{\idealwth} \le \overline{\theta}^*\in\R_{>0}$ and $\norm{\mv{\epsilon}} \le \overline{\epsilon} \in\R_{>0}$.
The DNN $\NN(\NNin;\wth)=\NN_k$ }\DIFaddend can be recursively \DIFdelbegin \DIFdel{represented, as
follows:
}%DIFDELCMD < \begin{equation}
%DIFDELCMD <     \NN_i :=
%DIFDELCMD <     \begin{cases}
%DIFDELCMD <         \wV_i^\top \act_i(\NN_{i-1}), 
%DIFDELCMD <         &
%DIFDELCMD <         i\in\{1,\dots ,k\},
%DIFDELCMD <         \\
%DIFDELCMD <         \wV_0^\top {\q}_n,
%DIFDELCMD <         &
%DIFDELCMD <         i=0
%DIFDELCMD <         ,
%DIFDELCMD <     \end{cases}
%DIFDELCMD <     \label{eq:DNN:architecture}
%DIFDELCMD < \end{equation}%%%
\DIFdelend \DIFaddbegin \DIFadd{defined as
}\begin{equation} \label{eq:dnn:def}
  \NN_i :=
  \begin{cases}
    \wV_i^\top \act_i(\NN_{i-1}), 
    &
    i\in\{1,\dots ,k\},
    \\
    \wV_0^\top \NNin,
    &
    i=0
    ,
  \end{cases}
\end{equation}\DIFaddend 
where \DIFdelbegin \DIFdel{$\wV_i=[w_{ij}]_{i\in\{1,\cdots,l_i+1\},j\in\{1,\cdots,l_{i+1}\}}\in\R^{(l_i+1)\times l_{i+1}}$ }\DIFdelend \DIFaddbegin \DIFadd{$\wV_i\in\R^{(l_i+1)\times l_{i+1}}$ }\DIFaddend and $\act_i: \R^{l_i}\to\R^{l_i+1}$ represent the weight matrix and the activation function of the $i$th layer \DIFdelbegin \DIFdel{, respectively}\DIFdelend \DIFaddbegin \DIFadd{for all $i\in\mathcal{K}$, respectively, and $\mathcal{K}:=\{0,1,\cdots,k\}$ denotes the index set of layers}\DIFaddend .
The DNN \DIFdelbegin \DIFdel{$\NN({\q}_n;\wth)$ }\DIFdelend \DIFaddbegin \DIFadd{$\NN(\NNin;\wth)$ }\DIFaddend has $k+1$ layers, where the input \DIFdelbegin \DIFdel{layer is indexed by }\DIFdelend \DIFaddbegin \DIFadd{and output layers are indexed by }\DIFaddend $i=0$ \DIFdelbegin \DIFdel{, and the output layer is indexed by }\DIFdelend \DIFaddbegin \DIFadd{and }\DIFaddend $i=k$\DIFdelbegin \DIFdel{.
For instance, }\DIFdelend \DIFaddbegin \DIFadd{, respectively, as illustrated by the example }\DIFaddend in Fig.~\ref{fig:DNN}\DIFdelbegin \DIFdel{, the DNN $\NN({\q}_n;\wth)$ with $k=2,l_0=2,l_1=l_2=3$, and $l_3=2$ is illustrated.
Notice that the output size of $\NN(\cdot)$ is the same as that of the control input $\rbu$, }%DIFDELCMD < \ie %%%
\DIFdel{$l_{k+1}=n$}\DIFdelend .
\DIFdelbegin %DIFDELCMD < 

%DIFDELCMD < %%%
\DIFdel{The DNN input vector ${\q}_n$ is defined as ${\q}_n=(\q^\top,\dq^\top,\fe^\top,1)^\top\in\R^{3n+1}$ to include the current state of the system, the filtered error, and the augmentation term $1$ to account for the bias term in the input weight matrix $\wV_0$.
The }\DIFdelend \DIFaddbegin \DIFadd{The }\DIFaddend activation function is defined as \DIFdelbegin \DIFdel{$\act_i(\mv x)=\left(\sigma(x_1),\sigma(x_2),\cdots, \sigma(x_{l_{i}}), 1\right)^\top,\ \forall\mv x\in\R^{l_i}$, where $\sigma: \R\to\R$ is a nonlinear function, and the augmentation of $1$ is used to account for bias terms in the weight matrices, similar to the input vector ${\q}_n$.
The weights $\wV_i,\ \forall i\in\{0,\cdots,k\}$, are initialized randomly using a uniform distribution at the beginning of the control process.
}%DIFDELCMD < 

%DIFDELCMD < %%%
%DIF <  \MSRY{생략 가능할 것 같습니다.}
%DIF <  \ctxt{red}
%DIF <  {
    \DIFdel{In selecting the nonlinear function $\sigma(\cdot)$, the ReLU function \mbox{%DIFAUXCMD
\cite{Maas:2013aa} }\hskip0pt%DIFAUXCMD
is one of the most widely used choices for large DNNs, since it effectively mitigates the gradient vanishing problem during error backpropagation. 
    However, for control applications, where relatively shallow DNNs are typically sufficient and the gradient vanishing issue is less severe, the sigmoid function or the hyperbolic tangent function is commonly used as 
the nonlinear function $\sigma(\cdot)$. 
    These functions simplify stability analysis, since they are continuously differentiable andtheir outputs and gradients are bounded. 
%DIF <  }
In this study, the hyperbolic tangent function $\tanh(\cdot)$ was selected as the nonlinear function, }%DIFDELCMD < \ie %%%
\DIFdel{$\sigma(\cdot) = \tanh(\cdot)$, which provides desirable boundedness with $\norm{\sigma(x)}<1$ and $0<\norm{\ddtfrac{\sigma(x)}{x}}\le 1$.
}%DIFDELCMD < 

%DIFDELCMD < %%%
\DIFdel{For simplicity, each layer's weights are vectorized as $\wth_i:=\myvec(\wV_i)\in\R^{\Xi_i},\ \forall i\in\{0,\cdots,k\}$, where $\Xi_i:= (l_i+1)l_{i+1}$ is the number of weights in the $i$th layer . 
The total weight vector $\wth\in\R^{\Xi}$ is defined by augmenting $\wth_i, \forall i\in \{0,\cdots,k\}$, }\DIFdelend \DIFaddbegin \DIFadd{$\act_i(\mv{z}) = \left(\tanh(z_1),\cdots,\tanh(z_{l_i}),1\right)^\top,$ for some $\mv{z}\in\R^{l_i}$.
We simplify the notations by vectorizing all weight matrices }\DIFaddend as 
\DIFdelbegin %DIFDELCMD < \begin{equation}
%DIFDELCMD <     \wth := 
%DIFDELCMD <     \begin{pmatrix}
%DIFDELCMD <         \wth_k\\
%DIFDELCMD <         \wth_{k-1}\\
%DIFDELCMD <         \vdots\\
%DIFDELCMD <         \wth_0
%DIFDELCMD <     \end{pmatrix}
%DIFDELCMD <     =
%DIFDELCMD <     \begin{pmatrix}
%DIFDELCMD <         \myvec(\wV_k)\\
%DIFDELCMD <         \myvec(\wV_{k-1})\\
%DIFDELCMD <         \vdots\\
%DIFDELCMD <         \myvec(\wV_0)
%DIFDELCMD <     \end{pmatrix},
%DIFDELCMD < \end{equation}
%DIFDELCMD < %%%
\DIFdelend \DIFaddbegin \begin{equation*}
  \wth_i 
  :=
  \myvec(\wV_i)
  \in\R^{(l_i+1)l_{i+1}}
  ,
  \qquad
  \forall 
  i\in\mathcal{K}
  ,
\end{equation*}
\DIFadd{and, finally, stacking them as $\wth := \left[\wth_i\right]_{i\in\mathcal{K}}\in\R^\Xi$, }\DIFaddend where \DIFdelbegin \DIFdel{$\Xi:={\sum_{i=0}^{k} \Xi_i}$ represents the total number of weights in the DNN}\DIFdelend \DIFaddbegin \DIFadd{$\Xi := \sum_{i\in\mathcal{K}} (l_i+1)l_{i+1}$.
Additionally, like $\overline{\theta}^*$, there exists a maximum norm bound for the weights of each layer $\overline{\theta}^*_i\in\R_{>0}$, such that $\norm{\idealwth_i}\le\overline{\theta}^*_i$, $\forall i\in\mathcal{K}$}\DIFaddend .

The gradient of \DIFdelbegin \DIFdel{$ \NN({\q}_n;\wth)$ }\DIFdelend \DIFaddbegin \DIFadd{$ \NN(\NNin;\wth)$ }\DIFaddend with respect to $\wth$ can be obtained using Proposition\DIFaddbegin \DIFadd{~}\DIFaddend \ref{propsit:kron} and the chain rule as follows:
\DIFdelbegin %DIFDELCMD < \begin{equation}
%DIFDELCMD <     \pptfrac{\NN}{\wth}=
%DIFDELCMD <     \begin{bmatrix}
%DIFDELCMD <         \pptfrac{\NN}{\wth_k}&
%DIFDELCMD <         \pptfrac{\NN}{\wth_{k-1}}&
%DIFDELCMD <     \cdots &
%DIFDELCMD <         \pptfrac{\NN}{\wth_0}
%DIFDELCMD <     \end{bmatrix}
%DIFDELCMD <     \in\R^{n \times \Xi}
%DIFDELCMD <     ,
%DIFDELCMD <     \label{eq:gradient:phi}
%DIFDELCMD < \end{equation}%%%
\DIFdelend \DIFaddbegin \begin{equation}
  \ppfrac{\NN}{\wth}=
  \begin{bmatrix}
    \ppdfrac{\NN}{\wth_k}&
    \ppdfrac{\NN}{\wth_{k-1}}&
\cdots &
    \ppdfrac{\NN}{\wth_0}
  \end{bmatrix}
  \in\R^{l_{k+1} \times \Xi}
  ,
  \label{eq:gradient:phi}
\end{equation}\DIFaddend 
where
\DIFdelbegin %DIFDELCMD < \begin{equation}
%DIFDELCMD <     \pptfrac{\NN}{\wth_i} = 
%DIFDELCMD <     \begin{cases}
%DIFDELCMD <         (
%DIFDELCMD <             \mm I_{l_{k+1}}
%DIFDELCMD <             \otimes 
%DIFDELCMD <             \act_{k}^\top  
%DIFDELCMD <         )
%DIFDELCMD <         , 
%DIFDELCMD <         &
%DIFDELCMD <         i=k
%DIFDELCMD <         ,
%DIFDELCMD <         \\
%DIFDELCMD <         \wV_k^\top\act_{k}' 
%DIFDELCMD <         (
%DIFDELCMD <             \mm I_{l_{k}}
%DIFDELCMD <             \otimes  
%DIFDELCMD <             \act_{k-1}^\top  
%DIFDELCMD <         )
%DIFDELCMD <         , 
%DIFDELCMD <         & 
%DIFDELCMD <         i=k-1
%DIFDELCMD <         ,
%DIFDELCMD <         \\
%DIFDELCMD <         \vdots 
%DIFDELCMD <         &
%DIFDELCMD <         \vdots 
%DIFDELCMD <         \\
%DIFDELCMD <         \wV_k^\top\act'_{k} 
%DIFDELCMD <         \cdots 
%DIFDELCMD <         \wV_1^\top\act_1' 
%DIFDELCMD <         (
%DIFDELCMD <             \mm I_{l_1}
%DIFDELCMD <             \otimes 
%DIFDELCMD <             {\q}_n^\top  
%DIFDELCMD <         )
%DIFDELCMD <         , 
%DIFDELCMD <         &
%DIFDELCMD <         i = 0
%DIFDELCMD <         ,
%DIFDELCMD <     \end{cases}
%DIFDELCMD <     \label{eq:NN:grad}
%DIFDELCMD < \end{equation}%%%
\DIFdelend \DIFaddbegin \begin{equation}
  \ppfrac{\NN}{\wth_i} = 
  \begin{cases}
    % (
      \mm I_{l_{k+1}}
      \otimes 
      \act_{k}^\top  
    % )
    , 
    &
    i=k
    ,
    \\
    \wV_k^\top\act_{k}' 
    (
      \mm I_{l_{k}}
      \otimes  
      \act_{k-1}^\top  
    )
    , 
    & 
    i=k-1
    ,
    \\
    \vdots 
    &
    \vdots 
    \\
    \wV_k^\top\act'_{k} 
    \cdots 
    \wV_1^\top\act_1' 
    (
      \mm I_{l_1}
      \otimes 
      \NNin^\top  
    )
    , 
    &
    i = 0
    ,
  \end{cases}
  \label{eq:NN:grad}
\end{equation}\DIFaddend 
and $\act_i:= \act_i(\NN_{i-1})$ and $\act_i':= \ddtfrac{\act_i}{\NN_{i-1}}$\DIFaddbegin \DIFadd{, for all $i\in\mathcal{K}\setminus\{0\}$}\DIFaddend .

\DIFdelbegin \DIFdel{In the following sections, let $\idealNN_i$ represent the output of the $i$th layer with the ideal weight vector $\idealwth$.
Define $\idealact_i=\act_i(\idealNN_{i-1})$ and ${\idealact}'_i= \ddtfrac{\idealact_i}{\idealNN_{i-1}}$.
Similarly, let $\estNN_i$ denote the output of the $i$th layer with the estimated weight vector $\estwth$, }\DIFdelend \DIFaddbegin \subsection{\DIFadd{Controller Design}} \label{sec:main:result:sub:ctrl:design}

\DIFadd{We define the tracking error as $\mv{e} = \mv{x} - \mv{x}_d$, and derive the error dynamics from \eqref{eq:sys} as
}\begin{equation} \label{eq:error:dyn}
    \ddtt\mv{e}
    =
    \mv{f}(\mv{x}) + \mm{B}(\mv{x})\mysat(\mv{u}) - \ddtt\mv{x}_d
    .
\end{equation}
\DIFadd{Let us consider the Lyapunov function $V_c:=\tfrac{1}{2}\mv{e}^\top\mm{B}^{-1}(\mv{x})\mv{e}$.
Invoking Assumptions~\ref{assum:B:bound:grad}, the time derivative of $V_c$ can be bounded as follows:
}\begin{align} \label{eq:Vc:dot:1} \nonumber
  \ddtt V_c
  =&
  \tfrac{1}{2}\mv{e}^\top
  \left(
    \ddtt\mm{B}^{-1}
  \right)
  \mv{e}
  +
  \mv{e}^\top\mm{B}^{-1}
  \left(
    \mv{f} + \mm{B}\mysat(\mv{u}) 
    - \ddtt\mv{x}_d
  \right)
  \\
  \le&
  \beta(\mv{x})\norm{\mv{e}}^2
  \!+\!
  \mv{e}^\top\mv{u}
  \!+\!
  \mv{e}^\top\mm{B}^{-1}
  \left(
    \mv{f} \!-\! \ddtt\mv{x}_d
  \right)
  \!+\!
  \mv{e}^\top
  \Delta\mv{u}
  ,
\end{align}
\DIFadd{where the arguments of $\mv{f}(\mv{x})$ }\DIFaddend and \DIFdelbegin \DIFdel{define $\estact_i=\act_i(\estNN_{i-1})$ and $\estact_i'= \ddtfrac{\estact_i}{\estNN_{i-1}}$}\DIFdelend \DIFaddbegin \DIFadd{$\mm{B}(\mv{x})$ are omitted for brevity and $\Delta\mv{u} := \mysat(\mv{u}) - \mv{u}$}\DIFaddend .

%DIF <   SECTION ADAPTATION LAW DERIVATION =======================
\DIFdelbegin \section{\DIFdel{Weight Adaptation Laws}}%DIFAUXCMD
\addtocounter{section}{-1}%DIFAUXCMD
%DIFDELCMD < \label{sec:adap_laws}
%DIFDELCMD < %%%
\DIFdelend \DIFaddbegin \DIFadd{A desired control law $\mv{u}_d$ such that the tracking error $\mv{e}$ is stable in the absence of input saturation is selected as 
}\begin{equation} \label{eq:ctrl:ideal}
  \mv{u}_d
  :=
  -
  \mm{P}
  \mv{e}
  - 
  \beta(\mv{x})\mv{e}
  +
  \mm{B}^{-1}
  \left(
    - \mv{f} 
    + \ddtt\mv{x}_d
  \right)
  % +
  % h(\mv{x},\mv{x}_d)
  ,
\end{equation}
\DIFadd{where $\underline{p}\mm{I}_n \preceq \mm{P}=\mm{P}^\top \preceq \overline{p}\mm{I}_n$ for known positive constants $\underline{p},\overline{p}\in\R_{>0}$.
%DIF >  The function $h:\R^n\times\R^n\to\R^n$ is correction term to satisfy the ideal control law $\mv{u}^*$ within the input constraints, \ie $\mv{u}^*\in\Omega_u$.
}\DIFaddend 

%DIF <  \subsection{Adaptation Law using Lagrangian Function}
\DIFdelbegin \subsection{\DIFdel{Weight Optimizer Design}}%DIFAUXCMD
\addtocounter{subsection}{-1}%DIFAUXCMD
%DIFDELCMD < \label{sec:sub:weight optimizer}
%DIFDELCMD < 

%DIFDELCMD < %%%
\DIFdel{The control objective can be represented as follows:
}%DIFDELCMD < \begin{equation}
%DIFDELCMD <     \begin{matrix}
%DIFDELCMD <         \min_{\estwth} \ J(\fe;\estwth):= 
%DIFDELCMD <         \tfrac{1}{2} \fe^\top \fe
%DIFDELCMD <         ,
%DIFDELCMD <         \vspace{.5em}
%DIFDELCMD <         \\
%DIFDELCMD <         \begin{aligned}
%DIFDELCMD <         \text{subject to }&c_{j}(\estwth) 
%DIFDELCMD <         \le0, \quad \forall j\in\mathcal{I},
%DIFDELCMD <         \end{aligned}
%DIFDELCMD <     \end{matrix}
%DIFDELCMD <     \label{eq:opt:problem1}
%DIFDELCMD < \end{equation}%%%
\DIFdelend \DIFaddbegin \DIFadd{By substituting $\mv{u}_d$ in \eqref{eq:ctrl:ideal} into \eqref{eq:Vc:dot:1}, we have
}\begin{align} \label{eq:Vc:dot:2} 
  \ddtt V_c
  \le &
  -\mv{e}^\top
  \mm{P}
  \mv{e}
  +
  \mv{e}^\top
  \left(
    \mv{u} - \mv{u}_d
  \right)
  +
  \mv{e}^\top
  \Delta\mv{u}
  \\ \nonumber
  \le & 
  -
  \underline{p}\norm{\mv{e}}^2
  +
  \mv{e}^\top
  \left(
    \mv{u} - \mv{u}_d       
  \right)
  +
  \mv{e}^\top
  % \left(
    \Delta\mv{u}
%     +
%     h(\mv{x},\mv{x}_d)
  % \right)
  .
\end{align}\DIFaddend 
\DIFaddbegin \DIFadd{Notably, in the ideal case, }\ie \DIFadd{$\mv{u} = \mv{u}_d$ and $\Delta\mv{u} = \mv{0}$, the tracking error $\mv{e}$ is asymptotically exponentially stable with decay rate $\underline{p}\overline{b}$, since 
$
  \tfrac{1}{\overline{b}}
  \ddtt
  \left(
    \norm{\mv{e}}^2
  \right)
  \le
  \ddtt V_c 
  \le 
  -\underline{p}\norm{\mv{e}}^2
$.
}

\DIFadd{Since the ideal control law \eqref{eq:ctrl:ideal} contains unknown functions $\beta(\mv{x})$, $\mv{f}$ and $\mm{B}$, we employ the DNN $\NN(\NNin;\wth)$ in \eqref{eq:dnn:def} to approximate the ideal control law as follows:
}\begin{equation*}
% \begin{aligned}
  \NN(\NNin;\idealwth)
  +
  \mv{\epsilon}
  =
  -\mv{e}
  -
  \mm{P}^{-1}
  \left(
  \beta(\mv{x})\mv{e}
  +
  \mm{B}^{-1}
  (
    \mv{f}
    -
    \ddtt \mv{x}_d
  )
  \right)
  .
% \end{aligned}
\end{equation*}
%DIF >  where $\mv{\epsilon}\in\R^n$ denotes an approximation error vector, such that 
%DIF >  $
%DIF >  \sup_{\NNin\in\Omega_{\NNin}}   
%DIF >  \left( 
%DIF >    \norm{\mv{\epsilon}}
%DIF >  \right)
%DIF >  =
%DIF >  \overline{\epsilon}
%DIF >  .
%DIF >  $
\DIFadd{The DNN input vector $\NNin$ can be defined in various ways, as long as it includes the arguments of the unknown functions.
In this formulation we define the DNN input vector as
$
  \NNin 
  = 
  \left(
    \mv{x}_d^\top, 
    \ddtt{\mv{x}}_d^\top, 
    \mv{x}^\top, 
    1 
  \right)^\top
  \in\R^{3n+1}
$.
Then, we have 
}\begin{equation} \label{eq:ctrl:NN:ideal}
  \mv{u}_d
  =
  \mm{P}
  \big(
    \NN(\NNin;\idealwth)
    +
    \mv{\epsilon}
  \big)
  .
\end{equation}
\DIFadd{The actual control law is designed as 
}\begin{equation} \label{eq:ctrl:NN:actual}
  \mv{u}
  :=
  \mm{P}
  \NN(\NNin;\estwth)
  ,
\end{equation}
\DIFaddend where \DIFdelbegin \DIFdel{$J(\fe;\estwth)$ }\DIFdelend \DIFaddbegin \DIFadd{$\estwth\in\R^{\Xi}$ denotes the estimated DNN weights.
}

\DIFadd{By substituting \eqref{eq:ctrl:NN:ideal} and \eqref{eq:ctrl:NN:actual} into \eqref{eq:Vc:dot:2}, we have
}\begin{equation} \label{eq:Vc:dot:3}
  \ddtt V_c
  \le
  - \underline{p}
  \norm{\mv{e}}^2
  +
  \mv{e}^\top\mm{P}
  \left(
    \estNN - \idealNN
    - \mv{\epsilon}
  \right)
  +
  \mv{e}^\top
  \Delta\mv{u}
  ,
\end{equation}
\DIFadd{where $\estNN\!:=\!\NN(\NNin;\estwth)$ and $\idealNN\!:=\!\NN(\NNin;\wth^*)$ are abbreviated.
}

\subsection{\DIFadd{Adaptation Law Derivation}} \label{sec:main:result:sub:adap:laws}

%DIF >  The control objective which minimizes the tracking error $\mv{e}$ while satisfying both weight boundedness and input constraints can be represented as follows:
\DIFadd{The control objectives defined in Section~\ref{sec:problem:formulation} can be formulated as a constrained optimization problem, as follows:
}\begin{subequations} \label{eq:opt:prob} 
  \begin{align} 
    \min_{\estwth} \quad 
    & 
    J(\mv{e};\estwth)
    = 
    \tfrac{1}{2} \mv{e}^\top \mv{e} 
    &
    \label{eq:opt:prob:obj} 
    \\ 
    \text{subject to} \quad  
    & 
    c^{\theta}_i(\estwth_i)=\tfrac{1}{2}
    \left(
      \norm{\estwth_i}^2
      -
      \overline{\theta}_i^2 
    \right)
    \le 0, 
    \!
    &
    \!
    \forall i\in\mathcal{K} 
    ,
    \label{eq:opt:prob:const1} 
    \\ 
    & 
    c^{u}_j(\estwth)
    % c_j^u
    % \left(
    %     \mm{P}\NN(\NNin;\estwth)
    % \right) 
    \le 0, 
    &
    \forall j\in\mathcal{I} 
    \label{eq:opt:prob:const2} 
    ,
  \end{align}%DIF > 
\end{subequations}
\DIFadd{where $J(\mv{e};\estwth)$ }\DIFaddend is the objective function.
\DIFdelbegin \DIFdel{Inequality constraints $c_j:=c_j(\estwth),\ \forall j\in\mathcal{I}$, are imposed during the weight adaptation process, where $\mathcal I$ denotes the set of the imposed inequality constraints }\DIFdelend \DIFaddbegin \DIFadd{In the optimization problem, inequality constraints of weight boundedness $c^{\theta}_i(\estwth_i),\ \forall i\in\mathcal{K},$ and input constraints $c^{u}_j(\estwth),\ \forall j\in\mathcal{I},$ are imposed}\DIFaddend .
%DIF <  Here, $\fe$ is considered a pre-defined data or static parameter for this optimization problem, so that the adaptation process does not consider the dynamics of $\fe$.
\DIFaddbegin \MSRY{inequality constraints of input constraints인데 괜찮을까요? 제가 의도한건 inequality constraints of weight boundedness and input saturation이긴했습니다.}

\DIFadd{Specifically, the input constraints used in \eqref{eq:opt:prob:const2} are the same convex constraint functions that characterize the admissible input set in Assumption~\ref{assum:sat}.
With a slight abuse of notation, the composed input constraints are denoted by
$c_j^u(\estwth):=c_j^u(\mv u(\estwth))=c_j^u(\mm P\NN(\NNin;\estwth)),\ \forall j\in\mathcal I$, to explicitly represent the dependence of the input constraints on the DNN weights $\estwth$.
Additionally, all active constraints are assumed to satisfy the Linear Independence Constraint Qualification (LICQ) to avoid degeneracy \mbox{%DIFAUXCMD
\cite[pp. 465--467]{Nocedal:2006aa}}\hskip0pt%DIFAUXCMD
.
}

\DIFaddend The Lagrangian function \DIFdelbegin \DIFdel{$L(\cdot)$ }\DIFdelend \DIFaddbegin \DIFadd{$L:=L(\mv{e},\estwth,[\sigma_i]_{i\in\mathcal{K}},[\lambda_j]_{j\in\mathcal I})$ }\DIFaddend is defined as
\DIFdelbegin %DIFDELCMD < \begin{equation}
%DIFDELCMD <     L(
%DIFDELCMD <         \fe,\estwth,[\lambda_j]_{j\in\mathcal I}
%DIFDELCMD <     ) 
%DIFDELCMD <     = 
%DIFDELCMD <     J(\fe;\estwth) 
%DIFDELCMD <     + 
%DIFDELCMD <     \textstyle\sum_{j\in\mathcal I}
%DIFDELCMD <     \lambda_{j}
%DIFDELCMD <     c_{j}(\estwth)
%DIFDELCMD <     ,
%DIFDELCMD <     \label{eq:lagrangian}
%DIFDELCMD < \end{equation}%%%
\DIFdelend \DIFaddbegin \begin{equation}
  L
  = 
  J(\mv{e};\estwth) 
  + 
  \textstyle\sum_{i\in\mathcal{K}}
  \sigma_{i}
  c^{\theta}_{i}(\estwth_i)
  +
  \textstyle\sum_{j\in\mathcal I}
  \lambda_{j}
  c^{u}_{j}(\estwth)
  ,
  \label{eq:lagrangian}
\end{equation}\DIFaddend 
where \DIFdelbegin \DIFdel{$\lambda_j,\ \forall j\in\mathcal{I}$, denotes the Lagrange multiplier for each constraint.
}\DIFdelend \DIFaddbegin \DIFadd{$[\sigma_i]_{i\in\mathcal{K}}$ and $[\lambda_j]_{j\in\mathcal{I}}$ denote the Lagrange multipliers for the weight constraints \eqref{eq:opt:prob:const1} and input constraints \eqref{eq:opt:prob:const2}, respectively.
Then, the primal-dual problem can be formulated as follows:
}\begin{equation} \label{eq:opt:prob:saddle}
  \min_{\estwth}
  \max_{
    [\sigma_i]_{i\in\mathcal{K}}
    ,
    [\lambda_j]_{j\in\mathcal{I}}
  } 
  L(\mv{e},\estwth,[\sigma_i]_{i\in\mathcal{K}},[\lambda_j]_{j\in\mathcal I})
  .
\end{equation}
\DIFaddend 

The adaptation laws for \DIFaddbegin \DIFadd{the adaptive variables, }\DIFaddend $\estwth$\DIFdelbegin \DIFdel{and $[\lambda_j]_{j\in\mathcal I}$ }\DIFdelend \DIFaddbegin \DIFadd{, $[\sigma_i]_{i\in\mathcal{K}}$, and $[\lambda_j]_{j\in\mathcal{I}}$, }\DIFaddend are derived to solve \DIFdelbegin \DIFdel{the dual problem of \eqref{eq:opt:problem1}}\DIFdelend \DIFaddbegin \DIFadd{\eqref{eq:opt:prob:saddle} by using a primal-dual gradient update law}\DIFaddend , \ie 
\DIFdelbegin \DIFdel{$\min_{\estwth} \max_{[\lambda_j]_{j\in\mathcal I}}L(\fe,\estwth,[\lambda_j]_{j\in\mathcal I})$, }\DIFdelend \DIFaddbegin \DIFadd{$
  \ddtt\estwth 
  = 
  -\alpha \pptfrac{L}{\estwth}
$, $
  \ddtt\sigma_i 
  = 
  \myproj_{\sigma_i\ge0}
  \left(
      \beta_i^\theta 
      \pptfrac{L}{\sigma_i}
  \right)
$, and $
  \ddtt\lambda_j 
  = 
  \myproj_{\lambda_j\ge0}
  \left(
    \beta_j^u 
    \pptfrac{L}{\lambda_j}
  \right)
$, }\DIFaddend as follows:
\begin{subequations} \DIFdelbegin %DIFDELCMD < \begin{align}
%DIFDELCMD <             \ddtt{\estwth}
%DIFDELCMD <             &
%DIFDELCMD <             =
%DIFDELCMD <             -\alpha \pptfrac{L}{\estwth}
%DIFDELCMD <             =
%DIFDELCMD <             -\alpha 
%DIFDELCMD <             \left(
%DIFDELCMD <                 \pptfrac{J}{\estwth}
%DIFDELCMD <                 +
%DIFDELCMD <                 \textstyle\sum_{j\in\mathcal{I}}
%DIFDELCMD <                 \lambda_j 
%DIFDELCMD <                 \pptfrac{c_j}{\estwth}
%DIFDELCMD <             \right),
%DIFDELCMD <         \label{eq:adap:th}
%DIFDELCMD <             \\
%DIFDELCMD <             \ddtt\lambda_j
%DIFDELCMD <             & 
%DIFDELCMD <             = 
%DIFDELCMD <             \beta_j\pptfrac{L}{\lambda_j} 
%DIFDELCMD <             = 
%DIFDELCMD <             \beta_j c_j ,
%DIFDELCMD <             \quad \forall j\in\mathcal I,
%DIFDELCMD <         \label{eq:adap:lbd}
%DIFDELCMD <             \\
%DIFDELCMD <             \lambda_j & \leftarrow \max(\lambda_j,0) ,
%DIFDELCMD <         \label{eq:adap:lbd:max}
%DIFDELCMD <     \end{align}%%%
\DIFdelend \DIFaddbegin \label{eq:adap}
  \begin{align}
    \ddtt{\estwth}
    &
    =
    -\alpha 
    \left(
      \pptfrac{J}{\estwth}
      +
      \textstyle\sum_{i\in\mathcal{K}}
      \sigma_i 
      \estwth_i
      +
      \textstyle\sum_{j\in\mathcal{I}}
      \lambda_j 
      \pptfrac{c^{u}_{j}}{\estwth}
    \right),
    \label{eq:adap:th}
    \\
    \ddtt\sigma_i
    & 
    = 
    \myproj_{\sigma_i\ge0}
      \left(
      \beta_i^{\theta}
      c^{\theta}_{i} 
    \right)
    ,
    \qquad \forall i\in\mathcal{K},
    \label{eq:adap:sig}
    \\
    \ddtt\lambda_j
    & 
    = 
    \myproj_{\lambda_j\ge0}
      \left(
      \beta_j^{u}
      c^{u}_{j} 
    \right)
    ,
    \qquad \forall j\in\mathcal{I},
    \label{eq:adap:lbd}
  \end{align}\DIFaddend \DIFdelbegin %DIFDELCMD < \label{eq:adaptation law}%%%
%DIF <              %%
\DIFdelend %DIF > 
\end{subequations}
where $\alpha\in\R_{>0}$ denotes the adaptation gain (learning rate) and \DIFdelbegin \DIFdel{$\beta_j\in\R_{>0}$ denotes the update rate }\DIFdelend \DIFaddbegin \DIFadd{$\beta_i^{\theta},\beta_j^{u}\in\R_{>0},\ \forall i\in\mathcal{K},j\in\mathcal{I}$, denote the update rates }\DIFaddend of the Lagrange multipliers\DIFdelbegin \DIFdel{in $\mathcal I$, and the arguments of $L(\cdot)$ and $J(\cdot)$ are suppressed for brevity}\DIFdelend . 
The Lagrange multipliers associated with inequality constraints are \DIFaddbegin \DIFadd{kept }\DIFaddend non-negative, \ie \DIFdelbegin \DIFdel{$\lambda_j\ge 0$}\DIFdelend \DIFaddbegin \DIFadd{$\sigma_i,\lambda_j\ge 0$}\DIFaddend , due to \DIFdelbegin \DIFdel{\eqref{eq:adap:lbd:max}.
}\DIFdelend \DIFaddbegin \DIFadd{the projection operator $\myproj_{(\cdot)\ge0}(\cdot)$ in \eqref{eq:adap:sig} and \eqref{eq:adap:lbd}, which is presented in \mbox{%DIFAUXCMD
\cite[Appendix E, eq. (E.4)]{Krstic:1995aa}}\hskip0pt%DIFAUXCMD
.
}

\DIFaddend When a constraint \DIFdelbegin \DIFdel{$c_j$ becomes active}\DIFdelend \DIFaddbegin \DIFadd{becomes violated}\DIFaddend , \ie \DIFdelbegin \DIFdel{$c_j > 0$}\DIFdelend \DIFaddbegin \DIFadd{positive constraint violation $c^{\theta}_{i},c^{u}_{j}\ge0$}\DIFaddend , the corresponding Lagrange multiplier \DIFdelbegin \DIFdel{$\lambda_j$ increases , see }\DIFdelend \DIFaddbegin \DIFadd{increases according to \eqref{eq:adap:sig} and }\DIFaddend \eqref{eq:adap:lbd}, to address the violation \DIFdelbegin \DIFdel{, see \eqref{eq:adap:th}.
Once the violation is resolved and the constraint is no longer active, }%DIFDELCMD < \ie %%%
\DIFdel{$c_j \le 0$, }\DIFdelend \DIFaddbegin \DIFadd{by updating weight $\estwth$ through \eqref{eq:adap:th}.
Ultimately, }\DIFaddend the Lagrange multiplier \DIFdelbegin \DIFdel{$\lambda_j$ decreases gradually until it returns to zero , see \eqref{eq:adap:lbd:max}.
%DIF <  \MSRY{생략 가능할 것 같습니다.}
%DIF <  \ctxt{red}
%DIF <  {
%DIF <      Note that this adaption law is similar to the ALM in \cite{Nocedal:2006aa}, where the adaptation law for Lagrange multipliers is given by $\lambda_j\leftarrow \max\left(\lambda_j+\tfrac{c_j}{\mu},0\right)$, with $\mu\in\R_{>0}$ being the penalty parameter. 
%DIF <  }
}\DIFdelend \DIFaddbegin \DIFadd{exhibits two distinct behaviors depending on the constraint status: it either decreases to zero (}\eg \DIFadd{via a projection operator) when the constraint becomes inactive, $c^{\theta}_i,c^{u}_j < 0$, or converges to a positive value when the constraint remains binding, $c^{\theta}_i,c^{u}_j = 0$, to enforce the limit.
}\DIFaddend 

\DIFdelbegin \DIFdel{At steady state, where $\ddtt{\estwth}=0$ and $\ddtt\lambda_j=0$, the KKT conditions are satisfied, }%DIFDELCMD < \ie %%%
\DIFdel{$\pptfrac{L}{\estwth}=0$, $c_j \le 0$, $\lambda_j \ge 0$, and $\lambda_j c_j=0$.
In other words, the proposed weight optimizer updates $\estwth$ and $\lambda_j$ in a way that satisfies the KKT conditions.
  These conditions represent the first-order necessary conditions for optimality, guiding the updates toward candidates for a locally optimal point.
%DIF <  The KKT condition is satisfied, when $\estwth$ converges to $\idealwth$, since $\partial L/\partial \estwth=0$ (i.e. the first-order necessary condition of the optimality is satisfied.).
}\DIFdelend \DIFaddbegin \begin{remark}[Comparison with Existing Methods for Weight Boundedness] \label{rem:weight:bound}
  \DIFadd{Unlike projection-based methods that directly restrict the adaptive weights to a prescribed set, or $\sigma$-modification, or $\epsilon$-modification that continuously bias them toward the origin \mbox{%DIFAUXCMD
\cite{Ryu:2025ab}}\hskip0pt%DIFAUXCMD
, CONAC handles weight boundedness through inequality constraints within the primal-dual adaptation framework like the input constraints.
  The corresponding constraint term is only activated when the prescribed bound is violated, thereby softly driving the weights toward the admissible set without direct clipping or persistent origin-directed bias.
}\end{remark}
\DIFaddend 

\DIFdelbegin \subsection{\DIFdel{Approximation of the Gradient of Objective Function}}
%DIFAUXCMD
\addtocounter{subsection}{-1}%DIFAUXCMD
%DIFDELCMD < 

%DIFDELCMD < %%%
\DIFdel{The adaptation law for }\DIFdelend \DIFaddbegin \DIFadd{Meanwhile, the adaptation law of }\DIFaddend $\estwth$ defined in \eqref{eq:adap:th} involves the partial derivative of the \DIFdelbegin \DIFdel{filtered }\DIFdelend error with respect to \DIFdelbegin \DIFdel{control input $\pptfrac{\fe}{\rbu}$}\DIFdelend \DIFaddbegin \DIFadd{the control input $\pptfrac{\mv{e}}{\mv{u}}$}\DIFaddend , \ie 
\DIFdelbegin \DIFdel{$
    \pptfrac{J}{\estwth}
    =
    \left(
        (\pptfrac{\fe}{\rbu})
        (\pptfrac{\rbu}{\estwth})
    \right)
    ^\top \fe
    =
    \left(
        (\pptfrac{\fe}{\rbu})
        (\pptfrac{\estNN}{\estwth})
    \right)
    ^\top \fe
$}\DIFdelend \DIFaddbegin \DIFadd{$
    \pptfrac{J}{\estwth}
    =
    \big(
        (\pptfrac{\mv{e}}{\mv{u}})
        (\pptfrac{\mv{u}}{\estwth})
    \big)
    ^\top \mv{e}
    =
    % \left(
    \big(
        (\pptfrac{\mv{e}}{\mv{u}})
        \mm{P}
        (\pptfrac{\estNN}{\estwth})
    \big)
    % \right)
    ^\top \mv{e}
$}\DIFaddend . 
Since the objective function depends on \DIFdelbegin \DIFdel{$\fe$ }\DIFdelend \DIFaddbegin \DIFadd{the tracking error $\mv{e}$ }\DIFaddend of a dynamic system, obtaining \DIFdelbegin \DIFdel{$\pptfrac{\fe}{\rbu}$ }\DIFdelend \DIFaddbegin \DIFadd{$\pptfrac{\mv{e}}{\mv{u}}$ }\DIFaddend is not straightforward. 
\DIFdelbegin \DIFdel{The recommended method to calculate the exact value of $\pptfrac{\fe}{\rbu}$ is to use the }\DIFdelend \DIFaddbegin \DIFadd{Exact computation methods, such as the }\DIFaddend forward sensitivity method \DIFdelbegin \DIFdel{\mbox{%DIFAUXCMD
\cite{Sengupta:2014aa,Ryu:2024aa} }\hskip0pt%DIFAUXCMD
by simulating the sensitivity equation as follows: 
}%DIFDELCMD < \begin{equation}
%DIFDELCMD <     \ddtt 
%DIFDELCMD <     \left(
%DIFDELCMD <         \pptfrac{\fe}{\rbu}
%DIFDELCMD <     \right)
%DIFDELCMD <     =
%DIFDELCMD <     \pptfrac{}{
%DIFDELCMD <         \rbu
%DIFDELCMD <     }
%DIFDELCMD <     \left(
%DIFDELCMD <         \rbM^{-1}(
%DIFDELCMD <             -\rbVm \fe
%DIFDELCMD <             -\mm K \fe
%DIFDELCMD <             + \mv f
%DIFDELCMD <             -\rbd + \mysat(\rbu)
%DIFDELCMD <         )
%DIFDELCMD <     \right)
%DIFDELCMD <     .
%DIFDELCMD < \end{equation}
%DIFDELCMD < %%%
\DIFdel{However, this cannot be realized, since we do not know the exact system dynamics, }%DIFDELCMD < \ie %%%
\DIFdel{$\rbM,\rbVm, \rbF$, and $\rbG$.
Additionally, the computational cost of the forward sensitivity method is high, as the number of DNN weightsis generally large.
}%DIFDELCMD < 

%DIFDELCMD < %%%
\DIFdel{In \mbox{%DIFAUXCMD
\cite{Douratsos:2007aa,Saerens:1991aa}}\hskip0pt%DIFAUXCMD
, the authors approximate $\pptfrac{\fe}{\rbu}$ }\DIFdelend \DIFaddbegin \DIFadd{\mbox{%DIFAUXCMD
\cite{Khalil:2002aa,Ryu:2025ab}}\hskip0pt%DIFAUXCMD
, can be used in principle, but they are often computationally demanding for real-time implementation with numerous DNN weights.
Therefore, several approximation methods have been proposed \mbox{%DIFAUXCMD
\cite{Douratsos:2007aa,Saerens:1991aa}}\hskip0pt%DIFAUXCMD
.
They approximate $\pptfrac{\mv{e}}{\mv{u}}$ }\DIFaddend by taking the sign of each entry \DIFaddbegin \DIFadd{of the control effectiveness matrix}\DIFaddend , \ie
\DIFdelbegin \DIFdel{$
    \pptfrac{\fe}{\rbu}
    \approx
    \left[
        \mysign
        \left(
            \pptfrac{r_i}{\tau_j}
        \right)
    \right]_{i,j\in\{1,\cdots,n\}}
$, 
assuming known controldirections}\DIFdelend \DIFaddbegin \DIFadd{$
  \pptfrac{\mv{e}}{\mv{u}}
  \approx
  \big[
    \mysign
    % \left(
    (
        % \pptfrac{e_i}{u_j}
    % \right)
    b_{i,j}
    )
  \big]_{i,j\in\{1,\cdots,n\}}
  ,
$
which corresponds to the MIT rule in the context of adaptive control}\DIFaddend .
However, \DIFdelbegin \DIFdel{this approach is not applicable to \eqref{eq:sys:filtered}, where control directions are unknown.
Instead, since the sign }\DIFdelend \DIFaddbegin \DIFadd{since the entry }\DIFaddend of the control \DIFdelbegin \DIFdel{input matrix is known—owing to the positive definiteness of $\rbM^{-1}$, see Property \ref{prop:M}—we approximate $\pptfrac{\fe}{\rbu}\approx \mm I_n$}\DIFdelend \DIFaddbegin \DIFadd{effectiveness matrix $\mm{B}(\mv{x})$ is unknown in this study, we instead approximate $\mm{B}(\mv{x})$ by its sign, }\ie \DIFadd{$\pptfrac{\mv{e}}{\mv{u}}\approx\mysign(\mm{B}(\mv{x}))=\mm{I}_{n}$, under Assumption~\ref{assum:B:bound}}\DIFaddend .
Consequently, the adaptation law in \eqref{eq:adap:th} becomes:
\DIFdelbegin %DIFDELCMD < \begin{equation}
%DIFDELCMD <     \ddtt{\estwth}
%DIFDELCMD <     \approx
%DIFDELCMD <     -\alpha 
%DIFDELCMD <     \left(
%DIFDELCMD <         {\pptfrac{\estNN}
%DIFDELCMD <         {\estwth}}^\top
%DIFDELCMD <         \fe
%DIFDELCMD <         +
%DIFDELCMD <         \textstyle\sum_{j\in\mathcal{I}}
%DIFDELCMD <         \lambda_j 
%DIFDELCMD <         \pptfrac{c_j}{\estwth}
%DIFDELCMD <     \right)
%DIFDELCMD <     .
%DIFDELCMD <     \label{eq:adap:th:approx}
%DIFDELCMD < \end{equation}%%%
\DIFdelend \DIFaddbegin \begin{equation}
  \ddtt{\estwth}
  \approx
  -\alpha 
  \left(
    {\pptfrac{\estNN}
    {\estwth}}^\top
    \mm{P}
    \mv{e}
    +
    \textstyle\sum_{i\in\mathcal{K}}
    \sigma_i 
    \estwth_i
    +
    \textstyle\sum_{j\in\mathcal{I}}
    \lambda_j 
    \pptfrac{c^{u}_{j}}{\estwth}
  \right)
  .
  \label{eq:adap:th:approx}
\end{equation}\DIFaddend 

\DIFdelbegin \subsection{\DIFdel{Essential and Potential Constraint Candidates}}%DIFAUXCMD
\addtocounter{subsection}{-1}%DIFAUXCMD
%DIFDELCMD < \label{sec:sub:cstr} 
%DIFDELCMD < 

%DIFDELCMD < %%%
\DIFdel{This section introduces essential constraints required to ensure the stability of the proposed CONAC, as well as conditions for potential input constraints.
}%DIFDELCMD < 

%DIFDELCMD < %%%
\DIFdel{First, weight norm constraints $c_{\theta_i}(\estwth),\ \forall i\in\{0,\cdots ,k\}$, defined in \eqref{eq:cstr:weight:ball}, are essential to prevent the weights from diverging during the adaptation process:
}%DIFDELCMD < \begin{equation}
%DIFDELCMD <     c_{\theta_i}
%DIFDELCMD <     =
%DIFDELCMD <     \tfrac{1}{2}
%DIFDELCMD <     \left(
%DIFDELCMD <         \Vert \estwth_i\Vert^2 
%DIFDELCMD <         -
%DIFDELCMD <         \overline\theta_i^2 
%DIFDELCMD <     \right)    
%DIFDELCMD <     % \le 0
%DIFDELCMD <     ,
%DIFDELCMD <     \label{eq:cstr:weight:ball}
%DIFDELCMD < \end{equation}
%DIFDELCMD < %%%
\DIFdel{with $\overline\theta_i\in\R_{>0}$ denoting the maximum allowable norm for $\estwth_i$, and the arguments of $c_{\theta_i}(\cdot)$ are omitted for brevity.
The gradient of $c_{\theta_i}$ with respect to $\estwth$ can be obtained by accumulating the gradients of each layer as follows:
}%DIFDELCMD < \begin{equation}
%DIFDELCMD <     \pptfrac{c_{\theta_i}}{\estwth_j} 
%DIFDELCMD <     =
%DIFDELCMD <     \begin{cases}
%DIFDELCMD <         \estwth_i,
%DIFDELCMD <         &
%DIFDELCMD <         \text{if } i=j,
%DIFDELCMD <         \\
%DIFDELCMD <         \mv{0},
%DIFDELCMD <         &
%DIFDELCMD <         \text{if } i\ne j
%DIFDELCMD <         .
%DIFDELCMD <     \end{cases} 
%DIFDELCMD <     \label{eq:cstr:weight:ball:grad}
%DIFDELCMD < \end{equation}
%DIFDELCMD < 

%DIFDELCMD < %%%
\DIFdel{Next, we present the following assumptions, which specify the conditions for the potential input constraints introduced in Appendix \ref{sec:appen:cstr}:
}%DIFDELCMD < 

%DIFDELCMD < \begin{assum}
%DIFDELCMD <     %%%
\DIFdel{The constraint function $c_j(\estwth)$ is convex in the (input-) $\tau$-space and satisfies $c_j(0) \le 0$ }\DIFdelend \DIFaddbegin \DIFadd{Additionally, the adaptation laws \eqref{eq:adap} provides an optimality interpretation at equilibrium.
At an equilibrium of \eqref{eq:adap}, where $\ddtt{\estwth}=\mv{0}$, $\ddtt\sigma_i=0$, }\DIFaddend and \DIFdelbegin \DIFdel{$c_j(\idealwth)\le 0$.
    }%DIFDELCMD < \label{assum:convex}
%DIFDELCMD < \end{assum}
%DIFDELCMD < %%%
\DIFdelend \DIFaddbegin \DIFadd{$\ddtt\lambda_j=0$, the resulting equilibrium conditions correspond to KKT conditions, }\ie \DIFadd{$\pptfrac{L}{\estwth}=0$, $c^{\theta}_i,c^{u}_j \le 0$, $\sigma_i,\lambda_j \ge 0$, and $\sigma_i c^{\theta}_i=0$, and $\lambda_j c^{u}_j=0,\ \forall i\in\mathcal{K},j\in\mathcal{I}$.
Thus, the equilibria characterize candidates for locally optimal solutions of the formulated constrained adaptation problem.
}\DIFaddend 

\DIFdelbegin %DIFDELCMD < \begin{assum}
%DIFDELCMD <     %%%
\DIFdel{The selected constraints satisfy the Linear Independence Constraint Qualification (LICQ) \mbox{%DIFAUXCMD
\cite[Chap.~12, Def.~12.1]{Nocedal:2006aa}}\hskip0pt%DIFAUXCMD
.
    }%DIFDELCMD < \label{assum:LICQ}
%DIFDELCMD < \end{assum}
%DIFDELCMD < %%%
\DIFdelend %DIF >  \begin{remark}[Verification of Approximated Adaptation Law]
%DIF >      In traditional adaptive control literature, by assuming known sign 

\DIFdelbegin %DIFDELCMD < \begin{remark}
%DIFDELCMD <     %%%
\DIFdel{Assumption \ref{assum:convex} is not restrictive, since, as aforementioned, the control input saturation is practically convex.
Moreover, Assumption \ref{assum:LICQ} is a standard assumption in optimization problems, ensuring that the gradients of the active constraints are linearly independent.
This assumption will be used in the stability analysis (see Lemma \ref{lem:convex:angle}).
}%DIFDELCMD < \end{remark}
%DIFDELCMD < %%%
\DIFdelend %DIF >      Furthermore, the approximated adaptation law \eqref{eq:adap:th:approx} coincides with the various existing adaptation laws in the literature when there are no constraints, \ie $\sigma_i=0,\ \forall i\in\mathcal{K}$ and $\lambda_j=0,\ \forall j\in\mathcal{I}$, derived using Lyapunov stability theory \cite{} and static approximation methods \cite{}.

\DIFdelbegin \DIFdel{The algorithm of the proposed CONAC is presented by Algorithm \ref{alg:CONAC}, implemented in a discrete-time controller with a sampling time of~$T_s$, where $(\cdot^{(k)})$ denotes the value at the $k$th time step.
}\DIFdelend %DIF >  \begin{remark}[Locally Lipschitz Continuity of the Adaptation Law] \label{rem:adap:lip}
%DIF >      The approximated adaptation law \eqref{eq:adap:th:approx} is locally Lipschitz continuous under the stated regularity assumptions.
%DIF >      Specifically, the error dynamics are locally Lipschitz, the DNN Jacobian $\pptfrac{\estNN}{\estwth}$ is locally bounded due to the use of $\tanh(\cdot)$ activation functions, and the input constraints $c_j^u(\mv{u})$ are $\mathcal{C}^1$ with respect to $\estwth$, while $\mv u=\mv P\estNN(\cdot;\estwth)$ is smooth in $\estwth$.
%DIF >      This property, in conjunction with the boundedness of signals established via subsequent Lyapunov analysis in Section~\ref{sec:stability:analysis}, precludes the finite escape time of the closed-loop system under the proposed CONAC
%DIF >  \end{remark}

\begin{algorithm}[!t]
  \caption{Proposed CONAC}\label{alg:CONAC}
  \begin{algorithmic}[1]
    \STATE \textbf{Initialize:} \DIFdelbegin \DIFdel{DNN weights }\DIFdelend $\estwth^{(0)}$, \DIFdelbegin \DIFdel{Lagrange multipliers $\lambda_j^{(0)}$
        }%DIFDELCMD < \FOR{$k = 0 $ to $\infty$}
%DIFDELCMD <             %%%
\DIFdelend \DIFaddbegin \DIFadd{$\sigma_i^{(0)},\ \forall i\in\mathcal{K}$, $\lambda_j^{(0)},\ \forall j\in\mathcal{I}$
    }\FOR{$r = 0 $ to $\infty$}
      \DIFaddend \STATE \DIFdelbegin \DIFdel{Measure $\q^{(k)},\dq^{(k)}$
            }%DIFDELCMD < \STATE %%%
\DIFdel{Obtain $\qd^{(k)},\dqd^{(k)}$ from given $\qd(t)$
            }\DIFdelend \DIFaddbegin \DIFadd{Obtain $\mv{x}^{(r)}$ and $\mv{x}_d^{(r)}$
      }\DIFaddend \STATE \DIFdelbegin \DIFdel{$\mv{e}^{(k)}\leftarrow\q^{(k)}-\qd^{(k)}$ and $\ddtt\mv{e}^{(k)}\leftarrow\ddtt\q^{(k)}-\dqd^{(k)}$
            }%DIFDELCMD < \STATE %%%
\DIFdel{$\fe^{(k)} \leftarrow \ddtt{\mv e^{(k)}}+\mm\Lambda \mv e^{(k)}$ based on \eqref{eq:err:filtered}
}\DIFdelend \DIFaddbegin \DIFadd{$\mv{e}^{(r)}\leftarrow\mv{x}^{(r)}-\mv{x}_d^{(r)}$
}\DIFaddend 

      \STATE \textit{// Neuro-adaptive control law (Section\DIFdelbegin \DIFdel{\ref{sec:ctrl design}}\DIFdelend \DIFaddbegin \DIFadd{~\ref{sec:main:result:sub:ctrl:design}}\DIFaddend )}

        %DIF >  \STATE $\NNin^{(r)}\leftarrow
        %DIF >    \left(
        %DIF >      {\mv{x}_d^{(r)}}^\top,
        %DIF >      {\ddtt{\mv{x}}_d^{(r)}}^\top,
        %DIF >      {\mv{x}^{(r)}}^\top,
        %DIF >      1
        %DIF >    \right)^\top$
        \STATE \DIFdelbegin \DIFdel{${\q}_n^{(k)}\leftarrow
                    \left(
                        {\q^{(k)}}^\top,{\dq^{(k)}}^\top,{\fe^{(k)}}^\top,1
                    \right)^\top$
                }\DIFdelend \DIFaddbegin \DIFadd{Define $\NNin^{(r)}$ from all arguments of $\mv{u}_d^{(r)}$
}\DIFaddend 

        \STATE \DIFdelbegin \DIFdel{$\rbu^{(k)}\leftarrow\NN({\q}_n^{(k)};\estwth^{(k)}) $ based on \eqref{eq:control:est}
                }\DIFdelend \DIFaddbegin \DIFadd{$\mv{u}^{(r)}\leftarrow\NN(\NNin^{(r)};\estwth^{(r)}) $ based on \eqref{eq:ctrl:NN:actual}
        }\DIFaddend \STATE Apply \DIFdelbegin \DIFdel{$\rbu^{(k)}$ }\DIFdelend \DIFaddbegin \DIFadd{$\mv{u}^{(r)}$ }\DIFaddend to the system in \DIFdelbegin \DIFdel{\eqref{eq:sys1}
}\DIFdelend \DIFaddbegin \DIFadd{\eqref{eq:sys}
}\DIFaddend 

      \STATE \textit{// \DIFdelbegin \DIFdel{Train }\DIFdelend \DIFaddbegin \DIFadd{Update }\DIFaddend DNN \DIFaddbegin \DIFadd{weights }\DIFaddend (Section\DIFdelbegin \DIFdel{\ref{sec:adap_laws}}\DIFdelend \DIFaddbegin \DIFadd{~\ref{sec:main:result:sub:adap:laws}}\DIFaddend )}

      \STATE Compute \DIFdelbegin \DIFdel{$\pptfrac{\NN}{\estwth^{(k)}}$ }\DIFdelend \DIFaddbegin \DIFadd{$\pptfrac{\NN}{\estwth^{(r)}}$ }\DIFaddend as defined in \eqref{eq:gradient:phi}
      \DIFdelbegin \DIFdel{, $c_j$ and $\pptfrac{c_j}{\estwth^{(k)}}$ for all $j\in\mathcal{I}$          
            }%DIFDELCMD < 

%DIFDELCMD <             %%%
\DIFdelend \STATE \DIFdelbegin \DIFdel{$\estwth^{(k+1)}
                    \leftarrow
                    \estwth^{(k)}
                    -\alpha 
                    \left(
                        {\pptfrac{\estNN}
                        {\estwth^{(k)}}}^\top
                        \fe^{(k)}
                        +
                        \textstyle\sum_{j\in\mathcal{I}}
                        \lambda_j^{(k)} 
                        \pptfrac{c_j}{\estwth^{(k)}}
                    \right)\cdot T_s$ based on \eqref{eq:adap:th}
}\DIFdelend \DIFaddbegin \DIFadd{Compute $c^{\theta}_i,c^{u}_j$ and $\pptfrac{c^{u}_{j}}{\estwth^{(r)}}$ for all $i\in\mathcal{K}, j\in\mathcal{I}$
      }\DIFaddend 

      \STATE \DIFdelbegin \DIFdel{$\lambda_j^{(k+1)}
            \leftarrow
            \lambda_j^{(k)}
            +
            \beta_j c_j\cdot T_s$ for all $j\in\mathcal{I}$\,based\,on\,\eqref{eq:adap:lbd}}%DIFDELCMD < \STATE %%%
\DIFdel{$\lambda_j^{(k+1)} \leftarrow \max(\lambda_j^{(k+1)},0)$ for all $j\in\mathcal{I}$ from \eqref{eq:adap:lbd:max} 
}%DIFDELCMD < 

%DIFDELCMD <         %%%
\DIFdelend \DIFaddbegin \DIFadd{Update $\estwth^{(r+1)}$, $\sigma_i^{(r+1)},\ \forall i\in\mathcal{K}$, and $\lambda_j^{(r+1)},\ \forall j\in\mathcal{I}$ using \eqref{eq:adap:th:approx}, \eqref{eq:adap:sig}, and \eqref{eq:adap:lbd}, respectively.
    }\DIFaddend \ENDFOR
  \end{algorithmic}
\end{algorithm}

%DIF <   SECTION STABILITY ANALYSIS ==============================
\DIFaddbegin \DIFadd{The proposed CONAC algorithm is presented in Algorithm~\ref{alg:CONAC}, implemented in a discrete-time controller, where $(\cdot^{(r)})$ denotes the value at the $r$th time step.
}

%DIF >  ------------------------------------
%DIF >        SECTION STABILITY ANALYSIS
%DIF >  ------------------------------------
\DIFaddend \section{Stability Analysis} \DIFdelbegin %DIFDELCMD < \label{sec:stability}
%DIFDELCMD < %%%
\DIFdelend \DIFaddbegin \label{sec:stability:analysis}
\DIFaddend 

\DIFdelbegin \DIFdel{Before conducting the stability analysis, let us define the weight estimation error as $\tilde\wth:= [\tilde\wth_i]_{i\in\{0,\cdots,k\}}$, where $\tilde\wth_i:=\estwth_i-\idealwth_i, \forall i\in\{0,\cdots,k\}$.
The following lemmas are introduced for }\DIFdelend \DIFaddbegin \DIFadd{In this section, the stability of the closed-loop system under the proposed CONAC is analyzed.
To facilitate }\DIFaddend the stability analysis\DIFaddbegin \DIFadd{, we introduce the following lemmas}\DIFaddend .

\DIFdelbegin %DIFDELCMD < \begin{lem}
%DIFDELCMD <     %%%
\DIFdel{If Assumptions \ref{assum:convex} and \ref{assum:LICQ} hold, then for each active constraint $c_j$, the angle between the output layer 's weight vector }\DIFdelend \DIFaddbegin \begin{lem}[Convexity of input constraints with respect to $\estwth_k$] \label{lem:convexity:input:constraint}
  \DIFadd{The input constraint function $c_j^{u}(\estwth_k),\ \forall j \in \mathcal{I},$ is convex with respect to the output layer weights }\DIFaddend $\estwth_k$\DIFdelbegin \DIFdel{and $\pptfrac{c_j}{\estwth_k}$ is positive; that is $\pptfrac{c_j}{\estwth_k}^\top\estwth_k>0$.
    }%DIFDELCMD < \label{lem:convex:angle}
%DIFDELCMD < %%%
\DIFdelend \DIFaddbegin \DIFadd{, resulting in the inequality 
  }\begin{equation*}
    c_j^{u}(\estwth_k) 
    - 
    \pptfrac{c_j^{u}}{\estwth_k}^\top 
    \estwth_k 
    \le 
    % c_j^{u}(\estwth)\big|_{\estwth_k=\mv{0}}
    - \underline{c}_j^{u}
    <
    0
    .
  \end{equation*}
\DIFaddend \end{lem}

\begin{proof}
  \DIFdelbegin %DIFDELCMD < 

%DIFDELCMD < %%%
\DIFdel{Since $\rbu = \estNN$, a linear mapping $\mm T(\cdot):\estwth_k\to\rbu$, which is independent of }\DIFdelend %DIF >  The dependence of the preceding layers and $\NNin$ on the output layer appears through the activation vector $\act_k$.
  %DIF >  According to the DNN gradient in \eqref{eq:gradient:phi}, we have
  \DIFaddbegin \DIFadd{According to Proposition~\ref{propsit:kron}, the output of the DNN $\estNN$ depends linearly on the output layer weights }\DIFaddend $\estwth_k$ \DIFdelbegin \DIFdel{, can be derived using Proposition \ref{propsit:kron}:
  }%DIFDELCMD < \begin{equation}\label{eq:linear:map}
%DIFDELCMD <     \begin{aligned}
%DIFDELCMD <     \rbu 
%DIFDELCMD <     = 
%DIFDELCMD <     &
%DIFDELCMD <     \estNN 
%DIFDELCMD <     % = 
%DIFDELCMD <     % \myvec(\estNN)
%DIFDELCMD <     =
%DIFDELCMD <     \myvec(
%DIFDELCMD <         \estwV_k^\top \estact_k
%DIFDELCMD <     ) 
%DIFDELCMD <     = 
%DIFDELCMD <     (
%DIFDELCMD <         \mm I_{l_{k+1}}
%DIFDELCMD <         \otimes 
%DIFDELCMD <         \estact_k^\top
%DIFDELCMD <     )^\top
%DIFDELCMD <     \myvec(\estwV_k)
%DIFDELCMD <     \\
%DIFDELCMD <     = &
%DIFDELCMD <     \underbrace{
%DIFDELCMD <         (
%DIFDELCMD <         \mm I_{l_{k+1}}
%DIFDELCMD <         \otimes 
%DIFDELCMD <         \estact_k^\top
%DIFDELCMD <     )^\top}_{=:\mm T(\estact_k)}
%DIFDELCMD <     \estwth_k 
%DIFDELCMD <     =
%DIFDELCMD <     \mm T(\estact_k) \estwth_k.
%DIFDELCMD <     \end{aligned}
%DIFDELCMD < \end{equation}%%%
\DIFdelend \DIFaddbegin \DIFadd{as follows:
  }\begin{equation} \label{eq:NN:linear}
    \estNN
    =
    \estwV_k^\top
    \estact_k
    = 
    \left(
      \mm{I}_{l_{k+1}}
      \!
      \otimes
      \! 
      \estact_k^\top
    \right)
    \myvec(\estwV_k)
    = 
    \left(
      \mm{I}_{l_{k+1}} 
      \!
      \otimes 
      \!
      \estact_k^\top
    \right)
    \estwth_k
    .
  \end{equation}\DIFaddend 
  \DIFaddbegin \DIFadd{This implies that the control input $\mv{u}=\mm{P}\estNN$ depends linearly on $\estwth_k$.
  }\DIFaddend Therefore, the convexity of \DIFdelbegin \DIFdel{the input constraints in the $\rbu$-space (see Assumption \ref{assum:convex}) is preserved in }\DIFdelend \DIFaddbegin \DIFadd{$c_j^u(\estwth),\ \forall j \in \mathcal{I}$, with respect to $\mv{u}$ is preserved with respect to }\DIFaddend $\estwth_k$\DIFdelbegin \DIFdel{-space, implying
that $\pptfrac{c_j}{\estwth_k}^\top\estwth_k>0$.
  }%DIFDELCMD < 

%DIFDELCMD < %%%
\DIFdelend \DIFaddbegin \DIFadd{, due to the linear dependence of $\mv{u}$ on $\estwth_k$.
  In addition, the convexity of $c_j^u(\estwth)$ with respect to $\estwth_k$ implies 
  %DIF >  \begin{equation*}
  $
    c_j^{u}(\estwth_k) 
    - 
    \pptfrac{c_j^{u}}{\estwth_k}^\top 
    \estwth_k 
    \le 
    c_j^{u}(\estwth)\big|_{\estwth_k=\mv{0}}
    % - \underline{c}_j^{u}
    % <
    % 0
    .
  $
  %DIF >  \end{equation*}
  %DIF >  , where $c_j^u(\mv{0})$ denotes the value at $\estwth_k=\mv{0}$
  Additionally, since $\estwth_k=\mv{0}$ implies $\mv{u}=\mm{P}\estNN=\mv{0}$, (}\ie \DIFadd{output layer weight is zero), we have $c_j^{u}(\estwth)\big|_{\estwth_k=\mv{0}}=-\underline{c}_j^{u}<0$ by Assumption~\ref{assum:sat}, which concludes the proof.
}\DIFaddend \end{proof}

\DIFdelbegin %DIFDELCMD < \begin{lem} 
%DIFDELCMD <     %%%
\DIFdel{If input constraint $c_j(\estwth),\ \forall j\in\mathcal I \setminus \{\theta_i\}_{i\in\{0,\cdots,k\}}$, satisfies Assumption \ref{assum:convex}, then $\norm{\pptfrac{c_j}{\estwth_i}},\ \forall i\in\{k-1,\cdots, 0\}$, is bounded, provided the norms of $\estwth_i,\ \forall i\in\{k,\cdots, i+1\}$, remain bounded.
}%DIFDELCMD < \label{lem:cstr:grad:bound}
%DIFDELCMD < %%%
\DIFdelend \DIFaddbegin \begin{lem}[Boundedness of $\Delta\mv{u}$ with respect to $\estwth_k$] \label{lem:delta:u:bound}
  \DIFadd{The saturation of control input $\Delta\mv{u}$ is bounded by $\estwth_k$ as $\norm{\Delta\mv{u}}\le \overline{p}\sqrt{l_k+1}\norm{\estwth_k}$.
}\DIFaddend \end{lem}

\DIFdelbegin \DIFdel{For instance, if $k=3$ and $i=1$, $\norm{\pptfrac{c_j}{\estwth_1}}$ is bounded , provided that $\norm{\estwth_3}$ and $\norm{\estwth_2}$ are bounded.
}%DIFDELCMD < 

%DIFDELCMD < %%%
\DIFdelend \begin{proof}
  \DIFdelbegin %DIFDELCMD < 

%DIFDELCMD < %%%
\DIFdel{The gradient of $c_j,\ \forall j\in\mathcal I \setminus \{\theta_i\}_{i\in\{0,\cdots,k\}}$, with respect to $\estwth_i$ is represented as
}%DIFDELCMD < \begin{equation}
%DIFDELCMD <     \pptfrac{c_j}{\estwth_i} 
%DIFDELCMD <     = 
%DIFDELCMD <     \pptfrac{c_j}{\rbu} 
%DIFDELCMD <     \pptfrac{\rbu}{\estNN} 
%DIFDELCMD <     \pptfrac{\estNN}{\estwth_i}
%DIFDELCMD <     = 
%DIFDELCMD <     \pptfrac{c_j}{\rbu} 
%DIFDELCMD <     \mm{I}_n
%DIFDELCMD <     \pptfrac{\estNN}{\estwth_i}
%DIFDELCMD <     .
%DIFDELCMD < \end{equation}%%%
\DIFdelend \DIFaddbegin \DIFadd{Since $\act_k(\cdot)$ employs the bounded activation function $\vert\tanh(\cdot)\vert < 1$, we have 
  }\begin{equation}
    \norm{\act_k(\cdot)} \le \sqrt{l_k+1}
    ,
  \end{equation}\DIFaddend 
  \DIFdelbegin \DIFdel{The boundedness of the first term $\pptfrac{c_j}{\rbu}$ is guaranteed, if the input argument $\rbu$ is bounded owing to the convexity of $c_j$ by Assumption \ref{assum:convex}.
}\DIFdelend \DIFaddbegin \DIFadd{for any output of the preceding layers and $\NNin$.
  Hence, invoking the relation in \eqref{eq:NN:linear}, the control input $\mv{u}=\mm{P}\estNN$ is bounded by $\estwth_k$ as follows:
  }\begin{equation} \label{eq:NN:bound}
    \norm{\mv{u}} 
    \le
    \overline{p}
    \norm{\estNN}
    \le
    \overline{p}
    \sqrt{l_k+1}\norm{\estwth_k}
    .   
  \end{equation}

  \DIFaddend In addition, \DIFdelbegin \DIFdel{the boundedness of $\rbu$ can be ensured by the boundedness of $\estwth_k$, since $\rbu=\estNN=(\mm I_{l_{k+1}}\otimes \estact_k^\top)\estwth_k$ }\DIFdelend \DIFaddbegin \DIFadd{since there exists $\underline{u}>0$ such that $\inf_{\mv u\notin\Omega_u}\allowbreak(\norm{\mysat(\mv u)})=\underline{u}$ by Assumption~\ref{assum:sat}, we have
  $
    \norm{\Delta\mv u}
    \le
    \norm{\mv u}-\underline{u}
    ,
  $ 
  while $\Delta\mv u=\mv 0$ when $\mv u\in\Omega_u$.
  Thus, we have
  }\begin{equation}
    \norm{\Delta\mv{u}}
    \le
    \max
    \left\{
      \overline{p}\sqrt{l_k+1}\norm{\estwth_k}
      -
      \underline{u}
      ,
      0
    \right\}
    \le
    \overline{p}\sqrt{l_k+1}\norm{\estwth_k}
    ,  
  \end{equation}
  \DIFadd{which concludes the proof.    
}\end{proof}

\begin{lem} \label{lem:proj:ineq}
  \DIFadd{For any $a,b\in\R_{\ge0}$ }\DIFaddend and \DIFdelbegin \DIFdel{$\norm{\estact_k}$ is bounded due to the properties of the activation functions.
  The third term, $\pptfrac{\estNN}{\estwth_i}$ is bounded, provided that the norms of $\estwth_i,\ \forall i \in\{k,\cdots,i+1\}$, are bounded.    
This }\DIFdelend \DIFaddbegin \DIFadd{$c\in\R_{>0}$, the following inequality holds:
  }\begin{equation}
    (a-b)\myproj_{a\ge0}(c) \le (a-b) c.
  \end{equation}
\end{lem}

\begin{proof}
  \DIFadd{The inequality }\DIFaddend can be verified by \DIFdelbegin \DIFdel{using the definition of $\pptfrac{\estNN}{\estwth_i}$ given in \eqref{eq:NN:grad}.
Consequently, the boundedness of $\norm{\pptfrac{c_j}{\estwth_i}},\ \forall j\in\mathcal I \setminus \{\theta_i\}_{i\in\{0,\cdots,k\}}$, depends on the boundedness of $\estwth_i,\ \forall i\in \{k,\cdots,i+1\}$.
}%DIFDELCMD < 

%DIFDELCMD < %%%
\DIFdelend \DIFaddbegin \DIFadd{considering the two cases of $a>0$ and $a=0$, as follows:
  }\begin{equation*}
    (a-b)\myproj_{a\ge0}(c) 
    =
    \begin{cases}
      (a-b) c, & \text{if } a>0,
      \\
      0, & \text{if } a=0,
    \end{cases}
  \end{equation*}
  \DIFadd{This implies that $(a-b)\myproj_{a\ge0}(c) \le (a-b) c$ always holds and concludes the proof.
}\DIFaddend \end{proof}

The following theorem \DIFdelbegin \DIFdel{states that $\fe$ and }\DIFdelend \DIFaddbegin \DIFadd{establishes the UUB of the tracking error $\mv{e}$, the DNN weights }\DIFaddend $\estwth$\DIFdelbegin \DIFdel{are bounded}\DIFdelend \DIFaddbegin \DIFadd{, and the Lagrange multipliers $\sigma_i,\ \forall i \in \mathcal{K}$, and $\lambda_j,\ \forall j \in \mathcal{I}$}\DIFaddend .

\begin{theorem} \DIFdelbegin \DIFdel{For the dynamical system described in \eqref{eq:sys1}, the NAC in \eqref{eq:control:est} with the weight adaptation laws in \eqref{eq:adaptation law} ensure the boundedness of the filtered error $\fe$ and }\DIFdelend \DIFaddbegin \label{thm:stability}
  \DIFadd{Consider the closed-loop system consisting of \eqref{eq:sys}, the control law \eqref{eq:ctrl:NN:actual}, and the adaptation laws \eqref{eq:adap}.
  Suppose that Assumptions~\ref{assum:B:bound}--\ref{assum:xd} hold and the initial condition belongs to the compact invariant set 
  $
  \bigcap_{i\in\mathcal K}
  \mathcal S_i(\overline V_i)
  $ 
  constructed in \eqref{eq:Sk:def} and \eqref{eq:Si:def}.
  Then, the tracking error $\mv e$, }\DIFaddend the \DIFdelbegin \DIFdel{weight estimate }\DIFdelend \DIFaddbegin \DIFadd{DNN weights }\DIFaddend $\estwth$, \DIFdelbegin \DIFdel{under the input constraints satisfying Assumption \ref{assum:convex} and \ref{assum:LICQ}, provided that the weight norm constraints \eqref{eq:cstr:weight:ball} are imposed}\DIFdelend \DIFaddbegin \DIFadd{and the Lagrange multipliers $\sigma_i,\ \forall i \in \mathcal{K},$ and $\lambda_j,\ \forall j \in \mathcal{I},$ are uniformly ultimately bounded}\DIFaddend .
\end{theorem}

\begin{proof}

%DIF <  The boundednesses of $\estwth$ and $\fe$ are analyzed recursively from the last $k$th layer to the first layer of $\estNN$. 
\DIFdelbegin \DIFdel{Invoking the fact that the amplitude of the control input depends only on }\DIFdelend \DIFaddbegin \DIFadd{The proof is divided into two steps.
We first analyze the UBB of the tracking error $\mv{e}$, the output-layer adaptive variables, }\DIFaddend $\estwth_k$, \DIFdelbegin %DIFDELCMD < \ie %%%
\DIFdel{$\rbu=\estNN=(\mm I_{l_{k+1}}\otimes \estact_k^\top)\estwth_k$ and $\norm{\estact_k}$ is bounded, the boundedness of inner layers' weights will be established after proving the boundedness of $\fe$ and $\estwth_k$.
In addition, without loss of generality, weight norm constraints $c_{\theta_i},\ \forall i\in\{0,\cdots,k\}$, are supposed to be active.
This assumption is justified because, even if $c_{\theta_i}$ is inactive, $\estwth$ is adapted to minimize the objective function $J$ as long as the constraint is inactive.
We consider two cases: (1) the control input is saturated, and (2) the control input is not saturated}\DIFdelend \DIFaddbegin \DIFadd{and $\sigma_k$, and the Lagrange multipliers associated with the input constraints $\lambda_j,\ \forall j\in\mathcal{I}$.
This step provides the key boundedness result for the tracking error and the controller output.
In the following step, we analyze the remaining inner-layer adaptive variables, $\estwth_i$ and $\sigma_i,\ \forall i\in\mathcal{K}\setminus\{k\}$}\DIFaddend .

\DIFdelbegin \subsection*{\DIFdel{Case 1: Control input is saturated.}}
%DIFAUXCMD
\DIFdelend \DIFaddbegin \hfill
\DIFaddend 

\DIFdelbegin \DIFdel{As a result of the }\DIFdelend \DIFaddbegin \textit{\DIFadd{Step~1: Boundedness of the tracking error, output-layer and input-constraint adaptive variables.}}

\DIFadd{Consider the following Lyapunov function 
}\begin{equation} \label{eq:Vk:def}
  V_{k}(\mv{z}_k)
  :=
  V_c
  +
  \tfrac{1}{2\alpha}\estwth_k^\top\estwth_k
  +
  \tfrac{1}{\beta_k^\theta}(\sigma_k-\sigma_k^o)^2
  +
  \textstyle\sum_{j\in\mathcal{I}}
  \tfrac{1}{2\beta_j^{u}}\lambda_j^2
  ,
\end{equation}
\DIFadd{where 
$
  \mv{z}_k
  :=
  \left(
    \mv{e}^\top,
    \estwth_k^\top,
    \sigma_k,
    \left[
      \lambda_j
    \right]_{j\in\mathcal{I}}^\top
  \right)^\top
$ and $\sigma_k^o\in\R_{>0}$ will be determined later.
The Lyapunov function $V_k:=V_k(\mv{z}_k)$ is lower bounded and radially unbounded with respect to $\mv{z}_k$.
In the following analysis, we consider the case where the output layer weight constraint is violated, }\ie \DIFadd{$c_k^{\theta}(\estwth_k)>0$.
%DIF >  If this condition does not hold, then $c_k^{\theta}(\estwth_k)\le0$ implies $\norm{\estwth_k}\le\overline{\theta}_k$, and the boundedness of $\estwth_k$ follows directly.
This is because the boundedness of $\estwth_k$ directly follows when the output layer weight constraint is not violated, }\ie \DIFadd{$c_k^{\theta}(\estwth_k)\le0$.
}

\DIFadd{The }\DIFaddend time derivative of \DIFdelbegin \DIFdel{$V_1$ in \eqref{eq:lya1:dot2}, the invariant set $\Theta_{\fe}^{1}$ of $\fe$, }%DIFDELCMD < \ie %%%
\DIFdel{if $\fe$ leaves $\Theta_{\fe}^{1}$, $\ddtt V_1$ is negative, can be obtained using Lemma \ref{lem:stable:set} }\DIFdelend \DIFaddbegin \DIFadd{$V_{k}$ can be analyzed }\DIFaddend as follows:
\DIFdelbegin %DIFDELCMD < \begin{equation}
%DIFDELCMD <     \Theta_{\fe}^{1} 
%DIFDELCMD <     = 
%DIFDELCMD <     \left\{ 
%DIFDELCMD <         \fe \in \R^n 
%DIFDELCMD <         \mid 
%DIFDELCMD <         \norm{\fe} 
%DIFDELCMD <         \le 
%DIFDELCMD <         \tfrac{
%DIFDELCMD <             \overline\tau_d
%DIFDELCMD <             +
%DIFDELCMD <             \overline\tau
%DIFDELCMD <             +
%DIFDELCMD <             \overline{\theta}_k^*\sqrt{l_{k}+1}
%DIFDELCMD <             +
%DIFDELCMD <             \overline{\epsilon}
%DIFDELCMD <         }{
%DIFDELCMD <             \lambda_{\min}(\mm K)
%DIFDELCMD <         }
%DIFDELCMD <     \right\}
%DIFDELCMD <     .
%DIFDELCMD <     \label{eq:invariant:c1:fe}
%DIFDELCMD < \end{equation}%%%
\DIFdelend \DIFaddbegin \begin{align} \label{eq:Vk:dot:1} \nonumber
  \ddtt V_{k}
  \le&
  -\underline{p}
  \norm{\mv{e}}^2
  +
  \mv{e}^\top\mm{P}
  \left(
    \estNN - \idealNN
    - \mv{\epsilon}
  \right)
  +
  \mv{e}^\top
  \Delta \mv{u}
  \\ \nonumber
  &
  -
  \underbrace{
    \estwth_k^\top
    \pptfrac{\estNN}{\estwth_k}^\top
    \mm{P}
    \mv{e}
  }_{
    = 
    \estNN^\top\mm{P}\mv{e}
    ,\ \text{see, \eqref{eq:gradient:phi}, \eqref{eq:NN:linear}}
  }
  -
  \sigma_k
  \estwth_k^\top
  \estwth_k
  -    
  \textstyle\sum_{j\in\mathcal{I}}
  \lambda_j
  \estwth_k^\top
  \pptfrac{c^{u}_j}{\estwth_k}
  \\ % \nonumber
  &
  +
  2(\sigma_k-\sigma_k^o) c_k^{\theta}(\estwth_k)
  +
  \textstyle\sum_{j\in\mathcal{I}}
  \lambda_j c^{u}_j(\estwth)
  \\ \nonumber
  \le&
  -\underline{p}
  \norm{\mv{e}}^2
  +
  \mv{e}^\top
  \mm{P}
  \left(
    - \idealNN
    - \mv{\epsilon}
  \right)
  +
  \mv{e}^\top
  \Delta \mv{u}
  -
  \sigma_k \norm{\estwth_k}^2
  \\ \nonumber
  &
  +
  (\sigma_k-\sigma_k^o) 
  \left(
    \norm{\estwth_k}^2
    -
    \overline{\theta}_k^{2}
  \right)
  -
  \underbrace{
    \textstyle\sum_{j\in\mathcal{I}}
    \lambda_j
    \underline{c}_j^u 
  }_{\text{see, Lemma~\ref{lem:convexity:input:constraint}}}
  \\ \nonumber
  =&
  -\underline{p}
  \norm{\mv{e}}^2
  +
  \mv{e}^\top
  \mm{P}
  \left(
    - \idealNN
    - \mv{\epsilon}
  \right)
  +
  \mv{e}^\top
  \Delta \mv{u}
  \\ \nonumber
  &
  -
  \sigma_k^o
  \norm{\estwth_k}^2
  -
  \sigma_k
  \overline{\theta}_k^2
  +
  \sigma_k^o
  \overline{\theta}_k^2
  -
  \textstyle\sum_{j\in\mathcal{I}}
  \lambda_j
  \underline{c}_j^u 
  .
\end{align}\DIFaddend 
\DIFdelbegin \DIFdel{It is notable that the ideal DNN $\idealNN_k={\idealwV_k}^\top\idealact_k$ is bounded by $\overline{\theta}_k^*\sqrt{l_{k}+1}$.
This is because, the ideal weight }\DIFdelend \DIFaddbegin \DIFadd{The projection operators in \eqref{eq:adap:sig} and \eqref{eq:adap:lbd} are omitted according to Lemma~\ref{lem:proj:ineq}.
}

\DIFadd{For the boundedness of }\DIFaddend $\idealwth_k$ \DIFdelbegin \DIFdel{is bounded by Assumption \ref{assum:feasible} such that $\norm{\idealwth_k}\le\overline{\theta}_k^*$ for some positive constant }\DIFdelend \DIFaddbegin \DIFadd{and $\mv{\epsilon}$, we assume that $\NNin$ remains in the approximation region $\Omega_{\NNin}$, which will be ensured later in the proof.
Hence, $\idealwth_k$ and $\mv{\epsilon}$ are bounded by }\DIFaddend $\overline{\theta}_k^*$ \DIFdelbegin \DIFdel{.
Hence, invoking $\norm{\idealwV_k}_F = \norm{\idealwth_k} \le \overline{\theta}_k^*$ and $\norm{\idealact_k}\le \sqrt{l_{k}+1}$, we have $\norm{\idealNN}\le \overline{\theta}_k^*\sqrt{l_{k}+1}$.
}\DIFdelend \DIFaddbegin \DIFadd{and $\overline{\epsilon}$, respectively.
Similarly to the bound of $\estNN$, $\idealNN$ is bounded by $\overline{\theta}_k^* \sqrt{l_k+1}$.
}\DIFaddend 

\DIFdelbegin \DIFdel{To investigate the boundedness of $\estwth_k$, consider the Lyapunov function candidate $V_2:=\tfrac{1}{2\alpha}\estwth_k^\top \estwth_k$.
Taking the time derivative of $V_2$ yields:
}%DIFDELCMD < \begin{equation}
%DIFDELCMD <     \begin{aligned}
%DIFDELCMD <         \ddtt V_2 
%DIFDELCMD <         =& 
%DIFDELCMD <         -\estwth_k^\top 
%DIFDELCMD <         \left(
%DIFDELCMD <             (\mm I_{l_{k+1}}\otimes \estact_k^\top)^\top
%DIFDELCMD <             \fe
%DIFDELCMD <             +
%DIFDELCMD <             \textstyle\sum_{j\in\mathcal{I}}
%DIFDELCMD <             \lambda_j 
%DIFDELCMD <             \pptfrac{c_j}{\estwth_k}
%DIFDELCMD <         \right)
%DIFDELCMD <         \\
%DIFDELCMD <         =&
%DIFDELCMD <         -\estwth_k^\top 
%DIFDELCMD <         \Big(
%DIFDELCMD <             (\mm I_{l_{k+1}}\otimes \estact_k^\top)^\top
%DIFDELCMD <             \fe
%DIFDELCMD <             +
%DIFDELCMD <             \lambda_{\theta_k}\estwth_k
%DIFDELCMD <         \\
%DIFDELCMD <         &
%DIFDELCMD <             +
%DIFDELCMD <             \textstyle\sum_{
%DIFDELCMD <                 j\in\mathcal{I}\setminus \{\theta_i\}_{i\in\{0,\cdots,k\}}
%DIFDELCMD <             }
%DIFDELCMD <             \lambda_j 
%DIFDELCMD <             \pptfrac{c_j}{\estwth_k}
%DIFDELCMD <         \Big)
%DIFDELCMD <         \\
%DIFDELCMD <         \le&
%DIFDELCMD <         -
%DIFDELCMD <         \lambda_{\theta_k}\norm{\estwth_k}^2
%DIFDELCMD <         +
%DIFDELCMD <         \underbrace{
%DIFDELCMD <             \norm{(\mm I_{l_{k+1}}\otimes {\estact}_k^\top)}
%DIFDELCMD <         }_{=:{\gamma}_1\in\R_{\ge0}}
%DIFDELCMD <         \norm{\fe}
%DIFDELCMD <         \norm{\estwth_k}
%DIFDELCMD <         \\
%DIFDELCMD <         &
%DIFDELCMD <         -
%DIFDELCMD <         \underbrace{
%DIFDELCMD <             \textstyle\sum_{
%DIFDELCMD <                 j\in\mathcal{I}\setminus \{\theta_i\}_{i\in\{0,\cdots,k\}}
%DIFDELCMD <             }
%DIFDELCMD <             \lambda_j 
%DIFDELCMD <             \estwth_k^\top \pptfrac{c_j}{\estwth_k}
%DIFDELCMD <         }_{=:{\gamma}_2\in\R_{\ge0},\ \text{by Lemma \ref{lem:convex:angle} and Assumption \ref{assum:LICQ}}}
%DIFDELCMD <         \\
%DIFDELCMD <         \le&
%DIFDELCMD <         -
%DIFDELCMD <         \lambda_{\theta_k}\norm{\estwth_k}^2
%DIFDELCMD <         +
%DIFDELCMD <         {\gamma}_1     
%DIFDELCMD <         \norm{\fe}
%DIFDELCMD <         \norm{\estwth_k}
%DIFDELCMD <         .
%DIFDELCMD <     \end{aligned}
%DIFDELCMD <     \label{eq:lya2:dot}
%DIFDELCMD < \end{equation}%%%
\DIFdelend \DIFaddbegin \DIFadd{Using Lemma~\ref{lem:delta:u:bound}, we can further analyze $\ddtt V_{k}$ as follows:
}\begin{align} \label{eq:Vk:dot:2} 
  \ddtt V_{k}
  \le&
    -\underline{p}
    \norm{\mv{e}}^2
    -
    \sigma_k^o
    \norm{\estwth_k}^2
  \\ \nonumber
  &
    +
    \overline{p}
    \norm{\mv{e}}
    \left(
      \overline{\theta}^*_k
      \sqrt{l_{k}+1}
      +
      \overline{\epsilon}
      +
      \sqrt{l_k+1}\norm{\estwth_k}
    \right)        
  \\ \nonumber
  &
    -
    \sigma_k 
    \overline{\theta}_k^2
    +
    \sigma_k^o 
    \overline{\theta}_k^2
    -
    \textstyle\sum_{j\in\mathcal{I}}
    \lambda_j
    \underline{c}_j^u
    .
\end{align}\DIFaddend 
\DIFdelbegin \DIFdel{According to \eqref{eq:lya2:dot} and Lemma\ref{lem:stable:set}, the invariant set of $\estwth_k$ is defined as }%DIFDELCMD < \begin{equation}
%DIFDELCMD <     \Theta_{\theta_k}^{1} 
%DIFDELCMD <     = 
%DIFDELCMD <     \left\{ 
%DIFDELCMD <         \estwth_k \in \R^{\Xi_{k}} 
%DIFDELCMD <         \mid 
%DIFDELCMD <         \norm{\estwth_k} 
%DIFDELCMD <         \le 
%DIFDELCMD <         \tfrac{
%DIFDELCMD <             {\gamma}_1
%DIFDELCMD <             (
%DIFDELCMD <                 \overline\tau_d
%DIFDELCMD <                 +
%DIFDELCMD <                 \overline\tau
%DIFDELCMD <                 +
%DIFDELCMD <                 \overline{\theta}^*_k\sqrt{l_{k}+1}
%DIFDELCMD <                 +
%DIFDELCMD <                 \overline{\epsilon}
%DIFDELCMD <             )
%DIFDELCMD <         }{
%DIFDELCMD <             \lambda_{\theta_k}\lambda_{\min}(\mm K)
%DIFDELCMD <         }
%DIFDELCMD <     \right\}
%DIFDELCMD <     .
%DIFDELCMD <     \label{eq:invariant:c1:wth}
%DIFDELCMD < \end{equation}%%%
\DIFdelend \DIFaddbegin 

\DIFadd{For further analysis, we apply Young's inequality which states 
%DIF >  \begin{equation}
$
  ab
  \le
  \eta a^2
  +
  \frac{1}{4\eta}b^2,
  \ 
  \forall a,b\in\R, \eta\in\R_{>0}
  .
$
%DIF >  \end{equation}
By applying Young's inequality to the third term in the right-hand side of \eqref{eq:Vk:dot:2}, we have
}\begin{align} \label{eq:Vk:dot:3} \nonumber
  \ddtt V_{k}
  \le &
    -(\underline{p}-\eta_1-\eta_2)
    \norm{\mv{e}}^2
    -
    \left(
      \sigma_k^o
      \!-\!
      \tfrac{1}{4\eta_2}
      \overline{p}^2
      (l_k+1)
    \right)
    \norm{\estwth_k}^2
  \\ 
  &
    -
    \sigma_k \overline{\theta}_k^2
    -
    \textstyle\sum_{j\in\mathcal{I}}
    \lambda_j
    \underline{c}_j^u
  \\ \nonumber
  &
      +
  \underbrace{
    \tfrac{1}{4\eta_1}
    \left[
    \overline{p}
      \left(
        \overline{\theta}_k^*
        \sqrt{l_{k}+1}
        +
        \overline{\epsilon}
      \right)
    \right]^2
    +
    \sigma_k^o
    \overline{\theta}_k^2
  }_{=: \delta_k \in \R_{>0}}
  ,
\end{align}\DIFaddend 
\DIFaddbegin \DIFadd{where $\eta_1,\eta_2\in\R_{>0}$ and $\sigma_k^o\in\R_{>0}$ are arbitrarily chosen such that $\underline{p}-\eta_1-\eta_2>0$ for given $\underline{p}$, and $\sigma_k^o - \frac{1}{4\eta_1} \overline{p}^2 (l_k+1) > 0$, respectively.
}\DIFaddend 

\DIFdelbegin \DIFdel{The satisfaction of the constraints, }%DIFDELCMD < \ie %%%
\DIFdel{$c_j,\ \forall j\in\mathcal{I}\setminus \{\theta_i\}_{i\in\{0,\cdots,k-1\}}$, can be verified through the boundedness of $\estwth_k$.
For the output layer's weight norm constraint $c_{\theta_k}$, if the corresponding Lagrange multiplier $\lambda_{\theta_k}$ increases sufficiently large due to the constraint violation, $\estwth_k$ approaches toward the origin as the invariant set $\Theta_{\theta_k}^{1}$ shrinks to a point.
The remaining input constraints can be validated implicitly through the term ${\gamma}_2$ in \eqref{eq:lya2:dot}.
Similarly to $c_{\theta_k}$, when an input constraint $c_j,\ \forall j\in\mathcal{I}\setminus \{\theta_i\}_{i\in\{0,\cdots,k\}}$, is violated, the corresponding Lagrange multiplier }\DIFdelend \DIFaddbegin \DIFadd{This result implies that $\ddtt V_{k}$ is negative outside the compact set defined by 
}\begin{equation} 
  \Omega_k
  \!
  =
  \!
  \left\{
    \!
    \mv{z}_k
    \!
    :
    \!
    a
    \!
    \norm{
      \left(
        \!
        \mv{e}^\top
        \!
        ,
        \!
        \estwth_k^\top
        \!
      \right)^\top
    }^2
    \!
    +
    \!
    b
    \!
    \norm{
      \left(
        \sigma_k
        % \!
        ,
        % \!
        [\lambda_j]_{j\in\mathcal{I}}^\top
      \right)^\top
    }_1
    \!
    \le
    \!
    \delta_k
    \!
  \right\}
  \!
  ,
\end{equation}
\DIFadd{where
$
    a
    :=
    \min
    \left\{
        \underline{p}
        - \eta_1
        - \eta_2
        ,
        \sigma_k^o
        -
        \tfrac{1}{4\eta_1}
        \overline{p}^2
        (l_k+1)
    \right\}
    \in\R_{>0}
$ 
and
$
    b
    :=
    \min
    \left\{
        \overline{\theta}_k^2
        ,
        \min_{j\in\mathcal{I}} 
        \left\{
            \underline{c}_j^u
        \right\}
    \right\}
    \in\R_{>0}
    .
$
%DIF >  and
%DIF >  $
%DIF >      \mv{z}_1 
%DIF >      := 
%DIF >      \left(
%DIF >          \mv{e}^\top
%DIF >          ,
%DIF >          \estwth_k^\top
%DIF >      \right)^\top
%DIF >      ,
%DIF >      \mv{z}_2
%DIF >      :=
%DIF >      \left(
%DIF >          \sigma_k
%DIF >          ,
%DIF >          [\lambda_j]_{j\in\mathcal{I}}^\top
%DIF >      \right)^\top
%DIF >      .
%DIF >  $
Moreover, since $V_k$ is lower bounded and radially unbounded, there exists a compact sublevel set
}\begin{equation} \label{eq:Sk:def}
  \mathcal S_k(\overline V_k)
  :=
  \left\{ 
    \mv{z}_k 
    : 
    V_k(\mv{z}_k)
    \le 
    \overline{V}_k 
  \right\}
  ,
\end{equation}
\DIFadd{where $\overline{V}_k\in\R_{>0}$ is arbitrarily chosen such that $\Omega_k \subset \mathcal{S}_k(\overline{V}_k)$.
Because $\ddtt V_k<0$ for all $\mv{z}_k\notin\Omega_k$, the system initialized in $\mathcal{S}_k(\overline{V}_k)$ will never escape $\mathcal{S}_k(\overline{V}_k)$ and eventually enter $\Omega_k$.
Therefore, UBB of $\mv e$, $\estwth_k$, $\sigma_k$, and }\DIFaddend $\lambda_j$\DIFdelbegin \DIFdel{increases sufficiently to a large value, leading to an increase in ${\gamma}_2$.
It makes ${\gamma}_2$ dominate the right-hand side of \eqref{eq:lya2:dot}, which makes $\ddtt V_2$ negative definite.
This drives $\estwth_k$ toward the origin until the constraint is satisfied.
%DIF <  By Assumption \ref{assum:convex} and Lemma \ref{lem:convex:angle}, the input constraints can be considered as convex constraint in the $\estwth_k$-space.
}\DIFdelend \DIFaddbegin \DIFadd{, $\forall j\in\mathcal I$, is ensured in $\mathcal{S}_k(\overline{V}_k)$.
}\DIFaddend 

\DIFdelbegin \DIFdel{Therefore}\DIFdelend \DIFaddbegin \DIFadd{In addition}\DIFaddend , the boundedness of \DIFdelbegin \DIFdel{$\fe$ and $\estwth_k$ is guaranteed by the invariant sets $\Theta_{\fe}^{1}$ }\DIFdelend \DIFaddbegin \DIFadd{$\mv{z}_k$ implies that $\mv{x}$ is bounded, since $\mv{e}=\mv{x}-\mv{x}_d$ and $\mv{x}_d$ is bounded by Assumption~\ref{assum:xd}.
Then, by defining the approximation region $\Omega_{\NNin}$ as $\mathcal{S}_k(\overline{V}_k)$, we can ensure that $\NNin$ remains in $\Omega_{\NNin}$.
This justifies the boundedness of $\idealwth_k$ }\DIFaddend and \DIFdelbegin \DIFdel{$\Theta_{\theta_k}^{1}$, and the satisfaction of $c_j,\ \forall j\in\mathcal{I}\setminus \{\theta_i\}_{i\in\{0,\cdots,k-1\}}$, is ensured}\DIFdelend \DIFaddbegin \DIFadd{$\mv{\epsilon}$ used in this analysis}\DIFaddend .

\DIFdelbegin \subsection*{\DIFdel{Case 2: Control input is not saturated.}}
%DIFAUXCMD
\DIFdelend \DIFaddbegin \hfill
\DIFaddend 

\DIFdelbegin \DIFdel{Since the input constraints are inactive, the saturation function $\mysat(\cdot)$ in \eqref{eq:lya1:dot2} and the input constraint functions $c_j,\ \forall j\in\mathcal I \setminus \{\theta_i\}_{i\in\{0,\cdots,k\}}$, in \eqref{eq:lya2:dot} can be neglected.
}\DIFdelend \DIFaddbegin \textit{\DIFadd{Step~2: Boundedness of the inner-layer adaptive variables.}}

\DIFadd{The remaining weights $[\estwth_i]_{i\in\mathcal{K}\setminus\{k\}}$ and the Lagrange multipliers $[\sigma_i]_{i\in\mathcal{K}\setminus\{k\}}$ are analyzed recursively from the $(k-1)$th layer to the input layer ($i=0$). 
Specifically, for an arbitrary layer index $i\in\{0,\cdots,k-1\}$, we assume that the subsequent adaptive variables, $\estwth_{i+1},\sigma_{i+1},\cdots,\estwth_k,\sigma_k,$ have already been shown to be bounded, and then prove the boundedness of $\estwth_i$ and $\sigma_i$.
}

\DIFaddend Consider the Lyapunov function \DIFdelbegin \DIFdel{candidate $V_3:=V_1+V_2$, whose time derivative is given by:
}%DIFDELCMD < \begin{equation}
%DIFDELCMD <     \begin{aligned}
%DIFDELCMD <         \ddtt V_3
%DIFDELCMD <         =&
%DIFDELCMD <         -\lambda_{\min}(\mm K)\norm{\fe}^2
%DIFDELCMD <         +\overline\tau_d\norm{\fe}
%DIFDELCMD <         +\fe^\top (\estNN - \idealNN - \mv\epsilon)
%DIFDELCMD <         \\
%DIFDELCMD <         &
%DIFDELCMD <         -\estwth_k^\top 
%DIFDELCMD <         \left(
%DIFDELCMD <             (\mm I_{l_{k+1}}\otimes \estact_k^\top)^\top
%DIFDELCMD <             \fe
%DIFDELCMD <             +
%DIFDELCMD <             \lambda_{\theta_k}\estwth_k
%DIFDELCMD <         \right)
%DIFDELCMD <         \\
%DIFDELCMD <         \le&
%DIFDELCMD <         -\lambda_{\min}(\mm K)\norm{\fe}^2
%DIFDELCMD <         +\fe^\top \estNN
%DIFDELCMD <         +
%DIFDELCMD <         \left(
%DIFDELCMD <             \overline\tau_d+\norm{
%DIFDELCMD <                 \idealNN+\mv\epsilon
%DIFDELCMD <             }
%DIFDELCMD <         \right)
%DIFDELCMD <         \norm{\fe}
%DIFDELCMD <         \\
%DIFDELCMD <         &
%DIFDELCMD <         -\estNN^\top\fe 
%DIFDELCMD <         -\lambda_{\theta_k}\norm{\estwth_k}^2
%DIFDELCMD <         \\
%DIFDELCMD <         \le&
%DIFDELCMD <         -
%DIFDELCMD <         \lambda_{\min}(\mm K)
%DIFDELCMD <         \norm{\fe}^2
%DIFDELCMD <         +
%DIFDELCMD <         \left(
%DIFDELCMD <             \overline\tau_d
%DIFDELCMD <             +
%DIFDELCMD <             \overline{\theta}^*_k\sqrt{l_{k}+1}
%DIFDELCMD <             +
%DIFDELCMD <             \overline{\epsilon}
%DIFDELCMD <         \right)
%DIFDELCMD <         \norm{\fe}
%DIFDELCMD <         \\
%DIFDELCMD <         &
%DIFDELCMD <         -
%DIFDELCMD <         \lambda_{\theta_k}\norm{\estwth_k}^2
%DIFDELCMD <         .
%DIFDELCMD <     \end{aligned}
%DIFDELCMD <     \label{eq:lya3:dot}
%DIFDELCMD < \end{equation}%%%
\DIFdelend \DIFaddbegin \DIFadd{of the $i$th layer defined as
}\begin{equation} \label{eq:Vi:def}
  V_{i}(\mv{z}_i)
  :=
  \tfrac{1}{2\alpha}\estwth_i^\top\estwth_i
  +
  \tfrac{1}{\beta_i^\theta}(\sigma_i-\sigma_i^o)^2
  ,
\end{equation}\DIFaddend 
\DIFdelbegin \DIFdel{Based on the same reasoning as
in Case 1, the boundedness of $\fe$ is ensured by the invariant set:
}%DIFDELCMD < \begin{equation}
%DIFDELCMD <     \Theta_{\fe}^{2}
%DIFDELCMD <     =
%DIFDELCMD <     \left\{ 
%DIFDELCMD <         \fe \in \R^n 
%DIFDELCMD <         \mid 
%DIFDELCMD <         \norm{\fe} 
%DIFDELCMD <         \le 
%DIFDELCMD <         \tfrac{
%DIFDELCMD <             \overline\tau_d
%DIFDELCMD <             +
%DIFDELCMD <             \overline{\theta}^*_k\sqrt{l_{k}+1}
%DIFDELCMD <             +
%DIFDELCMD <             \overline{\epsilon}
%DIFDELCMD <         }{
%DIFDELCMD <             \lambda_{\min}(\mm K)
%DIFDELCMD <         }
%DIFDELCMD <     \right\}
%DIFDELCMD <     .
%DIFDELCMD < \end{equation}
%DIFDELCMD < %%%
\DIFdel{The boundedness of $\estwth_k$ can also be ensured by the invariant set $\Theta_{\theta_k}^{2}=\{ \estwth_k \in \R^{\Xi_{k}} \mid \norm{\estwth_k} \le \overline\theta_k \}$, as $\lambda_{\theta_k}$ remains positive unless the weight norm constraint $c_{\theta_k}$ is satisfied.
}%DIFDELCMD < 

%DIFDELCMD < \hfill 
%DIFDELCMD < 

%DIFDELCMD < %%%
\DIFdel{The boundedness of }\DIFdelend \DIFaddbegin \DIFadd{where 
$
  \mv{z}_i:=
  \left(
    \estwth_i^\top,
    \sigma_i
  \right)^\top
  ,
$
and $\sigma_i^o\in\R_{>0}$ will be determined in }\DIFaddend the \DIFdelbegin \DIFdel{inner layer weights and the satisfaction of their corresponding constraints $c_j,\ \forall j\in\{\theta_i\}_{i\in\{k-1,\cdots,0\}}$, can be established recursively using \mbox{%DIFAUXCMD
\cite[Chap.~4, Thm.~1.9]{Desoer:2009aa}}\hskip0pt%DIFAUXCMD
.
The dynamics of $\estwth_i,\ \forall i\in\{k-1,\cdots,0\}$, are represented as
}%DIFDELCMD < \begin{equation}
%DIFDELCMD <     \ddtt \estwth_i 
%DIFDELCMD <     =
%DIFDELCMD <     -\alpha
%DIFDELCMD <     \left(
%DIFDELCMD <         {\pptfrac{\estNN}{\estwth_i}}^\top
%DIFDELCMD <         \fe
%DIFDELCMD <         +
%DIFDELCMD <         \lambda_{\theta_i}\estwth_i
%DIFDELCMD <         +
%DIFDELCMD <         \textstyle\sum_{
%DIFDELCMD <             j\in\mathcal{I}\setminus \{\theta_i\}_{i\in\{0,\cdots,k\}}
%DIFDELCMD <         }
%DIFDELCMD <         \lambda_j
%DIFDELCMD <         \pptfrac{c_j}{\estwth_i}
%DIFDELCMD <     \right)
%DIFDELCMD <     .
%DIFDELCMD <     \label{eq:estwth:dyn}
%DIFDELCMD < \end{equation}%%%
\DIFdelend \DIFaddbegin \DIFadd{subsequent analysis.
As in the analysis of $V_k$, the weight constraint $c_i^{\theta}(\estwth_i)$ is assumed to be violated, and the projection operator in \eqref{eq:adap:sig} is omitted according to Lemma~\ref{lem:proj:ineq}.
Taking the time derivative of $V_i$ yields
}\begin{align} \label{eq:Vi:dot:1} 
  \ddtt V_i
  \le&
    -\sigma_i \norm{\estwth_i}^2
    -
    \estwth_i^\top
    {\pptfrac{\estNN}{\estwth_i}}^\top
    \mm{P}
    \mv{e}
  \\ \nonumber
  &
    -
    \textstyle\sum_{j\in\mathcal{I}}
    \lambda_j 
    \estwth_i^\top
    {\pptfrac{\estNN}{\estwth_i}}^\top
    \mm{P}
    \pptfrac{c^{u}_{j}}{\mv{u}}
    +
    2(\sigma_i-\sigma_i^o) c_i^{\theta}(\estwth_i)
  \\ \nonumber
  \le&
    -\sigma_i \norm{\estwth_i}^2
    +
    (\sigma_i-\sigma_i^o) 
    \left(
      \norm{\estwth_i}^2
      -
      \overline{\theta}_i^2
    \right)
  \\ \nonumber
  &
    -
    \estwth_i^\top
    \underbrace{
    \left(
      {\pptfrac{\estNN}{\estwth_i}}^\top
      \mm{P}
      \left(
        \mv{e}
        +
        \textstyle\sum_{j\in\mathcal{I}}
        \lambda_j 
        \pptfrac{c^{u}_{j}}{\mv{u}}
      \right)
    \right)
    }_{=: \gamma_i\in\R^{\Xi_i}}
  \\ \nonumber
  \le&
    -\sigma_i^o \norm{\estwth_i}^2
    -\sigma_i \overline{\theta}_i^2
    +
    \norm{\estwth_i}
    \norm{\gamma_i}
    +
    \sigma_i^o \overline{\theta}_i^2
  .
\end{align}\DIFaddend 
\DIFdelbegin \DIFdel{By Lemma \ref{lem:cstr:grad:bound}, $\estwth_i$ }\DIFdelend \DIFaddbegin 

\DIFadd{For further analysis, the boundedness of $\gamma_i$ is analyzed.
According to the DNN gradient in \eqref{eq:gradient:phi}, $\pptfrac{\estNN}{\estwth_i}$ is bounded, since $\estwth_{i+1},\cdots,\estwth_k$ are bounded by assumption.
The remaining terms $\mv e$, $\estwth_k$, and $[\lambda_j]_{j\in\mathcal I}$ are bounded by the result of Step~1.
In addition, $\pptfrac{c_j^u}{\mv u}$ }\DIFaddend is bounded,\DIFdelbegin \DIFdel{provided that $\estwth_i,\ \forall i\in\{k,\cdots,i+1\}$, are bounded }\DIFdelend \DIFaddbegin \DIFadd{since $c_j^u(\cdot)$ is smooth by Assumption~\ref{assum:sat} and $\mv u$ is bounded by $\estwth_k$ in \eqref{eq:NN:bound}.
Therefore, $\gamma_i$ is bounded by an unknown positive constant $\overline{\gamma}_i\in\R_{>0}$.
}

\DIFadd{Then, by applying Young's inequality to the coupled term $\norm{\estwth_i}\norm{\gamma_i}\le\norm{\estwth_i}^2+\tfrac{1}{4}\overline{\gamma}_i^2$ in \eqref{eq:Vi:dot:1}, we have
}\begin{align} \label{eq:Vi:dot:2}
  \ddtt V_i
  \le&
    -
    \left(
      \sigma_i^o
      -
      1
    \right)
    \norm{\estwth_i}^2
    -
    \sigma_i \overline{\theta}_i^2
    +
    \delta_i
  ,
\end{align}
\DIFadd{where $\delta_i:=\frac{1}{4}\overline{\gamma}_i^2 + \sigma_i^o \overline{\theta}_i^2$ is a positive constant, and $\sigma_i^o\in\R_{>0}$ is arbitrarily chosen such that $\sigma_i^o - 1 > 0$.
As in Step~1, $\ddtt V_i$ is negative outside the compact set defined by 
}\begin{equation}
  \Omega_i
  =
  \left\{
    \mv{z}_i
    :
    \left(
      \sigma_i^o
      -
      \tfrac{\overline{\gamma}_i^2}{2}
    \right)
    \norm{\estwth_i}^2
    +
    \overline{\theta}_i^2
    \norm{\sigma_i}_1
    \le
    \delta_i
  \right\}
  , 
\end{equation}
\DIFadd{which implies the boundedness of $\estwth_i$ and $\sigma_i$}\DIFaddend .
\DIFdelbegin \DIFdel{This holds because the system matrix $-\lambda_{\theta_i} \mm I_{\Xi_i}$ is stable, and the residual terms—such as $\pptfrac{\estNN}{\estwth_i}, \fe$, and $\pptfrac{c_j}{\estwth_i}$ —are bounded.
}\DIFdelend Moreover, \DIFdelbegin \DIFdel{the Lagrange multiplier $\lambda_j$ does not diverge since the input constraints $c_j,\ \forall j\in\mathcal{I}\setminus \{\theta_i\}_{i\in\{0,\cdots,k\}}$ are satisfied before divergence can occur.
Therefore, starting from the $(k-1)$th layer}\DIFdelend \DIFaddbegin \DIFadd{a compact sublevel set $\mathcal{S}_i(\overline{V}_i)$ of $V_i$ can be constructed such that 
}\begin{equation} \label{eq:Si:def}
  \mathcal{S}_i(\overline{V}_i)
  :=
  \left\{
    \mv{z}_i
    :
    V_i(\mv{z}_i)
    \le
    \overline{V}_i
  \right\}
  ,
\end{equation}
\DIFadd{where $\overline{V}_i\in\R_{>0}$ is arbitrarily chosen such that $\Omega_i \subset \mathcal{S}_i(\overline{V}_i)$.
Since $\ddtt V_i<0$ for all $\mv{z}_i\notin\Omega_i$, the system initialized in $\mathcal{S}_i(\overline{V}_i)$ will remain in $\mathcal{S}_i(\overline{V}_i)$.
}

\DIFadd{By recursively applying the above analysis from $i=k-1$ to $i=0$}\DIFaddend , the boundedness of $\estwth_i$ \DIFaddbegin \DIFadd{and $\sigma_i$ for all $i\in\mathcal{K}\setminus\{k\}$ }\DIFaddend can be established\DIFdelbegin \DIFdel{recursively down to the input layer ($i=0$).
Moreover the satisfaction of the weight norm constraints for the inner layers can be verified, as the Lagrange multipliers $\lambda_{\theta_i}$, $\forall i\in\{k-1,\cdots,0\}$, increase sufficiently large to stabilize the $\estwth_i$ dynamics in \eqref{eq:estwth:dyn}, leading to its convergence toward the origin.
  }\DIFdelend \DIFaddbegin \DIFadd{, which concludes the proof.
}\DIFaddend 

\DIFdelbegin %DIFDELCMD < \hfill
%DIFDELCMD < %%%
\DIFdelend \DIFaddbegin \end{proof}
\DIFaddend 

\DIFdelbegin \DIFdel{In conclusion, the filtered tracking error $\fe$ and the weight estimate $\estwth$ are bounded, and all imposed constraints $c_j,\ \forall j\in\mathcal I$, are satisfied.
}\DIFdelend %DIF >  ——————————————————————————————————————————————
%DIF >  REMARKS
%DIF >  ——————————————————————————————————————————————

\DIFdelbegin %DIFDELCMD < \end{proof}
%DIFDELCMD < %%%
\DIFdelend \DIFaddbegin \DIFadd{Additionally, we introduce the following remarks regarding the initialization and the parameter tuning of CONAC.
}\DIFaddend 

%DIF <   SECTION SIMULATION ======================================
\DIFdelbegin \section{\DIFdel{Real-Time Implementation and Validation}}%DIFAUXCMD
\addtocounter{section}{-1}%DIFAUXCMD
%DIFDELCMD < \label{sec:sim}
%DIFDELCMD < %%%
\DIFdelend \DIFaddbegin \begin{remark}[Initialization of CONAC] \label{rem:init}
  \DIFadd{The invariant sublevel set 
  $
    \bigcap_{i\in\mathcal{K}} 
    \mathcal{S}_i(\overline{V}_i)
  $ 
  used in the stability analysis depends on unknown constants.
  Therefore, the exact admissible initial domain cannot be computed explicitly.
  In practice, we recommend initializing; the system state $\mv{x}(0)$ close to the initial reference state $\mv{x}_d(0)$ to reduce the initial tracking error $\mv{e}(0)$; DNN weights with small random values; and zero for the Lagrange multipliers, $[\sigma_i(0)]_{i\in\mathcal{K}}=\mv{0}$ and $[\lambda_j(0)]_{j\in\mathcal{I}}=\mv{0}$ without constraint violations, to help ensure that the system starts within
  $
    \bigcap_{i\in\mathcal{K}} 
    \mathcal{S}_i(\overline{V}_i)
    .
  $
}\end{remark}
\DIFaddend 

\DIFdelbegin \subsection{\DIFdel{Validation Setup}}
%DIFAUXCMD
\addtocounter{subsection}{-1}%DIFAUXCMD
\DIFdelend \DIFaddbegin \begin{remark}[Initial Desired Trajectory Design] \label{rem:ref}  
  \DIFadd{The initialization in Remark~\ref{rem:init} helps prevent initial unstable behavior, which generally results from poor initial performance of NAC due to the randomly initialized DNN weights.
  With small random weights, the initial control input is expected to be small, thereby avoiding abrupt control actions.
  However, if the system starts near an unstable or highly nonlinear operating region, small random weights may result in unstable behavior or slow adaptation.
  Hence, it is preferable to design the desired trajectory such that the system starts in a sufficiently stable operating region, and use a conservative desired trajectory during an initial warm-up phase, as introduced in \mbox{%DIFAUXCMD
\cite{Ryu:2025aa}}\hskip0pt%DIFAUXCMD
.
  %DIF >  This strategy is demonstrated across valuations in Section~\ref{sec:val}.
}\end{remark}
\DIFaddend 

\DIFaddbegin \begin{remark}[Tuning of CONAC Parameters] \label{rem:tuning}
  \DIFadd{The adaptation gain $\alpha$ and the multiplier update rates $\beta_i^\theta$ and $\beta_j^u$ affect the transient behavior of CONAC.
  The gain $\alpha$ determines the overall update speed of the DNN weights.
  A large $\alpha$ may improve the initial learning response but can also induce overshooting or oscillatory behavior, resulting in failure to learn the ideal weights.
  Conversely, an excessively small $\alpha$ may lead to slow adaptation and poor transient tracking performance.
  The update rates $\beta_i^\theta$ and $\beta_j^u$ determine how rapidly the Lagrange multipliers respond to weight and input constraint violations, respectively.
  Larger values accelerate the enforcement of the corresponding constraints by increasing the multiplier response when a violation occurs.
  However, excessively large update rates may introduce abrupt changes in the adaptation law, leading to chattering or high-frequency oscillations in the control input on the constraint boundary.
}

  \DIFadd{Therefore, $\alpha$, $\beta_i^\theta$, and $\beta_j^u$ should be selected to balance learning speed, prompt constraint enforcement, and smooth closed-loop behavior.
  The effect of the input constraint multiplier update rate $\beta_j^u$ is evaluated in the parametric study in Section~\ref{sec:val:results}.
}\end{remark}

\DIFaddend % ------------------------------------
%DIF <  Real-time Implementation Set-up
%DIF >        SECTION VALIDATIONS
% ------------------------------------
\DIFdelbegin %DIFDELCMD < \begin{figure}[t]
%DIFDELCMD <     \centering
%DIFDELCMD <         \includegraphics[width=.99\linewidth]
%DIFDELCMD <         %%%
%DIF <  \includegraphics[width=\figSizeOneCol\linewidth]
        %DIFDELCMD < {
%DIFDELCMD <             %%%
\DIFdelFL{
  figs/expSet.drawio.png
        }%DIFDELCMD < }%%%
%DIF < 
    %DIFDELCMD < \caption{%
{%DIFAUXCMD
\DIFdelFL{Experimental setup of the two-link robotic manipulator used for real-time validation.
    }}
    %DIFAUXCMD
%DIFDELCMD < \label{fig:ctrl:exp:set}
%DIFDELCMD < \end{figure}
%DIFDELCMD < %%%
\DIFdelend \DIFaddbegin \section{\DIFadd{Validations}} \label{sec:val}
\DIFaddend 

\DIFaddbegin \subsection{\DIFadd{Validation Setup}} \label{sec:val:setup}

%DIF >  -------------------------------------
%DIF >  Experiment Setup
%DIF >  -------------------------------------
\DIFaddend To validate the \DIFdelbegin \DIFdel{effectiveness of the }\DIFdelend proposed CONAC, a real-time experiment was conducted on a two-link robotic manipulator, as shown in Fig.~\DIFdelbegin \DIFdel{\ref{fig:ctrl:exp:set}}\DIFdelend \DIFaddbegin \DIFadd{\ref{fig:exp:real:set}}\DIFaddend . 
The OpenCR1.0 board \cite{opencr} with a 32-bit ARM Cortex \DIFdelbegin \DIFdel{and a 216 MHz }\DIFdelend \DIFaddbegin \DIFadd{processor and a $\qty{216}{\mega\hertz}$ }\DIFaddend clock frequency was used \DIFdelbegin \DIFdel{for the real-time implementation}\DIFdelend \DIFaddbegin \DIFadd{as the microcontroller unit (MCU)}\DIFaddend .
The two-link robotic manipulator was actuated by two Cubemars AK10-90 servo motors.
The servos were controlled using the \DIFdelbegin \DIFdel{OpenCR1.0 board}\DIFdelend \DIFaddbegin \DIFadd{MCU}\DIFaddend , which transmits the torque reference to the corresponding servo via controller area network (CAN) communication \DIFaddbegin \DIFadd{at a sampling rate of $\qty{500}{\hertz}$ using a PCAN-USB interface}\DIFaddend .
The system \DIFdelbegin \DIFdel{dynamics of the }\DIFdelend \DIFaddbegin \DIFadd{parameters and the schematic diagram are shown in Table~\ref{table:exp:system:params} and Fig.~\ref{fig:exp:robot:diagram}, respectively.
}

%DIF >  -------------------------------------
%DIF >  Robot System
%DIF >  -------------------------------------
\DIFadd{The }\DIFaddend two-link robotic manipulator \DIFdelbegin \DIFdel{can follow the }\DIFdelend \DIFaddbegin \DIFadd{dynamics are described by the following }\DIFaddend Euler-Lagrange equation\DIFdelbegin \DIFdel{\eqref{eq:sys1}, where $\q$ denotes the joint angles and $\rbu$ denotes the control input (joint torques), }\DIFdelend \DIFaddbegin \DIFadd{:
}\begin{equation} \label{eq:robot:dynamics}
  \mv{M}(\mv{q})\tfrac{\der^2}{\der t^2}{\mv{q}}
  +
  \mv{C}(\mv{q},\ddtt\mv{q})\ddtt\mv{q}
  +
  \mv{G}(\mv{q})
  =
  \mysat
  (
    \mv{\tau}
  )
  ,
\end{equation}
\DIFadd{where $\mv{q} := \left(q_1,q_2\right)^\top \in \R^2$ is the joint angle vector, $\mv{\tau} := \left(\tau_1,\tau_2\right)^\top \in \R^2$ is the joint torque vector, $\mv{M}(\mv{q}) \in \R^{2\times2}$ is the inertia matrix, $\mv{C}(\mv{q},\ddtt\mv{q}) \in \R^{2\times2}$ is the Coriolis matrix, and $\mv{G}(\mv{q}) \in \R^2$ is the gravity vector.
The inertia matrix $\mv{M}(\mv{q})$ whose inverse is the control effectiveness matrix, is a symmetric positive-definite matrix according to the physical properties of Euler-Lagrange equation, \mbox{%DIFAUXCMD
\cite[Lemma 3.1]{Ji:2025aa}}\hskip0pt%DIFAUXCMD
.
}

\DIFadd{To apply CONAC, the system was reformulated into the first-order dynamics form in \eqref{eq:sys} by defining the state as $\mv{x} := \ddtt\mv{q}+\Lambda\mv{q}$, and the control input as $\mv{u} := \mv{\tau}$, where $\Lambda\in\R^{2\times2}$ is a positive definite matrix, and $\mv{f}:=\mv{M}^{-1}(\mv{q})\left(-\mv{C}(\mv{q},\ddtt\mv{q})\ddtt\mv{q}-\mv{G}(\mv{q})\right)+\Lambda\ddtt\mv{q}$, and $\mm{B}:=\mv{M}^{-1}(\mv{q})$.
Moreover, the DNN input vector was defined as 
$
  \NNin 
  := 
  \left(
    (\mv{q}-\mv{q}_d)^\top
    ,
    \ddtt(\mv{q}-\mv{q}_d)^\top
    ,
    \tfrac{\der^2}{\der t^2}\mv{q}_d^\top
    ,
    1
  \right)^\top
  ,
$
since $\mv{f}$, $\mm{B}$ }\DIFaddend and \DIFdelbegin \DIFdel{is }\DIFdelend \DIFaddbegin \DIFadd{$\ddtt\mv{x}_d$ in \eqref{eq:ctrl:NN:actual} are functions of $\mv{q},\mv{q}_d,\ddtt\mv{q},\ddtt\mv{q}_d$ and $\tfrac{\der^2}{\der t^2}\mv{q}_d$.
}

\DIFadd{The input saturation $\mysat(\mv{\tau})$ was defined as a combination of the torque saturation of the second joint and the norm saturation of the control input, as follows:
}\begin{equation} \label{eq:exp:sat}
  \mysat(\mv{\tau})
  =
  \begin{cases}
    % if |tau_2| > \overline{\tau}_2 and |tau_1| > \overline{\tau}_1
    \begin{pmatrix}
      \mysign(\tau_1)\overline{\tau}_1
      \\
      \mysign(\tau_2)\overline{\tau}_2
    \end{pmatrix}
    , &
    \text{if }
    \vert\tau_2\vert > \overline{\tau}_2
    \text{ and }
    \vert\tau_1\vert > \overline{\tau}_1
    \\
    % if |tau_2| > \overline{\tau}_2 and |tau_1| <= \overline{\tau}_1
    \left(
      \tau_1
      , 
      \mysign(\tau_2)\overline{\tau}_2
    \right)^\top
    , &
    \text{if }
    \vert\tau_2\vert > \overline{\tau}_2
    \text{ and }
    \vert\tau_1\vert \le \overline{\tau}_1
    \\
    % if |tau_2| <= \overline{\tau}_2 and ||tau|| > \overline{\tau}
    \tfrac{\overline{\tau}}{\norm{\mv{\tau}}}\mv{\tau}
    , &
    \text{if }
    |\tau_2| \le \overline{\tau}_2
    \text{ and }
    \norm{\mv{\tau}} > \overline{\tau}
    \\
    % otherwise
    \mv{\tau}
    , &
    \text{otherwise}
    ,
  \end{cases}
  % \mysat(\mv{\tau})
  % =
  % \min\left\{
  %   1,\,
  %   \frac{\overline{\tau}}{\norm{\mv{v}}}
  % \right\}
  % \mv{v}
  % ,
  % % \text{where } 
  % \ 
  % \mv{v}
  % :=
  % \begin{pmatrix}
  %   \tau_1
  %   \\
  %   \operatorname{clip}(\tau_2,-\overline{\tau}_2,\overline{\tau}_2)
  % \end{pmatrix}
  % ,
\end{equation}
\DIFadd{where $\overline{\tau}_2=\qty{3.8}{\newton\meter}$, $\overline{\tau}_1=\qty{10.32}{\newton\meter}$ and $\overline{\tau}=\qty{11}{\newton\meter}$, as }\DIFaddend illustrated in Fig.~\DIFdelbegin \DIFdel{\ref{fig:robot:model}.
The measured and /or estimated system parameters are provided in Table~\ref{table:system:params}}\DIFdelend \DIFaddbegin \DIFadd{\ref{fig:exp:robot:sat}.
The saturation function can be decomposed into three convex input constraints, which can be physically interpreted as the upper and lower torque limit of the second joint and the power source limit, respectively, and are denoted by
}\begin{equation} \label{eq:val:cstr:input}
  c^u_{\overline{\tau}_2}(\mv{\tau}) 
  := 
  \tau_2 - \overline{\tau}_2
  ,
  \ 
  c^u_{\underline{\tau}_2}(\mv{\tau}) 
  := 
  -\tau_2 - \overline{\tau}_2
  ,
  \ 
  c^u_{\overline{\tau}}(\mv{\tau})
  := 
  \norm{\mv{\tau}}^2 - \overline{\tau}^2
  .
\end{equation}
\DIFadd{Accordingly, the input constraint index set is defined as $\mathcal{I} = \{\overline{\tau}_2, \underline{\tau}_2, \overline{\tau}\}$}\DIFaddend .

%DIF <  ------------------------------------
%DIF <  Robot Model Figure and Table
%DIF <  ------------------------------------
\begin{figure}[t]
\centering
  \DIFdelbeginFL %DIFDELCMD < \subfloat[Two-link robotic manipulator.]{
%DIFDELCMD <         \includegraphics[width=0.499\linewidth]{src/figures/RobotModel.drawio.pdf}
%DIFDELCMD <         \label{fig:robot:model}
%DIFDELCMD <     }
%DIFDELCMD <     %%%
%DIF <  \hfill
    %DIFDELCMD < \subfloat[Control input saturation function.]{
%DIFDELCMD <         \includegraphics[width=0.499\linewidth]{src/figures/saturation.drawio.pdf}
%DIFDELCMD <         \label{fig:robot:sat}
%DIFDELCMD <     }
%DIFDELCMD <     %%%
\DIFdelendFL \DIFaddbeginFL \includegraphics[width=.99\linewidth]
  %DIF >  \includegraphics[width=\figSizeOneCol\linewidth]
  {
    % \DIFaddFL{figs/expSet.drawio.png}
    figs/expSet.drawio.png
  }%DIF > 
\DIFaddendFL \caption{\DIFdelbeginFL \DIFdelFL{Two-link }\DIFdelendFL 
    \DIFaddbeginFL \DIFaddFL{Experimental setup of the two-link }\DIFaddendFL robotic manipulator \DIFdelbeginFL \DIFdelFL{and control saturation function}\DIFdelendFL \DIFaddbeginFL \DIFaddFL{used for real-time validation}\DIFaddendFL .
  }
\DIFdelbeginFL %DIFDELCMD < \label{fig:robot}
%DIFDELCMD < %%%
\DIFdelendFL \DIFaddbeginFL \label{fig:exp:real:set}
\DIFaddendFL \end{figure}

%DIF <  ------------------------------------
%DIF <  Robot Model Parameters
%DIF <  -----------------------------------
\DIFaddbegin \begin{figure}[t]
  \centering
  \subfloat[Schematic diagram of the two-link robotic manipulator.]{
    \includegraphics[width=0.46\linewidth]{figs/RobotModel.drawio.pdf}
    \label{fig:exp:robot:diagram}
  }
  %DIF >  \vfill
  \subfloat[Considered input Saturation.]{
    \includegraphics[width=0.48\linewidth]{figs/saturation.drawio.pdf}
    \label{fig:exp:robot:sat}
  }
  \caption{
    \DIFaddFL{The two-link robotic manipulator and the considered input saturation.
  }}
  \label{fig:exp:robot}
\end{figure}

\DIFaddend \begin{table}[t]
  \renewcommand{\arraystretch}{1.3}
  \caption{Two-link robotic manipulator's parameters.}
  \centering
  \DIFdelbeginFL %DIFDELCMD < \begin{tabular}{c m{9.5em} c c c }
%DIFDELCMD <     %%%
%DIF <  \begin{tabular}{c c c c c}
    \DIFdelendFL \DIFaddbeginFL \begin{tabular}{c m{11.5em} c}
  \DIFaddendFL \hline
  \textbf{Symbol} & \textbf{Description} & \textbf{Value} \\
  \hline
  \hline 
    $m_1,m_2$ & Mass \DIFaddbeginFL \DIFaddFL{of the links }\DIFaddendFL & \DIFdelbeginFL \DIFdelFL{2.465 $\mathrm{kg}$ }\DIFdelendFL \DIFaddbeginFL \qty{2.465}{\kilo\gram} \DIFaddendFL \\
  \hline
    $l_1,l_2$  & Length \DIFaddbeginFL \DIFaddFL{of the links }\DIFaddendFL & \DIFdelbeginFL \DIFdelFL{0.2 $\mathrm{m}$ }\DIFdelendFL \DIFaddbeginFL \qty{0.2}{\meter} \DIFaddendFL \\  
  \hline
    ${l_c}_1,{l_c}_2$ & \DIFdelbeginFL \DIFdelFL{Center }\DIFdelendFL \DIFaddbeginFL \DIFaddFL{Length to the center }\DIFaddendFL of mass & \DIFdelbeginFL \DIFdelFL{0.139 $\mathrm{m}$ }\DIFdelendFL \DIFaddbeginFL \qty{0.139}{\meter} \DIFaddendFL \\
  \hline
    $I_1,I_2$  & Inertia \DIFaddbeginFL \DIFaddFL{of the links }\DIFaddendFL & \DIFdelbeginFL \DIFdelFL{0.069 $\mathrm{kg\cdot m^{2}}$ }\DIFdelendFL \DIFaddbeginFL \qty{0.069}{\kilo\gram\meter\squared} \DIFaddendFL \\
  \hline
  \end{tabular}
  \DIFdelbeginFL %DIFDELCMD < \label{table:system:params}
%DIFDELCMD < %%%
\DIFdelendFL \DIFaddbeginFL \label{table:exp:system:params}
\DIFaddendFL \end{table}

\DIFaddbegin \begin{figure}[t]
  \centering
  \includegraphics[width=0.7\linewidth]{figs/desired_traj.drawio.pdf}
  \caption{
    \DIFaddFL{The desired trajectory in Episode~1 and Episode~2.
  }}
  \label{fig:exp:des:traj}
\end{figure}

\DIFaddend % ------------------------------------
% desired trajectory
% ------------------------------------
The desired trajectory \DIFdelbegin \DIFdel{is designed with $4$ segments as follows:
}%DIFDELCMD < \begin{equation}
%DIFDELCMD <     \qd(t) 
%DIFDELCMD <     =
%DIFDELCMD <     \begin{pmatrix}
%DIFDELCMD <         {q_d}_1
%DIFDELCMD <         \\
%DIFDELCMD <         {q_d}_2
%DIFDELCMD <     \end{pmatrix}
%DIFDELCMD <     =
%DIFDELCMD <     \begin{cases}
%DIFDELCMD <     \mathcal{P}^5(t;\mv{q}_d^{(0)},\mv{q}_d^{(1)},t_f),
%DIFDELCMD <         &
%DIFDELCMD <         \text{if } 0 \le t < t_f 
%DIFDELCMD <         ,
%DIFDELCMD <         \\
%DIFDELCMD <     \mathcal{P}^5(t;\mv{q}_d^{(1)},\mv{q}_d^{(2)},t_f),
%DIFDELCMD <         &
%DIFDELCMD <         \text{if } t_f \le t < 2t_f 
%DIFDELCMD <         ,
%DIFDELCMD <         \\
%DIFDELCMD <     \mathcal{P}^5(t;\mv{q}_d^{(2)},\mv{q}_d^{(1)},t_f),
%DIFDELCMD <         & 
%DIFDELCMD <         \text{if } 2t_f \le t < 3t_f 
%DIFDELCMD <         ,
%DIFDELCMD <         \\
%DIFDELCMD <     \mathcal{P}^5(t;\mv{q}_d^{(1)},\mv{q}_d^{(0)},t_f),
%DIFDELCMD <         &    
%DIFDELCMD <         \text{if } 3t_f \le t < 4t_f 
%DIFDELCMD <         ,
%DIFDELCMD <         \\
%DIFDELCMD <     \end{cases} 
%DIFDELCMD <     \label{eq:desired trajectory}
%DIFDELCMD < \end{equation}
%DIFDELCMD < %%%
\DIFdel{where $\mathcal{P}^5(t;\mv{x}_1,\mv{x}_2,t_f)$ denotes }\DIFdelend \DIFaddbegin \DIFadd{comprised 4 segments that move the system from 
$
  \mv{q}_d^{(0)}
  :=
  \left(
    -\qty{60}{\degree}
    ,
    \qty{60}{\degree}
  \right)^\top    
$
to 
$
  \mv{q}_d^{(1)}
  :=
  \left(
    \qty{45}{\degree}
    ,
    -\qty{90}{\degree}
  \right)
  ,
$ 
then to 
$
  \mv{q}_d^{(2)}
  :=
  \left(
    -\qty{45}{\degree}
    ,
    \qty{45}{\degree}
  \right)^\top
  ,
$
back to
$
  \mv{q}_d^{(3)}
  :=
  \mv{q}_d^{(1)}
  ,
$
and finally to the initial position $\mv{q}_d^{(0)}$, as illustrated in Fig.~\ref{fig:exp:des:traj}.
Each segment was designed using }\DIFaddend a quintic ($5$th-order) polynomial trajectory \DIFdelbegin \DIFdel{that moves the system from $\mv{x}_1\in\R^2$ to $\mv{x}_2\in\R^2$ in $t_f$ seconds}\DIFdelend \DIFaddbegin \DIFadd{to ensure the smoothness of the desired trajectory}\DIFaddend , with zero initial and final velocities.
The \DIFdelbegin \DIFdel{initial and final joint angles of the segments are defined as follows:
$
    \mv{q}_d^{(0)}=\left(-\tfrac{1}{3}\pi,\tfrac{1}{3}\pi\right)^\top
$, 
$
    \mv{q}_d^{(1)}=\left(\tfrac{1}{4}\pi,-\tfrac{1}{2}\pi\right)^\top
$, 
$
    \mv{q}_d^{(2)}=\left(-\tfrac{1}{4}\pi,\tfrac{1}{4}\pi\right)^\top
$, as illustrated in Fig.~\ref{fig:robot:model}.
The }\DIFdelend time duration of the segments \DIFdelbegin \DIFdel{is defined as $t_f=3$ s}\DIFdelend \DIFaddbegin \DIFadd{was defined as $\qty{2}{\second}$}\DIFaddend .
To demonstrate the learning capability of the DNN, the desired trajectory was applied twice\DIFdelbegin \DIFdel{. 
The first and second repetitions of the desired trajectory are }\DIFdelend \DIFaddbegin \DIFadd{, }\DIFaddend referred to as Episode $1$ and Episode $2$, respectively.
%DIF <  ------------------------------------

\DIFdelbegin \DIFdel{However, due to the random initialization of the DNN weights, CONAC's initial performance is expected to be poor, leading to unstable behavior.
Therefore, to }\DIFdelend \DIFaddbegin \DIFadd{As discussed in Remark~\ref{rem:ref}, a warm-up phase was introduced prior to applying the desired trajectory to }\DIFaddend avoid the instability caused by the initial poor performance of \DIFdelbegin \DIFdel{CONAC, a warm-up phase was introduced prior to applying the desired trajectory}\DIFdelend \DIFaddbegin \DIFadd{the DNN}\DIFaddend .
In the warm-up phase, the manipulator was initially positioned at a stable equilibrium point, denoted as 
\DIFdelbegin \DIFdel{${\qd}^{\rm{origin}} = \left(-\tfrac{1}{2}\pi, 0\right)^\top$}\DIFdelend \DIFaddbegin \DIFadd{$
    {\qd}^{\rm{origin}} 
    := 
    \left(
        % -\tfrac{1}{2}\pi
        -\qty{90}{\degree}
        , 
        % 0
        \qty{0}{\degree}
    \right)^\top
    ,
$
and remained at this position for $\qty{3}{\second}$}\DIFaddend .
Subsequently, a \DIFdelbegin \DIFdel{trajectory $\mathcal{P}^5(t; {\qd}^{\rm{origin}}, \mv{q}_d^{(0)},3)$ was used to move the manipulator }\DIFdelend \DIFaddbegin \DIFadd{conservative desired trajectory was applied to gently excite the system, which was designed as a quintic polynomial trajectory that moves the system }\DIFaddend from ${\qd}^{\rm{origin}}$ to $\mv{q}_d^{(0)}$ \DIFdelbegin \DIFdel{over a duration of 3 seconds.
Once the manipulator reached $\mv{q}_d^{(0)}$, a constant desired trajectory $\qd = \mv{q}_d^{(0)}$ was maintained for 8 seconds to allow the system to settle and eliminate the influence of transient dynamics}\DIFdelend \DIFaddbegin \DIFadd{in $\qty{3}{\second}$, followed by a constant reference that holds the system at $\mv q_d^{(0)}$ for the next $\qty{8}{\second}$}\DIFaddend .

%DIF <  ------------------------------------
%DIF <  Control Input Saturation
%DIF <  ------------------------------------
\DIFdelbegin \DIFdel{In the validation, we assume that the actuators are subject to a convex saturation function that projects the control input onto the convex set 
$
    \Theta_{\mysat} :=\left\{ \Theta_{\mysat, {\tau_2}} \cap \Theta_{\mysat, {\rbu}} \right\}
$, as illustrated in Fig.~\ref{fig:robot:sat}, where
}%DIFDELCMD < \begin{subequations}
%DIFDELCMD <     \begin{align}
%DIFDELCMD <         \Theta_{\mysat, {\tau_2}}
%DIFDELCMD <         :=&
%DIFDELCMD <         \left\{
%DIFDELCMD <             \tau_2
%DIFDELCMD <             \mid
%DIFDELCMD <             \vert{\tau_2}\vert \le 
%DIFDELCMD <             3.5
%DIFDELCMD <             % \overline\tau_2
%DIFDELCMD <         \right\}
%DIFDELCMD <         ,
%DIFDELCMD <         \label{eq:sat:tau2}
%DIFDELCMD <         \\    
%DIFDELCMD <         \Theta_{\mysat, {\rbu}}
%DIFDELCMD <         :=&
%DIFDELCMD <         \left\{
%DIFDELCMD <             \rbu
%DIFDELCMD <             \mid
%DIFDELCMD <             \norm{\rbu} \le 
%DIFDELCMD <             11
%DIFDELCMD <             % \overline\tau
%DIFDELCMD <         \right\}
%DIFDELCMD <         .
%DIFDELCMD <         \label{eq:sat:norm}
%DIFDELCMD <     \end{align}
%DIFDELCMD <     \label{eq:sat}%%%
%DIF <  %%
%DIFDELCMD < \end{subequations}
%DIFDELCMD < %%%
\DIFdel{The sets defined in \eqref{eq:sat:tau2} and \eqref{eq:sat:norm} can be physically interpreted as the current limit imposed by the torque limit of the second joint actuator and the power source, respectively.
The corresponding input constraints are mathematically described in Appendix~\ref{sec:appen:cstr} as the input bound constraint \eqref{eq:cstr:input:bound} and the input norm constraint \eqref{eq:cstr:input:norm}, respectively}\DIFdelend \DIFaddbegin \hfill

%DIF >  -------------------------------------
%DIF >  Comparative Controllers
%DIF >  -------------------------------------

%DIF >  왜 optimization-based가 배제되었나
%DIF >  - 우리는 structural methods를 기반으로힘
%DIF >  - point-wise constraint handling이 아닌 adaptation 

\DIFadd{For a comparative analysis, four controllers were implemented, with their properties summarized in Table~\ref{table:controller}.
In the comparative study, optimization-based methods, such as MPC and CBF-QP, were not included because the focus is on NAC methods that handle input constraints.
This choice isolates the proposed contribution, namely the integration of multiple convex input constraints into the online adaptation laws of NAC while retaining a Lyapunov-based stability analysis.
In addition, among the structural methods, the auxiliary system-based NAC was included as a comparative controller, since it has better input constraint handling capability than smooth approximation methods \mbox{%DIFAUXCMD
\cite{Min:2021aa}}\hskip0pt%DIFAUXCMD
}\DIFaddend .

%DIF <  ------------------------------------
%DIF <  Controller Table
%DIF <  ------------------------------------
\begin{table}[t]
  \renewcommand{\arraystretch}{1.3}
  \caption{\DIFdelbeginFL \DIFdelFL{Properties of the controllers}\DIFdelendFL \DIFaddbeginFL \DIFaddFL{Controllers and parameters used for comparative analysis.}\DIFaddendFL }
  \label{table:controller}
  \centering
  \DIFdelbeginFL %DIFDELCMD < \begin{tabular}{m{.25cm} m{4cm} l}
%DIFDELCMD <     %%%
\DIFdelendFL \DIFaddbeginFL \begin{tabular}{m{.6cm} m{4cm} m{2.7cm}}
  \DIFaddendFL % \begin{tabular}{c l l}
  \hline
    &\DIFdelbeginFL %DIFDELCMD < \bf{Description}%%%
\DIFdelendFL \DIFaddbeginFL \bf{Description and parameters}\DIFaddendFL &\DIFdelbeginFL %DIFDELCMD < \bf{Input Constraint Handling}%%%
\DIFdelendFL \DIFaddbeginFL \bf{Input constraint}\DIFaddendFL \\
  \hline
  \hline
    \DIFdelbeginFL %DIFDELCMD < \multirow{3}{*}{(C$_1$)}%%%
\DIFdelendFL \DIFaddbeginFL \multirow{4}{*}{
      \shortstack{
        {(C$_1$)}
        \\
        \protect\colorLine{blue}{solid}
      }
    }\DIFaddendFL & \DIFdelbeginFL \DIFdelFL{CONAC with }\DIFdelendFL \DIFaddbeginFL \DIFaddFL{Proposed CONAC  
    }\DIFaddendFL & 
      \DIFdelbeginFL \DIFdelFL{\eqref{eq:cstr:input:bound} $c_{\overline\tau_2}$ and $c_{\underline\tau_2}$ }\DIFdelendFL \DIFaddbeginFL \multirow{4}{*}{
        \shortstack[l]{
          $-\overline{\tau}_2 \le \tau_2 \le \overline{\tau}_2$, 
          \\
            $\norm{\mv{\tau}} \le \overline{\tau}$,
          \\
            ($\overline{\tau}_2=3.8, \overline{\tau}=11$)
        }
      }
    \DIFaddendFL \\
    & \DIFdelbeginFL \DIFdelFL{large update rate $\beta_j$. }\DIFdelendFL \DIFaddbeginFL \DIFaddFL{$\beta^{u}_{\overline{\tau}_2}=\beta^{u}_{\underline{\tau}_2}=1000$, $\beta^{u}_{\rbu}=100,$ }\DIFaddendFL & 
    \DIFdelbeginFL \DIFdelFL{with $\overline{\tau}_2=-\underline{\tau}_2=3.5$,
    and }\DIFdelendFL \\
    & 
    \DIFdelbeginFL \DIFdelFL{($\beta_{\overline{\tau}_2}=\beta_{\underline{\tau}_2}=100$ and $\beta_{\rbu}=10$) }\DIFdelendFL \DIFaddbeginFL \DIFaddFL{$\overline{\theta}_i=15,\beta_i^\theta=1,\forall i\in\{0,1,2\}$,
    }\DIFaddendFL & 
    \DIFdelbeginFL \DIFdelFL{\eqref{eq:cstr:input:norm} $c_{\rbu}$ with $\overline{\tau}=11$. }\DIFdelendFL \\
    \DIFdelbeginFL %DIFDELCMD < \hline
%DIFDELCMD <         \multirow{3}{*}{(C$_2$)}%%%
\DIFdelendFL & 
    \DIFdelbeginFL \DIFdelFL{CONAC with }\DIFdelendFL \DIFaddbeginFL \DIFaddFL{$\mm{P} = \mm{I}_2$
    }\DIFaddendFL &
    \DIFdelbeginFL \DIFdelFL{\eqref{eq:cstr:input:bound} $c_{\overline\tau_2}$ and $c_{\underline\tau_2}$ }\DIFdelendFL \\
  \DIFaddbeginFL \hline
    \multirow{3}{*}{
      \shortstack{
        {(C$_2$)}
        \\
        \protect\colorLine{my_cyan}{solid}
      }
    }\DIFaddendFL & \DIFdelbeginFL \DIFdelFL{small update rate }\DIFdelendFL \DIFaddbeginFL \DIFaddFL{Proposed CONAC (smaller }\DIFaddendFL $\beta_j$\DIFdelbeginFL \DIFdelFL{. }\DIFdelendFL \DIFaddbeginFL \DIFaddFL{) }\DIFaddendFL &  
    \DIFdelbeginFL \DIFdelFL{with $\overline{\tau}_2=-\underline{\tau}_2=3.5$, and }\DIFdelendFL \\
    & \DIFdelbeginFL \DIFdelFL{($\beta_{\overline{\tau}_2}=\beta_{\underline{\tau}_2}=1$ and $\beta_{\rbu}=0.1$}\DIFdelendFL \DIFaddbeginFL \DIFaddFL{$\beta^{u}_{\overline{\tau}_2}=\beta^{u}_{\underline{\tau}_2}=50$, $\beta^{u}_{\rbu}=10,$ }& \DIFaddFL{same as (C$_1$}\DIFaddendFL ) 
    \DIFaddbeginFL \\ 
    \DIFaddendFL & \DIFdelbeginFL \DIFdelFL{\eqref{eq:cstr:input:norm} $c_{\rbu}$ with $\overline{\tau}=11$. }\DIFdelendFL \DIFaddbeginFL \DIFaddFL{Others: same as (C$_1$) }& 
    \DIFaddendFL \\
\hline
    \DIFdelbeginFL %DIFDELCMD < \multirow{2}{*}{(C$_3$)}%%%
\DIFdelendFL \DIFaddbeginFL \multirow{5}{*}{
      \shortstack{  
      (C$_3$)
      \\
      \protect\colorLine{magenta}{solid}
      }
    }
    \DIFaddendFL & 
      \DIFdelbeginFL \DIFdelFL{CONAC without input constraints.}\DIFdelendFL \DIFaddbeginFL \DIFaddFL{NAC with auxiliary system 
    }\DIFaddendFL &
      \DIFdelbeginFL %DIFDELCMD < \multirow{2}{*}{ \shortstack[l]{No input constraints\\are imposed.} } %%%
\DIFdelendFL \DIFaddbeginFL \multirow{5}{*}{
        \shortstack[l]{  
        $-\overline{\tau}_1 \le \tau_1 \le \overline{\tau}_1$,
        \\
        $-\overline{\tau}_2 \le \tau_2 \le \overline{\tau}_2$, 
        \\
        ($\overline{\tau}_1=10.32, \overline{\tau}_2=3.8$)
        }
      }
    \DIFaddendFL \\
    & 
    \DIFdelbeginFL \DIFdelFL{($\beta_{\overline{\tau}_2}=\beta_{\underline{\tau}_2}=\beta_{\rbu}=0$) }\DIFdelendFL \DIFaddbeginFL \DIFaddFL{\mbox{%DIFAUXCMD
\cite{Rezaei:2025aa,Liu:2025aa,Yu:2022aa} 
    }\hskip0pt%DIFAUXCMD
}\DIFaddendFL & 
    \\
    \DIFaddbeginFL & 
    \DIFaddFL{$\mm{A}_{\zeta} = \mydiag(-120,-100),$ 
    }& 
    \\
    & 
    \DIFaddFL{$\mm{B}_{\zeta} = \mydiag(5\times10^3,5\times10^3),$ 
    }&
    \\
    &
    \DIFaddFL{$\overline{\theta}_i=15,\beta_i^\theta=1,\forall i\in\{0,1,2\}$
    }\\
  \DIFaddendFL \hline
    \DIFdelbeginFL %DIFDELCMD < \multirow{2}{*}{(C$_4$)}%%%
\DIFdelendFL \DIFaddbeginFL \multirow{2}{*}{
    \shortstack{  
      (C$_4$) 
      \\
      \protect\colorLine{gray}{solid}
    }
    }\DIFaddendFL & NAC \DIFdelbeginFL \DIFdelFL{with auxiliary system. }\DIFdelendFL \DIFaddbeginFL \DIFaddFL{without input constraints \mbox{%DIFAUXCMD
\cite{Patil:2022aa} }\hskip0pt%DIFAUXCMD
}\DIFaddendFL & 
    \DIFdelbeginFL \DIFdelFL{$\overline{\tau}_1=-\underline{\tau}_1=10.428$ }\DIFdelendFL \DIFaddbeginFL \multirow{2}{*}{Input is unconstrained}
    \DIFaddendFL \\
    & \DIFdelbeginFL \DIFdelFL{(conventional method) }\DIFdelendFL \DIFaddbeginFL \DIFaddFL{$\overline{\theta}=15$  }\DIFaddendFL & 
    \DIFdelbeginFL \DIFdelFL{$\overline{\tau_2}=-\underline{\tau_2}=3.5$. }\DIFdelendFL \\  
\hline
\DIFaddbeginFL \hline
    \multicolumn{3}{m{8.1cm}}{
      \DIFaddFL{DNN and other parameters for all controllers
    }}
    \\
    \multicolumn{3}{m{8.1cm}}{
      \DIFaddFL{$
        l_0=6,l_1=4,l_2=4,l_3=2,
        \alpha=0.2,
        \estwth(0)\sim\mathcal{U}\bigl([-0.5,0.5]^{\Xi}\bigr)
        ,
      $
    }}
    \\
    \multicolumn{3}{m{8.1cm}}{
      \DIFaddFL{$\NNin = 
      \left(
        (\mv{q}-\mv{q}_d)^\top,
        \ddtt(\mv{q}-\mv{q}_d)^\top,
        \tfrac{\der^2}{\der t^2}\mv{q}_d^\top,
        1
      \right)^\top$
      ,
      $\mm{\Lambda} = \mydiag(6,13)$
    }}
    \\
\hline
  \DIFaddendFL \end{tabular}
\DIFdelbeginFL %DIFDELCMD < \label{table:controllers}
%DIFDELCMD < %%%
\DIFdelendFL \end{table}

% ------------------------------------
%DIF <  Comparative Study
%DIF >  Controllers: CONAC 
% ------------------------------------
\DIFaddbegin \subsubsection*{\DIFadd{Controller 1 }\normalfont{(C$_1$)}\DIFadd{—}\textit{\DIFadd{Proposed CONAC}}}
\DIFaddend 

\DIFdelbegin %DIFDELCMD < \hfill
%DIFDELCMD < 

%DIFDELCMD < %%%
\DIFdel{For comparative analysis, four controllers were implemented, with their properties summarized in Table~\ref{table:controllers}.
}%DIFDELCMD < 

%DIFDELCMD < %%%
%DIF <  ------------------------------------
%DIF <  Controllers: CONAC with large beta 
%DIF <  ------------------------------------
\subsubsection*{\DIFdel{Controller 1 }%DIFDELCMD < \normalfont{(C$_1$)}%%%
\DIFdel{—}\textit{\DIFdel{CONAC with a large update rate $\beta_j,\ \forall j\in\mathcal{I}$, for the Lagrange Multipliers}}%DIFAUXCMD
}
%DIFAUXCMD
%DIFDELCMD < 

%DIFDELCMD < %%%
\DIFdel{To handle }\DIFdelend \DIFaddbegin \DIFadd{To handle all imposed }\DIFaddend input saturation, the proposed CONAC was implemented with the input bound constraints \DIFdelbegin \DIFdel{$c_{\overline\tau_2}$ and $c_{\underline\tau_2}$ \eqref{eq:cstr:input:bound} }\DIFdelend \DIFaddbegin \DIFadd{$c^u_{\overline\tau_2}$ and $c^u_{\underline\tau_2}$ }\DIFaddend for the second link, and the input norm constraint \DIFdelbegin \DIFdel{$c_{\rbu}$ \eqref{eq:cstr:input:norm}, as follows:
}%DIFDELCMD < \begin{equation}
%DIFDELCMD <     c_{\overline\tau_2}     
%DIFDELCMD <     =
%DIFDELCMD <     \tau_2-\overline\tau_2
%DIFDELCMD <     % \le 0
%DIFDELCMD <     ,
%DIFDELCMD <     \;
%DIFDELCMD <     c_{\underline\tau_2} 
%DIFDELCMD <     =
%DIFDELCMD <     \underline\tau_2-\tau_2
%DIFDELCMD <     % \le 0
%DIFDELCMD <     ,
%DIFDELCMD <     \;
%DIFDELCMD <     c_{\rbu}
%DIFDELCMD <     =
%DIFDELCMD <     \tfrac{1}{2}
%DIFDELCMD <     \left(
%DIFDELCMD <         \norm{\rbu}-\overline\tau
%DIFDELCMD <     \right)^2 
%DIFDELCMD <     % <0
%DIFDELCMD <     ,
%DIFDELCMD <     \label{eq:sim:cstr:CONAC}
%DIFDELCMD < \end{equation}
%DIFDELCMD < %%%
\DIFdel{where $\overline\tau_2=-\underline{\tau}_2=3.5$, and $\overline\tau=11$.
Therefore, controller (C$_1$)utilizes the full capability }\DIFdelend \DIFaddbegin \DIFadd{$c^u_{\rbu}$ (defined in \eqref{eq:val:cstr:input}), allowing the controller to fully utilize the authority }\DIFaddend of the actuator\DIFdelbegin \DIFdel{, see \eqref{eq:sat:tau2} and \eqref{eq:sat:norm}.
The weight norm constraints $c_{\theta_i}(\cdot),\ \forall i\in\{0,\cdots,k\}$, defined in \eqref{eq:cstr:weight:ball}, were also incorporated into the CONAC framework with $\overline\theta_0=\overline\theta_1=\overline\theta_2=6$}\DIFdelend .
As described in \DIFdelbegin \DIFdel{\eqref{eq:control:est}, the control }\DIFdelend \DIFaddbegin \DIFadd{\eqref{eq:ctrl:NN:actual}, the }\DIFaddend input (torque) corresponds to the output of the DNN, and the adaptation \DIFdelbegin \DIFdel{law is defined in \eqref{eq:adaptation law}.
}\DIFdelend \DIFaddbegin \DIFadd{laws are defined in \eqref{eq:adap}.
All parameters of (C$_1$) are listed in Table~\ref{table:controller}.
%DIF >  Additionally, the weight norm constraints $c^\theta_{i},\ \forall i\in \mathcal{K}$, defined in \eqref{eq:opt:prob:const1}, were also incorporated into CONAC framework with $\overline\theta_0=\overline\theta_1=\overline\theta_2=6$.
%DIF >  The update rates of the Lagrange multipliers were set to $\beta_{\theta_i}=1,\ \forall i\in\mathcal{K}$, and $\beta_{\overline{\tau}_2} = \beta_{\underline{\tau}_2} = 100$ and $\beta_{\rbu} = 10$, for the weight and input constraints, respectively.
}\DIFaddend 

\DIFdelbegin \DIFdel{The update rates of the Lagrange multipliers for the weight norm constraints were set to $\beta_{\theta_i}=1,\ \forall i\in\{0,\cdots,k\}$.
For the input constraints, large update rates were used: $\beta_{\overline{\tau}_2} = \beta_{\underline{\tau}_2} = 100$ and $\beta_{\rbu} = 10$, to investigate the effect of large update rates.
}%DIFDELCMD < 

%DIFDELCMD < %%%
\DIFdelend % ------------------------------------
% Controllers: CONAC with small beta
% ------------------------------------
\subsubsection*{Controller 2 \normalfont{(C$_2$)}—\DIFdelbegin \textit{\DIFdel{CONAC with a small update rate $\beta_j,\ \forall j\in\mathcal{I}$, for the Lagrange Multipliers}}%DIFAUXCMD
\DIFdelend \DIFaddbegin \textit{\DIFadd{Proposed CONAC with small update rate $\beta^u_j,\ \forall\lambda_j,\ \forall j\in\mathcal{I}$}}\DIFaddend }

\DIFdelbegin \DIFdel{The }\DIFdelend \DIFaddbegin \DIFadd{For the parametric study on $\beta_j^u$ mentioned in Remark~\ref{rem:tuning}, the }\DIFaddend proposed CONAC was implemented with the same input constraints as in (C$_1$), but with \DIFaddbegin \DIFadd{smaller }\DIFaddend update rates for the input constraints\DIFdelbegin \DIFdel{reduced by a factor of $100$: $\beta_{\overline{\tau}_2}=\beta_{\underline{\tau}_2}=1$ and $\beta_{\rbu}=0.1$}\DIFdelend \DIFaddbegin \DIFadd{, as listed in Table~\ref{table:controller}}\DIFaddend .
All other properties and parameters were kept identical to those of (C$_1$).

% ------------------------------------
%DIF <  Controllers: CONAC without ctrl cstr
%DIF <  ------------------------------------
\DIFdelbegin \subsubsection*{\DIFdel{Controller 3 }%DIFDELCMD < \normalfont{(C$_3$)}%%%
\DIFdel{—}\textit{\DIFdel{CONAC without input constraints}}%DIFAUXCMD
}
%DIFAUXCMD
%DIFDELCMD < 

%DIFDELCMD < %%%
\DIFdel{The third controller was identical to (C$_1$) and (C$_2$), except that the input constraints were removed to highlight the constraint-handling capability of the proposed CONAC.
In this case, the update rates of the Lagrange multipliers associated with the input constraints, }%DIFDELCMD < \ie %%%
\DIFdel{$\beta_{\overline{\tau}_2},\beta_{\underline{\tau}_2}$ and $\beta_{\rbu}$, were set to zero, such that the corresponding Lagrange multipliers $[\lambda_j]_{j\in\{\overline{\tau}_2,\underline{\tau}_2,\rbu\}}$ were not updated, see \eqref{eq:adap:lbd}.
All other properties and parameters were kept the same as in (C$_1$) and (C$_2$).
%DIF <  \MSRY{
%DIF <      weight constraint관련된 파라미터는 C1-3까지 모두 같습니다.
%DIF <      C2의 control input 관련 beta만 100배 작아졌다는 말 보충하였습니다.
%DIF <  }
}%DIFDELCMD < 

%DIFDELCMD < %%%
%DIF <  ------------------------------------
\DIFdelend % Controllers: Aux. system in CONAC-AUX
% ------------------------------------
\subsubsection*{Controller \DIFdelbegin \DIFdel{4 }%DIFDELCMD < \normalfont{(C$_4$)}%%%
\DIFdelend \DIFaddbegin \DIFadd{3 }\normalfont{(C$_3$)}\DIFaddend —\DIFdelbegin \textit{\DIFdel{NAC with auxiliary system (conventional method)}}%DIFAUXCMD
\DIFdelend \DIFaddbegin \textit{\DIFadd{NAC with auxiliary system}}\DIFaddend }

As a baseline for comparison \DIFaddbegin \DIFadd{in handling input constraints}\DIFaddend , the auxiliary \DIFdelbegin \DIFdel{system-based approach from \mbox{%DIFAUXCMD
\cite{Esfandiari:2014aa, Karason:1994aa, Esfandiari:2015aa} }\hskip0pt%DIFAUXCMD
was implemented.
%DIF <  in place of the proposed CONAC for handling input constraints.
}\DIFdelend \DIFaddbegin \DIFadd{system used in \mbox{%DIFAUXCMD
\cite{Rezaei:2025aa,Liu:2025aa,Yu:2022aa} }\hskip0pt%DIFAUXCMD
was adopted.
%DIF >  Gao:2017aa,Esfandiari:2014aa,Esfandiari:2015aa,Liu:2024aa,Liu:2025aa,Yu:2022aa,Rezaei:2025aa
It should be noted that the auxiliary system is limited to handling only the element-wise input bound constraints.
}\DIFaddend The auxiliary system is defined by 
\DIFdelbegin \DIFdel{$\ddtt{\mv\zeta} = \mm{A}_\zeta \mv{\zeta} + \mm{B}_\zeta \Delta\rbu$, with initial condition $\mv{\zeta}\vert_{t=0} = \mv{0}_{2\times1}$, }\DIFdelend \DIFaddbegin \begin{equation}
  \ddtt \mv{\zeta} 
  =
  \mm{A}_{\zeta}
  \mv{\zeta}
  +
  \mm{B}_{\zeta}
  \Delta_{\zeta} \mv{\tau}
  ,
  \quad
  \zeta(0) 
  = 
  \mv{0}_{2\times 1}
  ,
\end{equation}
\DIFaddend where $\mv{\zeta}\in\R^2$ \DIFdelbegin \DIFdel{represents }\DIFdelend \DIFaddbegin \DIFadd{denotes }\DIFaddend the auxiliary state \DIFdelbegin \DIFdel{variables, and $\mm{A}_\zeta\in\R_{<0}^{2\times 2}$ and $\mm{B}_\zeta\in\R^{2\times 2}$ are user-designed matrices.
The auxiliary state $\mv{\zeta}$ is driven by the overflow amount of the control input $\Delta\rbu$, where each element is defined as $\Delta\tau_{i} = \tau_{i}-\max(\underline\tau_i,\min(\overline\tau_i,\tau_i)),\ \forall i\in\{1,2\}$.
}%DIFDELCMD < 

%DIFDELCMD < %%%
\DIFdel{This }\DIFdelend \DIFaddbegin \DIFadd{which is initialized to zero, $\mm{A}_{\zeta}\in\R^{2\times 2}$ is a stable matrix, $\mm{B}_{\zeta}\in\R^{2\times 2}$ is a gain matrix, and $\Delta_{\zeta} \tau_i := \tau_i - \max\{\min\{\tau_i, \overline{\tau}_i\}, -\overline{\tau}_i\},\ \forall i\in\{1,2\}$.
%DIF >  \operatorname{clip}(u_i, \underline{u}_i, \overline{u}_i),\ \forall i\in\{1,2\}$.
It is notable that $\Delta_{\zeta}\mv{u}=\Delta_{\zeta}\mv{\tau}$ is defined as the element-wise input overflow, which is different from the actual input saturation $\Delta\mv{u}$ in \eqref{eq:Vc:dot:1}.
%DIF >  The limits of the control input are set as $\underline{u}_1=-\overline{u}_1=10.428$ and $\underline{u}_2=-\overline{u}_2=3.5$, to approximate the input norm constraint $\overline{\tau} = 11$.
The }\DIFaddend auxiliary state was \DIFdelbegin \DIFdel{incorporated into }\DIFdelend \DIFaddbegin \DIFadd{used in }\DIFaddend the adaptation law \DIFdelbegin \DIFdel{by replacing the filtered tracking error $\fe$ with $\fe+\mv{\zeta}$.
As a result, the weights were adapted to minimize both the filtered tracking error and the auxiliary state, thereby mitigating saturation in each control input.
}%DIFDELCMD < 

%DIFDELCMD < %%%
\DIFdel{The }\DIFdelend \DIFaddbegin \DIFadd{\eqref{eq:adap:th:approx}, by substituting the tracking error $\mv{e}$ with the }\DIFaddend auxiliary \DIFdelbegin \DIFdel{system matrices were chosen as $\mm{A}_\zeta=\mydiag(-10,-10)$ and $\mm{B}_\zeta=\mydiag(10^3,10^3)$, 
to balance }\DIFdelend \DIFaddbegin \DIFadd{error $\mv{e}_\zeta := \mv{e} + \mv{\zeta}$.
The parameters of the auxiliary system were chosen to balance the }\DIFaddend responsiveness to input overflow with the convergence speed of the auxiliary dynamics\DIFdelbegin \DIFdel{.
Since the auxiliary system only accounts for element-wise input saturation—corresponding to the input bound constraint \eqref{eq:cstr:input:bound}—the upper bound of the first actuator torque $\overline{\tau}_1=-\underline{\tau}_1$ was set to $10.428$ to approximate the input norm constraint $\overline{\tau} = 11$, given $\overline{\tau}_1^2 + \overline{\tau}_2^2 = \overline{\tau}^2$.
}%DIFDELCMD < 

%DIFDELCMD < %%%
\DIFdelend \DIFaddbegin \DIFadd{, as listed in Table~\ref{table:controller}.
}\DIFaddend All other properties and parameters were kept the same as in (C$_1$) \DIFdelbegin \DIFdel{, }\DIFdelend \DIFaddbegin \DIFadd{and }\DIFaddend (C$_2$)\DIFdelbegin \DIFdel{, and (C}\DIFdelend \DIFaddbegin \DIFadd{.
}

%DIF >  ------------------------------------
%DIF >  Controllers: NAC without ctrl cstr
%DIF >  ------------------------------------
\subsubsection*{\DIFadd{Controller 4 }\normalfont{(C$_4$)}\DIFadd{—}\textit{\DIFadd{NAC without input constraints}}}

\DIFadd{The fourth controller, adopted from \mbox{%DIFAUXCMD
\cite{Patil:2022aa}}\hskip0pt%DIFAUXCMD
, does not handle input constraints.
The controller structure is identical to (C$_1$)--(C}\DIFaddend $_3$), except that the \DIFaddbegin \DIFadd{input constraints were removed and }\DIFaddend weight boundedness was handled using the projection operator\DIFdelbegin \DIFdel{used in \mbox{%DIFAUXCMD
\cite{Patil:2022aa}}\hskip0pt%DIFAUXCMD
}\DIFdelend \DIFaddbegin \DIFadd{, as follows:
}\begin{equation}
    \ddtt\estwth
    =
    \myproj_{\norm{\estwth}\le\overline{\theta}}
    \left(
        -\alpha
        \pptfrac{\estNN}{\estwth}^\top
        \mm{P}
        \mv{e}
    \right)
    .
\end{equation}
\DIFadd{The weight norm bound was set to $\overline{\theta}=15$.
All other properties and parameters were kept the same as in (C$_1$)--(C$_3$), as summarized in Table~\ref{table:controller}}\DIFaddend .

\hfill

% ------------------------------------
% Controller Parameters
% ------------------------------------
For all controllers, \DIFdelbegin %DIFDELCMD < \ie %%%
\DIFdel{(C$_1$), (C$_2$), (C$_3$), and (C$_4$), the DNNs shared the same architecture, consisting of two hidden layers with four nodes each, }%DIFDELCMD < \ie %%%
\DIFdel{$k=2$, and $l_0=6, l_1=4, l_2=4$, and $l_3=2$ in \eqref{eq:DNN:architecture}.
The input vector to the DNNs was defined as ${\q}_n=\left(\q^\top,\dq^\top,\fe^\top,1\right)^\top\in\R^{7}$, where the scalar $1$ was included to account for the bias term in the weight matrix.
The adaptation gain was set to $\alpha = 0.2$, and the filtering matrix in \eqref{eq:err:filtered} was chosen as $\mm{\Lambda}=\mydiag(5,15)$.
To support real-time implementation, the controller operated at a sampling rate of $250$ Hz, with control results collected via CAN communication using a PCAN-USB interface and recorded on a computer at the same rate.
}%DIFDELCMD < 

%DIFDELCMD < %%%
\DIFdel{The tracking performance was quantitatively evaluated using the $L_2$-norm of the each joint angle's tracking error, defined as $\sqrt{\int_0^T\norm{e_i(t)}\,\der t},\ \forall i\in\{1,2\}$, where $T\in\R_{>0}$ denotes the duration of each episode}\DIFdelend \DIFaddbegin \DIFadd{the remaining parameters, including the DNN architecture, were kept identical to ensure a fair comparison, as listed in Table~\ref{table:controller}.
Specifically, invoking Remark~\ref{rem:init}, the DNN weights were initialized with small random values according to $\estwth(0)\sim\mathcal{U}\bigl([-0.5,0.5]^{\Xi}\bigr)$, where $\mathcal{U}(\cdot)$ denotes a uniform random distribution}\DIFaddend .

\DIFaddbegin \subsection{\DIFadd{Validation Results}} \label{sec:val:results}

\DIFaddend \begin{figure*}[t]
  \centering
  \DIFdelbeginFL %DIFDELCMD < \subfloat[Tracking performance of $1$st joint angle $q_1$.]{
%DIFDELCMD <     %   \includegraphics
%DIFDELCMD <         \includegraphics[width=\figSizeTwoCol\linewidth]
%DIFDELCMD <         {
%DIFDELCMD <         % fig/Fig1.pdf
%DIFDELCMD <         src/measurement/figures/compare/Fig1.eps
%DIFDELCMD <         }%
%DIFDELCMD <         \label{fig:ctrl:result:q1}}
%DIFDELCMD <     %%%
\DIFdelendFL \DIFaddbeginFL \subfloat[Tracking performance of $1$st joint angle $q_1$.]{
    \includegraphics[width=.49\linewidth]
    {
      \resultFig/Fig1.eps
    }%
    \label{fig:ctrl:result:general:q1}}
  \DIFaddendFL \hfill
  %DIF <  \vfill
    \DIFdelbeginFL %DIFDELCMD < \subfloat[Tracking performance of $2$nd joint angle $q_2$.]{
%DIFDELCMD <         \includegraphics[width=\figSizeTwoCol\linewidth]
%DIFDELCMD <         {
%DIFDELCMD <         % src/measurement/figures/compare/Fig11.eps
%DIFDELCMD <         src/measurement/figures/compare/Fig2.eps
%DIFDELCMD <         }%
%DIFDELCMD <         \label{fig:ctrl:result:q2}}
%DIFDELCMD <     %%%
\DIFdelendFL \DIFaddbeginFL \subfloat[Tracking performance of $2$nd joint angle $q_2$.]{
    \includegraphics[width=.49\linewidth]
    {
      \resultFig/Fig2.eps
    }%
    \label{fig:ctrl:result:general:q2}}
  \DIFaddendFL \vfill
  \DIFdelbeginFL %DIFDELCMD < \subfloat[Control input of the $1$st joint $\tau_1$.]{
%DIFDELCMD <     \includegraphics[width=\figSizeTwoCol\linewidth]
%DIFDELCMD <         {
%DIFDELCMD <             src/measurement/figures/compare/Fig3.eps
%DIFDELCMD <         }%
%DIFDELCMD <         \label{fig:ctrl:result:tau:1}}
%DIFDELCMD <     \hfill
%DIFDELCMD <     %%%
\DIFdelendFL % \vfill
  \DIFdelbeginFL %DIFDELCMD < \subfloat[Control input of the 2nd joint $\tau_2$.]{
%DIFDELCMD <         \includegraphics[width=\figSizeTwoCol\linewidth]
%DIFDELCMD <         {
%DIFDELCMD <             src/measurement/figures/compare/Fig4.eps
%DIFDELCMD <         }%
%DIFDELCMD <         \label{fig:ctrl:result:tau:2}}
%DIFDELCMD <     %%%
\DIFdelendFL \DIFaddbeginFL \subfloat[Commanded control input of the $1$st joint $\tau_1$.]{
  \includegraphics[width=.49\linewidth]
    {
      \resultFig/Fig5.eps
    }%
    \label{fig:ctrl:result:general:tau:1}}
  \hfill
    \subfloat[Commanded control input of the 2nd joint $\tau_2$.]{
    \includegraphics[width=.49\linewidth]
    {
      \resultFig/Fig6.eps
    }%
    \label{fig:ctrl:result:general:tau:2}}
  \DIFaddendFL \vfill
  \DIFdelbeginFL %DIFDELCMD < \subfloat[Norm of the joint torque input $\rbu$.]{
%DIFDELCMD <         \includegraphics[width=\figSizeTwoCol\linewidth]
%DIFDELCMD <         {
%DIFDELCMD <             src/measurement/figures/compare/Fig5.eps
%DIFDELCMD <         }%
%DIFDELCMD <         \label{fig:ctrl:result:tau:norm}}
%DIFDELCMD <     %%%
\DIFdelendFL \DIFaddbeginFL \subfloat[Norm of the commanded control input $\rbu$.]{
    \includegraphics[width=.49\linewidth]
    {
      \resultFig/Fig7.eps
    }%
    \label{fig:ctrl:result:general:tau:norm}}
  \DIFaddendFL \hfill
  % \vfill
  \DIFdelbeginFL %DIFDELCMD < \subfloat[Norm of DNN weights ${\estwth}_i$.]{
%DIFDELCMD <         % \includegraphics
%DIFDELCMD <         \includegraphics[width=\figSizeTwoCol\linewidth]
%DIFDELCMD <         {
%DIFDELCMD <             % fig/Fig7.pdf
%DIFDELCMD <             src/measurement/figures/compare/Fig7.eps
%DIFDELCMD <         }%
%DIFDELCMD <         \label{fig:ctrl:result:weight}}        
%DIFDELCMD <     %%%
\DIFdelendFL \DIFaddbeginFL \subfloat[Norm of DNN weights ${\estwth}_i$.]{
    % \includegraphics
    \includegraphics[width=.49\linewidth]
    {
      \resultFig/Fig10.eps
    }%
    \label{fig:ctrl:result:general:weight}}        
  \DIFaddendFL \caption{
      Real-time implementation results of (C$_1$) \protect\colorLine{blue}{solid}, (C$_2$) \DIFdelbeginFL %DIFDELCMD < \protect\colorLine{cyan}{solid}%%%
\DIFdelendFL \DIFaddbeginFL \protect\colorLine{my_cyan}{solid}\DIFaddendFL , (C$_3$) \DIFdelbeginFL %DIFDELCMD < \protect\colorLine{gray}{solid}%%%
\DIFdelendFL \DIFaddbeginFL \protect\colorLine{magenta}{solid}\DIFaddendFL , (C$_4$) \DIFdelbeginFL %DIFDELCMD < \protect\colorLine{magenta}{solid}%%%
\DIFdelendFL \DIFaddbeginFL \protect\colorLine{gray}{solid}\DIFaddendFL , and desired trajectory $\qd$ \protect\colorLine{red}{dashed}.
      % \vspace{-5mm}
  }
\DIFdelbeginFL %DIFDELCMD < \label{fig:ctrl:result}
%DIFDELCMD < %%%
\DIFdelendFL \DIFaddbeginFL \label{fig:ctrl:result:general}
\DIFaddendFL \end{figure*}

%DIF <  ------------------------------------
%DIF <  Norm of error in Episode 1: 
%DIF <  C1 e1 ep1: 25.130
%DIF <  C1 e2 ep1: 25.016
%DIF <  C2 e1 ep1: 21.719
%DIF <  C2 e2 ep1: 10.530
%DIF <  C3 e1 ep1: 23.168
%DIF <  C3 e2 ep1: 13.061
%DIF <  C4 e1 ep1: 24.195
%DIF <  C4 e2 ep1: 12.927
%DIF <  Norm of error in Episode 2: 
%DIF <  C1 e1 ep2: 7.066
%DIF <  C1 e2 ep2: 2.198
%DIF <  C2 e1 ep2: 8.021
%DIF <  C2 e2 ep2: 2.435
%DIF <  C3 e1 ep2: 8.408
%DIF <  C3 e2 ep2: 3.554
%DIF <  C4 e1 ep2: 19.721
%DIF <  C4 e2 ep2: 4.256
%DIF <  Improvement in Episode 2: 
%DIF <  C1 e1: 0.719
%DIF <  C1 e2: 0.912
%DIF <  C2 e1: 0.631
%DIF <  C2 e2: 0.769
%DIF <  C3 e1: 0.637
%DIF <  C3 e2: 0.728
%DIF <  C4 e1: 0.185
%DIF <  C4 e2: 0.671
%DIF <  ------------------------------------
\DIFdelbegin %DIFDELCMD < 

%DIFDELCMD < \begin{table}[!t]
%DIFDELCMD <     \renewcommand{\arraystretch}{1.3}
%DIFDELCMD <     %%%
%DIFDELCMD < \caption{%
{%DIFAUXCMD
\DIFdelFL{Quantitative\,comparison\,of\,performances'\,$L_2$\,norm.}}
    %DIFAUXCMD
\DIFdelendFL \DIFaddbeginFL \begin{figure}[t]
  \DIFaddendFL \centering
  \DIFdelbeginFL %DIFDELCMD < \begin{tabular}{c c c c c}
%DIFDELCMD <     \hline
%DIFDELCMD < 		& %%%
%DIFDELCMD < \multicolumn{4}{c}{%
{%DIFAUXCMD
\textbf{\DIFdelFL{Episode 1}}%DIFAUXCMD
} %DIFAUXCMD
%DIFDELCMD < \\
%DIFDELCMD <     \hline
%DIFDELCMD < 	\hline 
%DIFDELCMD < 		& %%%
\DIFdelFL{(C$_1$) }%DIFDELCMD < & %%%
\DIFdelFL{(C$_2$) }%DIFDELCMD < & %%%
\DIFdelFL{(C$_3$) }%DIFDELCMD < & %%%
\DIFdelFL{(C$_4$) }%DIFDELCMD < \\
%DIFDELCMD < 	\hline
%DIFDELCMD < 		%%%
\DIFdelFL{$e_1\times10^{3}$ / rad }%DIFDELCMD < & %%%
\DIFdelFL{$25.130$ }%DIFDELCMD < & %%%
\DIFdelFL{$21.719$ }%DIFDELCMD < & %%%
\DIFdelFL{$23.168$ }%DIFDELCMD < & %%%
\DIFdelFL{$24.195$ }%DIFDELCMD < \\ 
%DIFDELCMD < 	\hline
%DIFDELCMD <         %%%
\DIFdelFL{$e_2\times10^{3}$ / rad }%DIFDELCMD < & %%%
\DIFdelFL{$25.016$ }%DIFDELCMD < & %%%
\DIFdelFL{$10.530$ }%DIFDELCMD < & %%%
\DIFdelFL{$13.061$ }%DIFDELCMD < & %%%
\DIFdelFL{$12.927$ }%DIFDELCMD < \\
%DIFDELCMD < 	\hline
%DIFDELCMD <         & %%%
%DIFDELCMD < \multicolumn{4}{c}{%
{%DIFAUXCMD
\textbf{\DIFdelFL{Episode 2}}%DIFAUXCMD
} %DIFAUXCMD
%DIFDELCMD < \\
%DIFDELCMD <     \hline
%DIFDELCMD <     \hline
%DIFDELCMD <         & %%%
\DIFdelFL{(C$_1$) }%DIFDELCMD < & %%%
\DIFdelFL{(C$_2$) }%DIFDELCMD < & %%%
\DIFdelFL{(C$_3$) }%DIFDELCMD < & %%%
\DIFdelFL{(C$_4$) }%DIFDELCMD < \\
%DIFDELCMD < 	\hline
%DIFDELCMD <     \multirow{2}{*}{$e_1\times10^{3}$ / rad} 
%DIFDELCMD <         & %%%
\DIFdelFL{$7.066$ }%DIFDELCMD < & %%%
\DIFdelFL{$8.021$ }%DIFDELCMD < & %%%
\DIFdelFL{$8.408$ }%DIFDELCMD < & %%%
\DIFdelFL{$19.721$ }%DIFDELCMD < \\
%DIFDELCMD <         & %%%
\DIFdelFL{($-71.6\%$) }%DIFDELCMD < & %%%
\DIFdelFL{($-63.1\%$) }%DIFDELCMD < & %%%
\DIFdelFL{($-63.7\%$) }%DIFDELCMD < & %%%
\DIFdelFL{($-18.5\%$) }%DIFDELCMD < \\
%DIFDELCMD <     \hline
%DIFDELCMD <     \multirow{2}{*}{$e_2\times10^{3}$ / rad} 
%DIFDELCMD <         & %%%
\DIFdelFL{$2.198$ }%DIFDELCMD < & %%%
\DIFdelFL{$2.435$ }%DIFDELCMD < & %%%
\DIFdelFL{$3.554$ }%DIFDELCMD < & %%%
\DIFdelFL{$4.256$ }%DIFDELCMD < \\
%DIFDELCMD <         & %%%
\DIFdelFL{($-88.4\%$) }%DIFDELCMD < & %%%
\DIFdelFL{($-76.9\%$) }%DIFDELCMD < & %%%
\DIFdelFL{($-72.8\%$) }%DIFDELCMD < & %%%
\DIFdelFL{($-67.1\%$) }%DIFDELCMD < \\
%DIFDELCMD <     \hline
%DIFDELCMD <     \end{tabular}
%DIFDELCMD <     \label{tab:sim:L2}
%DIFDELCMD < \end{table}
%DIFDELCMD < %%%
\DIFdelend \DIFaddbegin \includegraphics[width=0.99\linewidth]{\resultFig/Fig21.eps}
  \caption{
    \DIFadd{Quantitative comparison of tracking performance.
  }}
  \label{fig:ctrl:results:tracking}
\end{figure}
\DIFaddend 

\DIFdelbegin \subsection{\DIFdel{Validation Results}}
%DIFAUXCMD
\addtocounter{subsection}{-1}%DIFAUXCMD
\DIFdelend %DIF >  \begin{table}[!t]
%DIF >    \renewcommand{\arraystretch}{1.3}
%DIF >    \caption{Quantitative comparison of tracking performance.}
%DIF >    \centering
%DIF >    \begin{tabular}{c c c c c}
%DIF >    \hline
%DIF >    & \multicolumn{4}{c}{\textbf{Episode 1}} \\
%DIF >    \hline
%DIF >    \hline 
%DIF >    \qty{}{\deg} & {(C$_1$)} \protect\colorLine{blue}{solid} & {(C$_2$)} \protect\colorLine{my_cyan}{solid} & {(C$_3$)} \protect\colorLine{magenta}{solid} & {(C$_4$)} \protect\colorLine{gray}{solid} \\
%DIF >    \hline
%DIF >      $q_1-{q_d}_1$
%DIF >      & \CONAChighEpiOneAngOne & \CONAClowEpiOneAngOne & \AUXEpiOneAngOne & \NACEpiOneAngOne \\ 
%DIF >    \hline
%DIF >      $q_2-{q_d}_2$
%DIF >      & \CONAChighEpiOneAngTwo & \CONAClowEpiOneAngTwo & \AUXEpiOneAngTwo & \NACEpiOneAngTwo \\
%DIF >    \hline
%DIF >      & \multicolumn{4}{c}{\textbf{Episode 2}} \\
%DIF >    \hline
%DIF >    \hline
%DIF >    \qty{}{\deg} & {(C$_1$)} \protect\colorLine{blue}{solid} & {(C$_2$)} \protect\colorLine{my_cyan}{solid} & {(C$_3$)} \protect\colorLine{magenta}{solid} & {(C$_4$)} \protect\colorLine{gray}{solid} \\
%DIF >  \hline
%DIF >    \multirow{2}{*}{$q_1-{q_d}_1$} 
%DIF >      & \CONAChighEpiTwoAngOne & \CONAClowEpiTwoAngOne & \qty{\AUXEpiTwoAngOne}{\radian} & \NACEpiTwoAngOne \\
%DIF >      & (\CONAChighImpAngOne\%) & (\CONAClowImpAngOne\%) & (\AUXImpAngOne\%) & (\NACImpAngOne\%) \\
%DIF >    \hline
%DIF >    \multirow{2}{*}{$q_2-{q_d}_2$} 
%DIF >      & \CONAChighEpiTwoAngTwo & \CONAClowEpiTwoAngTwo & \AUXEpiTwoAngTwo & \NACEpiTwoAngTwo \\
%DIF >      & (\CONAChighImpAngTwo\%) & (\CONAClowImpAngTwo\%) & (\AUXImpAngTwo\%) & (\NACImpAngTwo\%) \\
%DIF >    \hline
%DIF >    \end{tabular}
%DIF >    \label{table:results:tracking}
%DIF >  \end{table}

%DIF <  ------------------------------------
%DIF <  Tracking Performance
%DIF <  ------------------------------------
\DIFdelbegin \subsubsection{\DIFdel{Tracking Performance}}
%DIFAUXCMD
\addtocounter{subsubsection}{-1}%DIFAUXCMD
\DIFdelend %DIF >  --------------------------------------------
%DIF >  val. result: tracking performance
%DIF >  --------------------------------------------

\DIFaddbegin \subsubsection{\DIFadd{Tracking and Adaptation Performance}}

%DIF >  초기화 + 웜업 페이즈 + 일반적인 학습의 성능

\DIFaddend The real-time implementation results are \DIFdelbegin \DIFdel{presented }\DIFdelend \DIFaddbegin \DIFadd{illustrated }\DIFaddend in Fig.~\DIFdelbegin \DIFdel{\ref{fig:ctrl:result}, and the quantitative comparison of tracking performance is summarized in Table~\ref{tab:sim:L2}.Across the two episodes}\DIFdelend \DIFaddbegin \DIFadd{\ref{fig:ctrl:result:general}.
As shown in Fig.~\ref{fig:ctrl:result:general:q1} and Fig.~\ref{fig:ctrl:result:general:q2}}\DIFaddend , all controllers\DIFdelbegin \DIFdel{improved their tracking performance. 
Specifically, the $L_2$-norm of $e_1$ was reduced by $71.6 \%$ for (C$_1$), $63.1 \%$ for (C$_2$), $63.7 \%$ for (C$_3$), and $18.5 \%$ for (C$_4$); while the $L_2$-norm of $e_2$ was reduced by $88.4 \%$}\DIFdelend , \DIFdelbegin \DIFdel{$76.9 \%$, $72.8 \%$, and $67.1 \%$, respectively.
Qualitatively, as illustrated }\DIFdelend \DIFaddbegin \DIFadd{except (C$_4$), tracked the desired trajectory with their tracking performance improving from
Episode~1 to Episode~2.
It should be noted that, due to the input saturation effect, the tracking performance of all controllers is limited.
A quantitative comparison is summarized }\DIFaddend in Fig.~\DIFdelbegin \DIFdel{\ref{fig:ctrl:result:q1} and Fig.~\ref{fig:ctrl:result:q2}, oscillations were suppressed during the second episode.
These results indicate that the DNNs in all controllers successfully learned the ideal control input $\rbu^*$ using the filtered tracking error $\fe$, despite $\rbu^*$ being completely unknown a priori.
Furthermore, the stable learning observed under random weight initialization and without prior knowledge of the system demonstrated the effectiveness of the proposed CONAC in enabling online learning capability.}%DIFDELCMD < 

%DIFDELCMD < %%%
\DIFdel{Notably, in both episodes, all controllersfailed to track the desired trajectory of $q_1$ during certain time intervals: from $17.5$ s to $19.9$ s in }\DIFdelend \DIFaddbegin \DIFadd{\ref{fig:ctrl:results:tracking}, using the root mean square error (RMSE) of the tracking error for each joint angle during 2 episodes.
As shown in Fig.~\ref{fig:ctrl:results:tracking}, all controllers, except (C$_4$), improved their tracking performance from }\DIFaddend Episode $1$ \DIFdelbegin \DIFdel{, and from $23.5$ s to $25.9$ s and from $29.5$ s to $31.9$ s in }\DIFdelend \DIFaddbegin \DIFadd{to }\DIFaddend Episode $2$ \DIFdelbegin \DIFdel{, as shown in Fig.~\ref{fig:ctrl:result:q1}.
This degradation in tracking performance was attributed to actuator saturation during those intervals, as observed in Fig.
~\ref{fig:ctrl:result:tau:2} and Fig.~\ref{fig:ctrl:result:tau:norm}.
}\DIFdelend \DIFaddbegin \DIFadd{by around 30\% and 70\% for the first and second joint angles, respectively, demonstrating the adaptation capability of the DNN.
}\DIFaddend 

\DIFdelbegin \DIFdel{The tracking performance of CONAC with input constraints—}%DIFDELCMD < \ie %%%
\DIFdelend \DIFaddbegin \DIFadd{Remarkably, the proposed CONAC }\DIFaddend (C$_1$) and (C$_2$) \DIFdelbegin \DIFdel{—outperformed the auxiliary system-based NAC (C$_4$) in the second episode, as shown in Table~\ref{tab:sim:L2}.
This improvement was attributed to the ability of (C$_1$) and (C$_2$) to }\DIFdelend \DIFaddbegin \DIFadd{demonstrated superior tracking performance compared to the other controllers.
This is because (C$_3$) could not }\DIFaddend fully utilize the actuator \DIFdelbegin \DIFdel{'s capacity, whereas (C}\DIFdelend \DIFaddbegin \DIFadd{authority due to the auxiliary system design, which will be discussed in detail in Section~\ref{sec:val:results:constraint}.
For (C}\DIFaddend $_4$\DIFdelbegin \DIFdel{) was limited by the design of the auxiliary system , which constrained the first joint actuator.
While (C$_3$) also showed better performance than (C$_4$}\DIFdelend ), \DIFdelbegin \DIFdel{this was primarily because it did not consider input saturation and thus used more of the actuator's capacity.
However, since (C$_3$) did not handle the control input saturation, it generated excessively large control efforts to minimize }\DIFdelend \DIFaddbegin \DIFadd{since it does not consider the input saturation in the adaptation process, }\DIFaddend the \DIFaddbegin \DIFadd{DNN weights were continuously increased in an attempt to reduce the }\DIFaddend tracking error, even though \DIFdelbegin \DIFdel{these efforts }\DIFdelend \DIFaddbegin \DIFadd{the corresponding control input }\DIFaddend could not be realized \DIFdelbegin \DIFdel{due to actuator limits, }\DIFdelend \DIFaddbegin \DIFadd{by the saturated actuators.
Consequently, the DNN weights exhibited persistent growth rather than convergence }\DIFaddend as shown in Fig.~\DIFdelbegin \DIFdel{\ref{fig:ctrl:result:tau:1}–\ref{fig:ctrl:result:tau:norm}.
In addition, the severity of saturation violations in (C$_3$) increased during periods when the input constraints were active.
These large violations pose a significant risk to the physical system, potentially causing actuator damage or failure. 
Furthermore, once the control saturation was lifted—typically due to changes in the desired trajectory—this resulted in oscillatory behavior, as CONAC's weights had been adapted to generate excessively high control inputs, see Fig.~\ref{fig:ctrl:result:tau:norm}, where the input constraints were lifted at $25.9$ s and $31.9$ s.
}\DIFdelend \DIFaddbegin \DIFadd{\ref{fig:ctrl:result:general:weight}, which may eventually lead to parameter drift or unbounded weight growth and is therefore undesirable.
}\DIFaddend 

%DIF <  ------------------------------------
%DIF <  Control Input Saturation
%DIF <  ------------------------------------
%DIF >  --------------------------------------------
%DIF >  val. result: constraint performance
%DIF >  --------------------------------------------
\DIFaddbegin 

\subsubsection{\DIFadd{Constraint Handling Performance}} \label{sec:val:results:constraint}

\DIFaddend \begin{figure}[t]
  \centering
  \DIFdelbeginFL %DIFDELCMD < \subfloat[Control input $\rbu$ locus during the time interval from $29.4$ s to $31.9$ s.]{
%DIFDELCMD <     % \subfloat[Auxiliary states $\mv{\zeta}$ of CONAC-AUX.]{
%DIFDELCMD <         \includegraphics[width=\figSizeOneCol\linewidth]
%DIFDELCMD <         {
%DIFDELCMD <             src/measurement/figures/compare/Fig6.eps
%DIFDELCMD <         }%
%DIFDELCMD <         \label{fig:ctrl:result:scope:control}}
%DIFDELCMD <         %%%
\DIFdelendFL \DIFaddbeginFL \subfloat[Trajectory of the tracking error of each joint angle $\mv{q}$ during the time interval {$[\zoomStartTime, \qty{\zoomEndTime}{\second}]$}.]{
    \includegraphics[width=0.97\linewidth]
    {
      \resultFig/Fig16.eps
    }
    \label{fig:ctrl:result:detail:q:top:view}
  }
  \DIFaddendFL \vfill
  \DIFdelbeginFL %DIFDELCMD < \subfloat[Lagrange multipliers $\lambda_j$ during the time interval from $29.4$ s to $31.9$ s.]{
%DIFDELCMD <         \includegraphics[width=\figSizeOneCol\linewidth]
%DIFDELCMD <         {
%DIFDELCMD <             src/measurement/figures/compare/Fig8.eps
%DIFDELCMD <         }%
%DIFDELCMD <         \label{fig:ctrl:result:scope:beta}}
%DIFDELCMD <         %%%
\DIFdelendFL \DIFaddbeginFL \subfloat[Trajectory of the commanded control input $\mv{\tau}$ during the time interval {$[\zoomStartTime, \qty{\zoomEndTime}{\second}]$}.]{
    \includegraphics[width=0.97\linewidth]
    {
      \resultFig/Fig8.eps
    }
    \label{fig:ctrl:result:detail:u:top:view}
  }
  \DIFaddendFL \vfill
  \DIFdelbeginFL %DIFDELCMD < \subfloat[Auxiliary states $\mv{\zeta}$ during the time interval from $29.4$ s to $31.9$ s.]{
%DIFDELCMD <         \includegraphics[width=\figSizeOneCol\linewidth]
%DIFDELCMD <         {
%DIFDELCMD <             src/measurement/figures/compare/Fig9.eps
%DIFDELCMD <         }%
%DIFDELCMD <         \label{fig:ctrl:result:scope:aux}}
%DIFDELCMD <     %%%
%DIFDELCMD < \caption{%
{%DIFAUXCMD
\DIFdelFL{Constraint handling behavior of (C$_1$) }%DIFDELCMD < \protect\colorLine{blue}{solid}%%%
\DIFdelFL{, (C$_2$) }%DIFDELCMD < \protect\colorLine{cyan}{solid}%%%
\DIFdelFL{, (C$_3$) }%DIFDELCMD < \protect\colorLine{gray}{solid}%%%
\DIFdelFL{, and (C$_4$) }%DIFDELCMD < \protect\colorLine{magenta}{solid} %%%
\DIFdelFL{during the time interval from $29.4$ s to $31.9$ s.
        %DIF <  \vspace{-5mm}
    }}
%DIFAUXCMD
%DIFDELCMD < \label{fig:ctrl:result:scope}
%DIFDELCMD < %%%
\DIFdelendFL \DIFaddbeginFL \subfloat[Lagrange multiplier associated with the input norm constraint for (C$_1$) and (C$_2$) during the time interval {$[\zoomStartTime, \qty{\zoomEndTime}{\second}]$}.]{
    \includegraphics[width=0.97\linewidth]
    {
      \resultFig/Fig9.eps
    }
    \label{fig:ctrl:result:detail:multipliers}
  }
  \caption{  
    \DIFaddFL{Zoomed-in view of the tracking error, commanded control input, and
Lagrange multiplier $\lambda_{\overline{\tau}}$ during the time interval
$[\zoomStartTime, \qty{\zoomEndTime}{\second}]$ of (C$_1$) }\protect\colorLine{blue}{solid}\DIFaddFL{, (C$_2$) }\protect\colorLine{my_cyan}{solid}\DIFaddFL{, (C$_3$) }\protect\colorLine{magenta}{solid}\DIFaddFL{, and (C$_4$) }\protect\colorLine{gray}{solid}\DIFaddFL{.
  }}
  \label{fig:ctrl:result:detail}
\DIFaddendFL \end{figure}

\DIFdelbegin %DIFDELCMD < \begin{figure}[!t]
%DIFDELCMD < 	\centering
%DIFDELCMD < 	\includegraphics[width=0.65\linewidth]{
%DIFDELCMD < 		src/figures/lambda_effect.drawio.pdf
%DIFDELCMD < 		}
%DIFDELCMD < 	%%%
%DIFDELCMD < \caption{%
{%DIFAUXCMD
\DIFdelFL{The effect of the Lagrange multiplier $\lambda_{j}$ on the adaptation direction of $\widehat{\mv{\theta}}_2$ is illustrated. 
        The notations $(\cdot^{\prime})$ and $(\cdot^{\prime\prime})$ denote two distinct cases corresponding to large and small values of the Lagrange multiplier, respectively, }%DIFDELCMD < \ie %%%
\DIFdelFL{$\lambda_{j}^{\prime} > \lambda_{j}^{\prime\prime}$.
        %DIF <  \vspace{-5mm}
    }}
	%DIFAUXCMD
%DIFDELCMD < \label{fig:lambda_effect}
%DIFDELCMD < \end{figure}
%DIFDELCMD < 

%DIFDELCMD < \hfill
%DIFDELCMD < 

%DIFDELCMD < %%%
\subsubsection{\DIFdel{Constraint Handling}} %DIFAUXCMD
\addtocounter{subsubsection}{-1}%DIFAUXCMD
%DIFDELCMD < \label{sec:sim:constraint}
%DIFDELCMD < 

%DIFDELCMD < %%%
\DIFdel{In this section, the }\DIFdelend %DIF >  여러 제약조건을 다루는 것에 대한 유연성
%DIF >  제약조건을 지속적으로 위반되더라도 최적점을 찾아가는 것
\DIFaddbegin 

\DIFadd{In this section, the }\DIFaddend constraint-handling \DIFdelbegin \DIFdel{capability of all controllers is investigated}\DIFdelend \DIFaddbegin \DIFadd{performance of the controllers is analyzed.
For a detailed analysis, we focus on the time interval $[\zoomStartTime, \qty{\zoomEndTime}{\second}]$ of Episode $2$ which is highlighted by the green dashed lines in Fig}\DIFaddend .\DIFaddbegin \DIFadd{~\ref{fig:ctrl:result:general}, where the difference in the considered input constraints among the controllers is most pronounced.
}\DIFaddend As shown in Fig.~\DIFdelbegin \DIFdel{\ref{fig:ctrl:result:tau:1}–\ref{fig:ctrl:result:tau:norm}, all controllers except (C$_3$), which did not consider the control input saturation , satisfied the control input saturation illustrated in Fig.~\ref{fig:robot:sat}.
In }\DIFdelend \DIFaddbegin \DIFadd{\ref{fig:ctrl:result:detail:q:top:view}, the tracking errors of all controllers were observed to be increased due to the input saturation effect, but }\DIFaddend (C$_1$) and (C$_2$) \DIFdelbegin \DIFdel{, control input saturation was addressed through corresponding input constraints within the constrained optimization framework, whereas (C}\DIFdelend \DIFaddbegin \DIFadd{exhibited smaller tracking errors compared to (C$_3$) and (C}\DIFaddend $_4$)\DIFdelbegin \DIFdel{employed an auxiliary system, as summarized in Table~\ref{table:controllers}.
}%DIFDELCMD < 

%DIFDELCMD < %%%
\DIFdel{To investigate the constraint handling process, the time interval from $29.4$ s to $31.9$ s was selected, during which the control input saturation was active for all controllers.
}\DIFdelend \DIFaddbegin \DIFadd{.
}\DIFaddend In Fig.~\DIFdelbegin \DIFdel{\ref{fig:ctrl:result:scope}, the control input locus, }\DIFdelend \DIFaddbegin \DIFadd{\ref{fig:ctrl:result:detail:u:top:view}, the commanded control inputs of (C$_1$) and (C$_2$) were observed to properly handle the input constraints, using the full authority of the actuator.
In contrast, }\DIFaddend the \DIFdelbegin \DIFdel{Lagrange multipliers, and the auxiliary states are illustrated over this time interval.
The properties of the controllers are illustrated in Fig.
~\ref{fig:ctrl:result:scope:control}, where the control input loci }\DIFdelend \DIFaddbegin \DIFadd{commanded control input }\DIFaddend of (C\DIFdelbegin \DIFdel{$_1$) and (C$_2$) fully explore the feasible input domain shown }\DIFdelend \DIFaddbegin \DIFadd{$_3$) was saturated at the box constraints, which is attributed to the design of the auxiliary system that only constrains the element-wise input overflow, rather than the actual input saturation, as illustrated }\DIFaddend in Fig.~\DIFdelbegin \DIFdel{\ref{fig:robot:sat}, whereas (C$_3$) exceeds the input saturation limits, and }\DIFdelend \DIFaddbegin \DIFadd{\ref{fig:exp:sat:cmp}.
Finally, the commanded control input of (C$_4$) was observed to exceed the input limits, which is expected since }\DIFaddend (C$_4$) \DIFdelbegin \DIFdel{is restricted in its first actuator input limit due to the auxiliary system design}\DIFdelend \DIFaddbegin \DIFadd{does not consider input constraints in its adaptation process}\DIFaddend .

%DIF <  beta
\DIFdelbegin \DIFdel{The input constraint handling from $29.4$ s to $31.9$ s is examined in detail. 
For }\DIFdelend \DIFaddbegin \begin{figure}[t]
\centering
  \caption{\DIFaddFL{Considered input saturation in (C$_1$), (C$_2$), and (C$_3$).}}
  \includegraphics[width=.89\linewidth]
  %DIF >  \includegraphics[width=\figSizeOneCol\linewidth]
  {
    % \DIFaddFL{
      figs/saturation_cmp.drawio.pdf
  % }
  }%DIF > 
  \label{fig:exp:sat:cmp}
\end{figure}

\DIFadd{The parametric study of the Lagrange multiplier update rates is conducted by comparing }\DIFaddend (C$_1$) and (C$_2$), \DIFdelbegin \DIFdel{activation of input constraints $c_j$ dynamically generated the corresponding Lagrange multipliers $\lambda_j$, as shown in Fig.~\ref{fig:ctrl:result:scope:beta}. 
These multipliers guided the adaptation of the weights $\estwth_i,\ \forall i\in\{0,1,2\}$, to reduce constraint violations in accordance with \eqref{eq:adap:th}.
}%DIFDELCMD < 

%DIFDELCMD < %%%
%DIF <  comparison; big, small beta
\DIFdel{To further illustrate this effect, }\DIFdelend \DIFaddbegin \DIFadd{focusing on the multiplier associated with the input norm constraint $\lambda_{\overline{\tau}}$.
In }\DIFaddend Fig.~\DIFdelbegin \DIFdel{\ref{fig:lambda_effect} compares the adaptation direction of the output layer's weight $\estwth_2$ under two different magnitudes of $\lambda_j$. 
Under Assumptions \ref{assum:convex} and \ref{assum:LICQ}, the convexity of $c_j$ in the $\estwth_2$-space (see Lemma~\ref{lem:convex:angle}) is illustrated in Fig.~\ref{fig:lambda_effect}.
The figure highlights the influence of the update rate $\beta_j$: a larger $\beta_j$ (as in (}\DIFdelend \DIFaddbegin \DIFadd{\ref{fig:ctrl:result:detail:multipliers}, $\lambda_{\overline{\tau}}$ of (}\DIFaddend C$_1$) \DIFdelbegin \DIFdel{) yields a larger $\lambda_j$, resulting in a more aggressive adjustment of $\estwth_2$ via the term $-\lambda_j\pptfrac{c_j}{\estwth_2}$ in \eqref{eq:adap:th}, thereby accelerating constraint satisfaction—though at the cost of oscillatory learning behavior. 
In contrast, the smaller update rate in }\DIFdelend \DIFaddbegin \DIFadd{and }\DIFaddend (C$_2$) \DIFdelbegin \DIFdel{leads to more conservative weight updates and smoother, but slower, convergence.
These behaviors are clearly observed in Fig.~\ref{fig:ctrl:result:scope:control} and Fig.
~\ref{fig:ctrl:result:scope:beta}, where }\DIFdelend \DIFaddbegin \DIFadd{was observed to be non-zero, indicating that the input norm constraint was active and enforced during this time interval.
Specifically, the Lagrange multiplier $\lambda_{\overline{\tau}}$ of }\DIFaddend (C$_1$)\DIFdelbegin \DIFdel{satisfies the input constraint $c_{\rbu}$ more quickly than (C$_2$), which shows repeated growth and reset of $\lambda_{\rbu}$. 
}%DIFDELCMD < 

%DIFDELCMD < %%%
%DIF <  zeta
\DIFdel{For (C$_4$) }\DIFdelend , \DIFdelbegin \DIFdel{the auxiliary state $\zeta_1$ was generated in }\DIFdelend \DIFaddbegin \DIFadd{which had a larger update rate, repeatedly grew and decayed, while that of (C$_2$) remained relatively constant.
Consequently, a larger update rate $\beta^u_{\overline{\tau}}$ provides a faster }\DIFaddend response to the \DIFdelbegin \DIFdel{control input saturation of the first link actuator $\tau_1$, as shown in Fig.
~\ref{fig:ctrl:result:scope:aux}.
Since the weights were adapted to reduce both the filtered error $\fe$ and the auxiliary state $\mv{\zeta}$, the control input saturation was subsequently reduced, as evident in Fig.~\ref{fig:ctrl:result:scope:control}.
Once the saturation was resolved at around $31.2$ s (see Fig.~\ref{fig:ctrl:result:tau:1}), the auxiliary state $\zeta_1$ converged to zero due to the stable state matrix $\mm{A}_\zeta$ in the auxiliary system dynamics, given by $\ddtt\mv{\zeta} = \mm{A}_\zeta \mv{\zeta} + \mm{B}_\zeta \Delta\rbu$ with $\Delta\rbu=\mv{0}_{2\times1}$, as shown again in }\DIFdelend \DIFaddbegin \DIFadd{input constraint but induces greater oscillations in the Lagrange multiplier, whereas a smaller update rate yields smoother but less responsive constraint enforcement.
This result confirms the trade-off discussed in Remark~\ref{rem:tuning}.
}

\DIFadd{Additionally, the nearly constant value of $\lambda_{\overline{\tau}}$ in }\DIFaddend Fig.~\DIFdelbegin \DIFdel{\ref{fig:ctrl:result:scope:aux}.
}%DIFDELCMD < 

%DIFDELCMD < \hfill 
%DIFDELCMD < 

%DIFDELCMD < %%%
\DIFdel{Finally, the handling of the weight norm constraints $c_{\theta_i},\ \forall i\in\{0,1,2\}$, is investigated. 
Among all controllers, the weight norm constraint was only activated for the output layer weights $\estwth_2$ in (C$_3$), as shown in }\DIFdelend \DIFaddbegin \DIFadd{\ref{fig:ctrl:result:detail:multipliers} and the small oscillations of $c_{\overline{\tau}}^u$ in }\DIFaddend Fig.~\DIFdelbegin \DIFdel{\ref{fig:ctrl:result:weight}. 
This is because the weights in (C$_1$), (C}\DIFdelend \DIFaddbegin \DIFadd{\ref{fig:ctrl:result:detail:u:top:view} of (C}\DIFaddend $_2$) \DIFdelbegin \DIFdel{, and (C$_4$) were implicitly suppressed by the input constraints, which were not applied in (C3).
%DIF <  , \ie by limiting the control input, the weight norms were also limited.
%DIF <  \MSRY{필요하다면 왜 제어입력 제약하는게 가중치 제약도 되는지 추가 가능}
Furthermore, the constraints on the inner and hidden layer weights were never violated; that is, both $\norm{\estwth_0}$ and $\norm{\estwth_1}$ consistently remained below their limits.
This was attributed to the fact that the backpropagated error was scaled by the gradients of the activation function, which are bounded in $(0,1]$, }%DIFDELCMD < \ie %%%
\DIFdel{$\sigma(\cdot)=\tanh(\cdot)$.
}%DIFDELCMD < 

%DIFDELCMD < %%%
\DIFdel{Additional details on the constraint handling process are provided in Fig.~\ref{fig:ctrl:result:scope:beta}. 
Similar to the case of the input constraints, the Lagrange multipliers $\lambda_{\theta_2}$ increased at around $30.4$ s, as shown in Fig.~\ref{fig:ctrl:result:scope:beta}, when $\estwth_2$ exceeded the prescribed norm constraint $\overline\theta_2$ (see Fig.
~\ref{fig:ctrl:result:weight}), thereby adjusting the adaptation direction of $\estwth_2$ to mitigate the violation of the constraint $c_{\theta_2}$}\DIFdelend \DIFaddbegin \DIFadd{suggest that the adaptation dynamics remain near an equilibrium characterized by a nonzero Lagrange multiplier.
This behavior is qualitatively consistent with the KKT-like optimality conditions discussed in Section~\ref{sec:main:result:sub:adap:laws}, according to which, at an equilibrium, the positive Lagrange multiplier $\lambda_{\overline{\tau}}$ is associated with active constraint satisfying $c_{\overline{\tau}}^u=0$ through the complementary slackness relation $\lambda_{\overline{\tau}}c_{\overline{\tau}}^u=0$.
Although the adaptation dynamics do not converge exactly to a stationary KKT equilibrium because of the practical limitations in the real-time implementation, the observed behavior can be interpreted as a near-equilibrium behavior}\DIFaddend .

%DIF <  In addition, it is worth noting that the input constraints were repeatedly violated even after the constraint handling mechanisms were applied, as shown in Fig.~\ref{fig:ctrl:result:scope:control}. 
%DIF <  This occurred because the controllers were implemented in discrete time with a relatively low sampling rate of $250$ Hz.
%DIF >  --------------------------------------------
%DIF >  val. result: real-time feasibility
%DIF >  --------------------------------------------

%DIF <  % ------------------------------------
    %DIF <  % Karush-Kuhn-Tucker (KKT) conditions}
    %DIF <  % ------------------------------------
    %DIF <  \hfill
\DIFaddbegin \subsubsection{\DIFadd{Feasibility of CONAC in Real-Time Implementation}}
\DIFaddend 

%DIF <  \subsubsection{Satisfaction of Karush-Kuhn-Tucker (KKT) conditions}
\DIFaddbegin \begin{figure}[t]
  \centering
  \includegraphics[width=0.99\linewidth]{\resultFig/Fig12.eps}
  \caption{
    \DIFaddFL{Computation time of (C$_1$) }\protect\colorLine{blue}{solid}\DIFaddFL{, (C$_2$) }\protect\colorLine{my_cyan}{solid}\DIFaddFL{, (C$_3$) }\protect\colorLine{magenta}{solid}\DIFaddFL{, and (C$_4$) }\protect\colorLine{gray}{solid} \DIFaddFL{on the OpenCR1.0 board.
  }}
  \label{fig:ctrl:result:cmp:time}
\end{figure}
\DIFaddend 

%DIF <  Since the proposed CONAC was implemented in discrete time, a direct demonstration of the KKT conditions at steady state is not straightforward. 
    %DIF <  % % \MSRY{
    %DIF <  % %     그림을 잘 확인해보니 제약조건이 위반되는 구간동안 정상상태에 도달하지 못하더라구요... 이전 논문 결과에서와 마찬가지로 그러할것으로 생각했는데 다른 요인이 있는 것 같습니다. 삭제하던지 아니면 보여질수없는 이유를 더 설명해야할 것으로 보입니다.
    %DIF <  % % }
    %DIF <  % % \ctxt{red}
    %DIF <  % % {
    %DIF <  % %     However, the near-constant values of the weights and Lagrange multipliers during constraint violation suggest an implicit satisfaction of the KKT conditions, as shown in Fig.~\ref{fig:ctrl:result:weight} and Fig.~\ref{fig:ctrl:result:scope:beta}
    %DIF <  % %     Specifically, the Lagrange multipliers corresponding to active constraints were positive, while those associated with inactive constraints remained zero.
    %DIF <  % % }
    %DIF <  % % \ctxt{blue}
    %DIF <  % % {
    %DIF <      This is because, the KKT conditions are typically defined for continuous-time systems, and their direct application to discrete-time systems requires careful consideration of the sampling rate.
    %DIF <      Although exact satisfaction of the KKT conditions was not observed—due to the lack of a clear steady state in the discrete-time implementation—the control results remain sufficiently close to the optimal behavior implying that the KKT conditions were approximately satisfied.
    %DIF <  % % }
\DIFdelbegin %DIFDELCMD < 

%DIFDELCMD < %%%
%DIF <  ------------------------------------
%DIF <  Computational Time
%DIF <  ------------------------------------
%DIFDELCMD < \hfill
%DIFDELCMD < 

%DIFDELCMD < %%%
\subsubsection{\DIFdel{Effectiveness of CONAC in Real-time Applications}}
%DIFAUXCMD
\addtocounter{subsubsection}{-1}%DIFAUXCMD
%DIFDELCMD < 

%DIFDELCMD < %%%
\DIFdel{The effectiveness of the proposed CONAC in }\DIFdelend \DIFaddbegin \DIFadd{The feasibility of the }\DIFaddend real-time \DIFdelbegin \DIFdel{applications was evaluated by measuring the computational timeusing CPU clock ticks on the OpenCR1.0 board}\DIFdelend \DIFaddbegin \DIFadd{implementation is analyzed in terms of the computation time}\DIFaddend .
As shown in Fig.~\DIFdelbegin \DIFdel{\ref{fig:ctrl:real:result:cmp:time}, the computational times }\DIFdelend \DIFaddbegin \DIFadd{\ref{fig:ctrl:result:cmp:time}, the computation time }\DIFaddend of all controllers \DIFdelbegin \DIFdel{remained below $4$ ms, demonstrating that the proposed CONAC is suitable for }\DIFdelend \DIFaddbegin \DIFadd{is consistently below the sampling time of $\qty{500}{\micro\second}$, confirming the feasibility of }\DIFaddend real-time implementation at a sampling \DIFdelbegin \DIFdel{rate of $250$ Hz.
}\DIFdelend \DIFaddbegin \DIFadd{frequency of $\qty{500}{\hertz}$ ($\qty{2}{\milli\second}$ maximum sampling time).
This is attributed to the fact that the proposed CONAC does not require solving an optimization problem at each time step, but rather continuously adapts the adaptive variables to solve the optimization problem in an online manner.
}\DIFaddend 

\DIFdelbegin %DIFDELCMD < \begin{figure}[t]
%DIFDELCMD <     \centering
%DIFDELCMD <         \includegraphics[width=\figSizeOneCol\linewidth]
%DIFDELCMD <         {
%DIFDELCMD <             %%%
\DIFdelFL{src/measurement/figures/compare/Fig10.
eps
        }%DIFDELCMD < }%%%
%DIF < 
    %DIFDELCMD < \caption{%
{%DIFAUXCMD
\DIFdelFL{Computational time $T_{\rm{comp}}$ of (C$_1$) }%DIFDELCMD < \protect\colorLine{blue}{solid}%%%
\DIFdelFL{, (C$_2$) }%DIFDELCMD < \protect\colorLine{cyan}{solid}%%%
\DIFdelFL{, (C$_3$) }%DIFDELCMD < \protect\colorLine{gray}{solid}%%%
\DIFdelFL{, and (C$_4$) }%DIFDELCMD < \protect\colorLine{magenta}{solid}%%%
\DIFdelFL{.
        %DIF <  \vspace{-5mm}
    }}
    %DIFAUXCMD
%DIFDELCMD < \label{fig:ctrl:real:result:cmp:time}
%DIFDELCMD < \end{figure}
%DIFDELCMD < %%%
\DIFdelend \DIFaddbegin \subsection{\DIFadd{Reproducibility and Video Demonstration}}
\DIFaddend 

%DIF <   SECTION CONCLUSION ======================================
\DIFaddbegin \DIFadd{For reproducibility, the source code and the experimental data are available in \mbox{%DIFAUXCMD
\cite{Ryu2026Software}}\hskip0pt%DIFAUXCMD
.
A supplementary video demonstrating the experimental results is provided with this paper.
}

%DIF >  ------------------------------------
%DIF >        SECTION CONCLUSION
%DIF >  ------------------------------------
\DIFaddend \section{Conclusion} \label{sec:conclusion}

This paper proposed \DIFdelbegin \DIFdel{a constrained optimization-based neuro-adaptive controller (CONAC ) for unknown Euler-Lagrange }\DIFdelend \DIFaddbegin \DIFadd{CONAC for unknown MIMO }\DIFaddend systems, addressing \DIFdelbegin \DIFdel{both weight norm and }\DIFdelend \DIFaddbegin \DIFadd{multiple }\DIFaddend convex input constraints within a unified \DIFaddbegin \DIFadd{constrained }\DIFaddend optimization framework. 
The \DIFaddbegin \DIFadd{adaptation laws were derived from the constrained optimization problem, which was formulated to minimize the tracking error while satisfying the multiple convex input constraints.
The }\DIFaddend stability of CONAC was analyzed using Lyapunov theory, demonstrating \DIFdelbegin \DIFdel{bounded }\DIFdelend \DIFaddbegin \DIFadd{the boundedness of the }\DIFaddend tracking errors and \DIFdelbegin \DIFdel{weight estimation under real-time adaptation.
}%DIFDELCMD < 

%DIFDELCMD < %%%
\DIFdel{CONAC effectively incorporated both convex input constraints and weight norm constraints, ensuring that actuator limitations and neural network weights remained within predefined bounds.
By embedding these constraints into the optimization process}\DIFdelend \DIFaddbegin \DIFadd{adaptive variables, including the DNN weights and Lagrange multipliers.
By deriving adaptation laws from the constrained optimization problem}\DIFaddend , CONAC ensured weight convergence in accordance with the \DIFdelbegin \DIFdel{Karush-Kuhn-Tucker (KKT ) }\DIFdelend \DIFaddbegin \DIFadd{KKT optimality }\DIFaddend conditions, thereby ensuring both stability and optimality.

Real-time implementation \DIFdelbegin \DIFdel{validated the superior performance of CONAC compared to a conventional method.
CONAC successfully handled complex input constraintswhile rigorously managing weight norms, resulting in improved tracking accuracy and stable behavior without significant oscillations}\DIFdelend \DIFaddbegin \DIFadd{demonstrated the capability of CONAC in handling multiple convex input constraints simultaneously, while achieving superior tracking performance compared to other NAC methods.
Moreover, the parametric study revealed the effect of the update rates of the Lagrange multipliers on the enforcement of input constraints.
In addition, the experiments validated the feasibility of CONAC in real-time implementation}\DIFaddend .

Future work may extend this approach to include constraints \DIFdelbegin \DIFdel{on both control inputs and }\DIFdelend \DIFaddbegin \DIFadd{ons }\DIFaddend system states, \DIFaddbegin \DIFadd{and minimize the infinite horizon error cost, thereby }\DIFaddend further enhancing the \DIFdelbegin \DIFdel{flexibility and robustness of neuro-adaptive control }\DIFdelend \DIFaddbegin \DIFadd{capability and optimality of NAC }\DIFaddend systems within constrained optimization frameworks.

%DIF < % The Appendices part is started with the command \appendix;
%DIF < % appendix sections are then done as normal sections
\DIFdelbegin %DIFDELCMD < \appendix
%DIFDELCMD < %%%
\DIFdelend %DIF >  -------------------------------------
%DIF >  Appendix
%DIF >  -------------------------------------
%DIF >  \appendix

\DIFdelbegin \subsection{\DIFdel{Input Constraint Candidates}}%DIFAUXCMD
\addtocounter{subsection}{-1}%DIFAUXCMD
%DIFDELCMD < \label{sec:appen:cstr} 
%DIFDELCMD < 

%DIFDELCMD < %%%
\DIFdel{This section introduces potential input constraints that can be incorporated into the proposed CONAC framework.
}%DIFDELCMD < 

%DIFDELCMD < %%%
\subsubsection{\DIFdel{Input Bound Constraint}}%DIFAUXCMD
\addtocounter{subsubsection}{-1}%DIFAUXCMD
%DIFDELCMD < \label{sec:appen:cstr:input:bound}
%DIFDELCMD < 

%DIFDELCMD < %%%
\DIFdel{Most physical systems are subject to control input limitations due to inherent electrical and mechanical constraints.
These are expressed as $\mv{c}_{\overline \tau}:= [c_{\overline \tau_i}]_{i\in\{1,\cdots,n\}}$ and $\mv{c}_{\underline\tau}:= [c_{\underline\tau_i}]_{i\in\{1,\cdots,n\}}$, where
}%DIFDELCMD < \begin{equation}
%DIFDELCMD <     \begin{aligned}
%DIFDELCMD <         c_{\overline \tau_i}=\tau_{(i)} - {\tau_{\overline \tau_i}}
%DIFDELCMD <         ,
%DIFDELCMD <         \quad
%DIFDELCMD <         c_{\underline\tau_i}={\tau_{\underline\tau_i}}-\tau_{(i)}
%DIFDELCMD <         ,
%DIFDELCMD <     \end{aligned}
%DIFDELCMD <     \label{eq:cstr:input:bound}
%DIFDELCMD < \end{equation}
%DIFDELCMD < %%%
\DIFdel{with $\tau_{\overline \tau_i}\in\R$ and $\tau_{\underline\tau_i}\in\R$ representing the maximum and minimum control input bounds, respectively.
The gradients of $\mv{c}_{\overline \tau}$ and $\mv{c}_{\underline\tau}$ with respect to $\estwth$ are given by
}%DIFDELCMD < \begin{equation}
%DIFDELCMD <     \begin{aligned}
%DIFDELCMD <         \pptfrac{\mv c_{\overline \tau}}{\estwth}
%DIFDELCMD <         =& 
%DIFDELCMD <         % \left[
%DIFDELCMD <         %     \pptfrac{c_{\overline \tau_i}}{\estwth}^\top
%DIFDELCMD <         % \right]_{i\in\{1,\cdots,n\}}
%DIFDELCMD <         % \begin{bmatrix}
%DIFDELCMD <         %     \pptfrac{c_{\overline \tau_1}}{\estwth}^\top \\
%DIFDELCMD <         %     \vdots \\
%DIFDELCMD <         %     \pptfrac{c_{\overline \tau_n}}{\estwth}^\top
%DIFDELCMD <         % \end{bmatrix}
%DIFDELCMD <         = 
%DIFDELCMD <             +\pptfrac{\estNN}{\estwth}
%DIFDELCMD <             % \\
%DIFDELCMD <         % =&
%DIFDELCMD <         =
%DIFDELCMD <         +
%DIFDELCMD <         \begin{bmatrix}
%DIFDELCMD <             (\mm{I}_{l_{k+1}}\otimes \estact_{k}^\top)&
%DIFDELCMD <             \!\cdots\! &
%DIFDELCMD <             (
%DIFDELCMD <                 \cdot
%DIFDELCMD <             )
%DIFDELCMD <         \end{bmatrix} 
%DIFDELCMD <         \in
%DIFDELCMD <         \mathbb R^{n\times \Xi}
%DIFDELCMD <         , 
%DIFDELCMD <         \\
%DIFDELCMD <         \pptfrac{\mv c_{\underline\tau}}{\estwth}         
%DIFDELCMD <         =
%DIFDELCMD <         & 
%DIFDELCMD <         % \left[
%DIFDELCMD <         %     \pptfrac{c_{\underline\tau_i}}{\estwth}^\top
%DIFDELCMD <         % \right]_{i\in\{1,\cdots,n\}}
%DIFDELCMD <         % \begin{bmatrix}
%DIFDELCMD <         %     \pptfrac{c_{\underline\tau_1}}{\estwth}^\top \\
%DIFDELCMD <         %     \vdots \\
%DIFDELCMD <         %     \pptfrac{c_{\underline\tau_n}}{\estwth}^\top
%DIFDELCMD <         % \end{bmatrix}
%DIFDELCMD <         = 
%DIFDELCMD <         -\pptfrac{\estNN}{\estwth}
%DIFDELCMD <         % \\
%DIFDELCMD <         % =&
%DIFDELCMD <         =
%DIFDELCMD <         -\begin{bmatrix}
%DIFDELCMD <             (\mm{I}_{l_{k+1}}\otimes \estact_{k}^\top)&
%DIFDELCMD <             % V_k^\top \phi_{k}' (I_{l_{k}}\otimes  \phi_{k-1}^\top)&
%DIFDELCMD <             \!\cdots\! &
%DIFDELCMD <             (
%DIFDELCMD <                 \cdot
%DIFDELCMD <             )
%DIFDELCMD <         \end{bmatrix} 
%DIFDELCMD <         \in
%DIFDELCMD <         \mathbb R^{n\times \Xi}
%DIFDELCMD <         .
%DIFDELCMD <     \end{aligned}
%DIFDELCMD <     \label{eq:cstr:input:bound:grad}
%DIFDELCMD < \end{equation}
%DIFDELCMD < 

%DIFDELCMD < %%%
\subsubsection{\DIFdel{Input Norm Constraint}}%DIFAUXCMD
\addtocounter{subsubsection}{-1}%DIFAUXCMD
%DIFDELCMD < \label{sec:appen:cstr:input:norm}
%DIFDELCMD < 

%DIFDELCMD < %%%
\DIFdel{Consider the control input $\rbu$ as the torque corresponding to its generalized coordinate. 
Since torque is typically linearly proportional to current, actuators that share a common power source are often subject to total current limitations. 
This can be captured by the following inequality constraint: 
}%DIFDELCMD < \begin{equation}
%DIFDELCMD <     c_{\rbu}
%DIFDELCMD <     =
%DIFDELCMD <     \tfrac{1}{2}\
%DIFDELCMD <     \left(
%DIFDELCMD <         \norm{
%DIFDELCMD <         \mv{\tau}
%DIFDELCMD <         }^2 
%DIFDELCMD <         -
%DIFDELCMD <         \overline\tau^2
%DIFDELCMD <     \right)
%DIFDELCMD <     ,
%DIFDELCMD <     \label{eq:cstr:input:norm}
%DIFDELCMD < \end{equation}
%DIFDELCMD < %%%
\DIFdel{with the maximum allowable control input magnitude $\overline\tau\in\R_{>0}$. This input norm constraint is also commonly applied in current and torque control problems for electric motors \mbox{%DIFAUXCMD
\cite{Choi:2024aa}}\hskip0pt%DIFAUXCMD
.
%DIF <  The corresponding Lagrangian multiplier is $\lambda_{u_b}$. 
The gradient of $c_{\rbu}$ with respect to $\estwth$ is given by
}%DIFDELCMD < \begin{equation}
%DIFDELCMD <     \pptfrac{\mv c_{\rbu}}{\estwth}
%DIFDELCMD <     \!=\! 
%DIFDELCMD <     \textstyle\sum_{i=1}^n \tau_{i} 
%DIFDELCMD <     \left(
%DIFDELCMD <         \myrow_i
%DIFDELCMD <         \left(
%DIFDELCMD <             \pptfrac{\estNN}{\estwth}
%DIFDELCMD <         \right)
%DIFDELCMD <     \right)^\top  
%DIFDELCMD <     \!=\! 
%DIFDELCMD <     {\pptfrac{\estNN}{\estwth}}^\top
%DIFDELCMD <     \mv{\tau}
%DIFDELCMD <     % \tau^\top (I_{l_{k+1}}\otimes \estNN_k^\top)
%DIFDELCMD <     \in \mathbb R^{\Xi}.
%DIFDELCMD <     \label{eq:cstr:input:norm:grad}
%DIFDELCMD < \end{equation}
%DIFDELCMD < 

%DIFDELCMD < \hfill
%DIFDELCMD < 

%DIFDELCMD < %%%
\DIFdel{It should be noted that constraints \eqref{eq:cstr:input:bound} and \eqref{eq:cstr:input:norm} can be imposed simultaneously, as their gradients \eqref{eq:cstr:input:bound:grad} and \eqref{eq:cstr:input:norm:grad} are linearly independent, satisfying Assumption \ref{assum:LICQ}, }%DIFDELCMD < \ie %%%
\DIFdel{the LICQ condition.
}%DIFDELCMD < 

%DIFDELCMD < %%%
\DIFdelend % ------------------------------------
% Bibliography
% ------------------------------------
\bibliographystyle{IEEEtran}
\DIFdelbegin %DIFDELCMD < \bibliography{\template/refs}
%DIFDELCMD < %%%
%DIF <  \bibliography{\template/refs_abbr}
\DIFdelend %DIF >  \bibliographystyle{apalike}
\DIFaddbegin \bibliography{localRefs, fromLR}
\DIFaddend 

% ------------------------------------
% Biography
% ------------------------------------
% \vspace{-10mm}
\DIFdelbegin %DIFDELCMD < \begin{IEEEbiography}[{
%DIFDELCMD <     \includegraphics[width=1in,height=1.25in,clip,keepaspectratio]{\template/people/msRyu.jpg}}]{Myeongseok Ryu}
%DIFDELCMD <     %%%
\DIFdel{received the B.S. degree in Mechanical Engineering from Incheon National University, South Korea, in 2023, and the M.S. degree in Mechanical Engineering from GIST, South Korea, in 2025. 
    He is currently pursuing the Ph.D. degree in the Cho Chun Shik Graduate School of Mobility at Korea Advanced Institute of Science and Technology (KAIST), South Korea.
    His research interests include neuro-adaptive control, online optimization, and contraction theory.
    %DIF <  \vspace{-10mm}
}%DIFDELCMD < \end{IEEEbiography}
%DIFDELCMD < %%%
\DIFdelend %DIF >  \begin{IEEEbiography}[{
%DIF >      \includegraphics[width=1in,height=1.25in,clip,keepaspectratio]{\template/people/msRyu.jpg}}]{Myeongseok Ryu}
%DIF >      received the B.S. degree in Mechanical Engineering from Incheon National University, South Korea, in 2023, and the M.S. degree in Mechanical Engineering from GIST, South Korea, in 2025. 
%DIF >      He is currently pursuing the Ph.D. degree in the Cho Chun Shik Graduate School of Mobility at Korea Advanced Institute of Science and Technology (KAIST), South Korea.
%DIF >      His research interests include neuro-adaptive control, online optimization, and contraction theory.
%DIF >      % \vspace{-10mm}
%DIF >  \end{IEEEbiography}

\DIFdelbegin %DIFDELCMD < \begin{IEEEbiography}[{
%DIFDELCMD <     \includegraphics[width=1in,height=1.25in,clip,keepaspectratio]{\template/people/dhHong.jpg}}]{Donghwa Hong}
%DIFDELCMD <     %%%
\DIFdel{received his B.S. degree in Physics and Engineering Physics from Yonsei University, Wonju, South Korea, in 2023. He is currently pursuing the M.S. degree in Mechanical Engineering from GIST, Gwangju, South Korea.  
His research interests include robot control, neural-networks, and model identification. 
    %DIF <  \vspace{-10mm}
}%DIFDELCMD < \end{IEEEbiography}
%DIFDELCMD < %%%
\DIFdelend %DIF >  \begin{IEEEbiography}[{

\DIFdelbegin %DIFDELCMD < \begin{IEEEbiography}[{
%DIFDELCMD <     \includegraphics[width=1in,height=1.25in,clip,keepaspectratio]{\template/people/khChoi.jpg}}]{Kyunghwan Choi}
%DIFDELCMD <     %%%
\DIFdel{received his B.S., M.S., and Ph.D. degrees in Mechanical Engineering from KAIST, Daejeon, South Korea, in 2014, 2016, and 2020, respectively. Following his Ph.D., he worked at the Center for Eco-Friendly \& Smart Vehicles at KAIST as a Postdoctoral Fellow and later as a Research Assistant Professor. In 2022, he joined GIST as an Assistant Professor in the Department of Mechanical and Robotics Engineering. Currently, he serves as an Assistant Professor at the Cho Chun Shik Graduate School of Mobility at KAIST, where he is also the Director of the Mobility Intelligence and Control Laboratory. His research focuses on optimal and learning-based control for connected, automated, and electrified vehicles (CAEVs).
}%DIFDELCMD < \end{IEEEbiography}
%DIFDELCMD < %%%
\DIFdelend %DIF >  \includegraphics[width=1in,height=1.25in,clip,keepaspectratio]{\template/people/dhHong.jpg}}]{Donghwa Hong}
%DIF >      received his B.S. degree in Physics and Engineering Physics from Yonsei University, Wonju, South Korea, in 2023, and the M.S. degree in Mechanical and Robotics Engineering from GIST, Gwangju, South Korea, in 2026. Following his M.S., he is working in the Cho Chun Shik Graduate School of Mobility at KAIST as a researcher. His research interests include online learning and model identification. 
%DIF >      % \vspace{-10mm}
%DIF >  \end{IEEEbiography}
\DIFaddbegin 

%DIF >  \begin{IEEEbiography}[{
%DIF >      \includegraphics[width=1in,height=1.25in,clip,keepaspectratio]{\template/people/khChoi.jpg}}]{Kyunghwan Choi}
%DIF >      received his B.S., M.S., and Ph.D. degrees in Mechanical Engineering from KAIST, Daejeon, South Korea, in 2014, 2016, and 2020, respectively. Following his Ph.D., he worked at the Center for Eco-Friendly \& Smart Vehicles at KAIST as a Postdoctoral Fellow and later as a Research Assistant Professor. In 2022, he joined GIST as an Assistant Professor in the Department of Mechanical and Robotics Engineering. Currently, he serves as an Assistant Professor at the Cho Chun Shik Graduate School of Mobility at KAIST, where he is also the Director of the Mobility Intelligence and Control Laboratory. His research focuses on optimal and learning-based control for connected, automated, and electrified vehicles (CAEVs).
%DIF >  \end{IEEEbiography}
\DIFaddend 

\end{document}
